\documentclass[11pt]{article}
\usepackage[utf8]{inputenc}
\usepackage{amsmath,amssymb,amsthm}
\usepackage{hyperref}
\usepackage[capitalise,noabbrev]{cleveref}
\usepackage{geometry}
\usepackage{booktabs}
\usepackage{graphicx}
\usepackage{array}
\usepackage{xcolor}
\geometry{margin=1in}

%----------------------------------------------------------------------
% Theorem environments
%----------------------------------------------------------------------
\newtheorem{theorem}{Theorem}[section]
\newtheorem{proposition}[theorem]{Proposition}
\newtheorem{corollary}[theorem]{Corollary}
\newtheorem{lemma}[theorem]{Lemma}
\theoremstyle{definition}
\newtheorem{definition}[theorem]{Definition}
\newtheorem{conjecture}[theorem]{Conjecture}
\theoremstyle{remark}
\newtheorem{remark}[theorem]{Remark}

%----------------------------------------------------------------------
% Notation macros
%----------------------------------------------------------------------
% Fisher information matrix
\newcommand{\Fish}{F}
% Sign matrix
\newcommand{\sgn}{S}
% Signed metric M = F^{1/2} S F^{1/2}
\newcommand{\Met}{M}
% Observer vector (timelike eigenvector)
\newcommand{\obs}{u}
% Vorticity tensor
\newcommand{\vort}{\omega}
% Coupling vector
\newcommand{\coupl}{\mathbf{J}}
% Partition function
\newcommand{\partf}{Z}
% Christoffel symbol shorthand
\newcommand{\Chr}[2]{\Gamma^{#1}_{#2}}
% Covariance operator
\newcommand{\Cov}{\operatorname{Cov}}
% Variance operator
\newcommand{\Var}{\operatorname{Var}}
% Diagonal matrix
\newcommand{\diag}{\operatorname{diag}}
% sech
\newcommand{\sech}{\operatorname{sech}}
% Edge observable
\newcommand{\edgeo}{\sigma}
% Graph cycle rank (first Betti number)
\newcommand{\betti}{\beta_1}
% Einstein tensor
\newcommand{\Ein}{G}
% Ricci tensor
\newcommand{\Ric}{R}
% Riemann tensor (full tensor shorthand, same letter R)
\newcommand{\Riem}{R}
% Scalar curvature
\newcommand{\Scal}{\mathcal{R}}
% Spatial projector
\newcommand{\proj}{h}
% Expansion scalar
\newcommand{\expan}{\theta}
% Shear tensor
\newcommand{\shear}{\varsigma} % Use ς (varsigma) to avoid collision with σ_e (edge observable)

%----------------------------------------------------------------------
% Title and authors
%----------------------------------------------------------------------
\title{The Flatness of Statistical Manifolds for Acyclic Graphs:\\
Characterizing the Statistical Vacuum Einstein Condition}

\author{Max Zhuravlev%
\thanks{Independent researcher. Email: \texttt{max@vibecodium.ai}}}

\date{February 2026}

%----------------------------------------------------------------------
\begin{document}
%----------------------------------------------------------------------

\maketitle

\noindent\textbf{arXiv category:} cond-mat.stat-mech (cross-list: gr-qc, math-ph)

%----------------------------------------------------------------------
\begin{abstract}
%----------------------------------------------------------------------
% Condensed abstract (~190 words). Full result list in Section 1.4.
%----------------------------------------------------------------------
We study the curvature of the \emph{signed Fisher metric}
$\Met = \Fish^{1/2} \sgn \Fish^{1/2}$ on the statistical manifold of
Ising and $O(N)$ models on graphs, where $\Fish$ is the Fisher
information matrix of edge observables and $\sgn$ selects Lorentzian
signature via spectral gap maximization.
We establish a \emph{statistical vacuum Einstein equation}: the Einstein
tensor $\Ein_{ab} = 0$ for all couplings and sign assignments if the graph
is a tree (proved in full generality, at $h = 0$). The converse---that
non-trees admit $\Ein_{ab} \neq 0$ at some couplings---is posed as a
conjecture and verified on 839 connected graphs with $n \leq 7$ vertices
(zero counterexamples), with analytic results for small cycles.
For $C_4$ we derive the first explicit scalar curvature formula for a
signed Fisher metric (to our knowledge).
A universal identity
$\Ric^a{}_{bcd}|_{\mathrm{lin}} = -2\,\Ric^a{}_{bcd}|_{\mathrm{quad}}$
is verified numerically for all single-block exponential families tested
(49 cases, zero violations), with analytic proof for cycle graphs $C_n$.
The observer vorticity vanishes whenever every biconnected component has
cycle rank $\leq 1$ (proved); the converse is verified on 839 graphs.
On square grids with open boundary conditions, $|\Scal(J_c)| \sim n^{d_{\Scal}}$ with
$d_{\Scal} \approx 2.83$ at the 2D Ising critical point (equivalently
$|\Scal| \sim m^{2.15}$ where $m = |E|$ is the number of edges).
For periodic boundary conditions (tori), a different exponent governs the
scaling: $|\Scal(J_c)| \sim m^{d_R^{\mathrm{PBC}}}$ with the conjectured
formula $d_R^{\mathrm{PBC}} = (d\nu + 2\eta)/(d\nu + \eta)$.
For 2D Ising this predicts $d_R^{\mathrm{PBC}} = 10/9 \approx 1.111$, confirmed
by exact transfer-matrix computation on tori $L = 3, \ldots, 8$
(measured effective exponent $1.1110$ at $L = 7{\to}8$).
Continuous $O(N \geq 2)$ spins eliminate the non-Lorentzian gap:
$100\%$ Lorentzian selection on 719 tested graphs.
Full results are listed in \cref{subsec:main-results}.

\medskip
\noindent\textbf{Keywords:} Fisher information geometry, Lorentzian
signature, Ising model, Einstein vacuum equation, signed metric,
biconnected components, curvature scaling, periodic boundary conditions,
torus, vorticity.
\end{abstract}

%======================================================================
\tableofcontents
%======================================================================

%======================================================================
\section{Introduction}
\label{sec:intro}
%======================================================================

\subsection{The Problem: Curvature of a Statistical Metric}

The Fisher information matrix $\Fish_{ab} = \Cov_{\coupl}(\edgeo_a, \edgeo_b)$
is the canonical Riemannian metric on the statistical manifold of an Ising
observer, where $\edgeo_e = \sigma_i \sigma_j$ denotes the edge observable
for edge $e = (i,j)$. Matsueda~\cite{Matsueda2013} computed the Christoffel
symbols and geodesics for the unsigned Fisher metric on small Ising
systems, finding nontrivial geometric structure even in low-dimensional
examples. The natural question left open was: \emph{what is the curvature
of the signed variant} $\Met = \Fish^{1/2} \sgn \Fish^{1/2}$ that carries
Lorentzian signature?
Earlier finite-size information-geometric curvature calculations
focused on low-dimensional thermodynamic or uniform-coupling parameter spaces
(e.g., Brody--Ritz~\cite{BrodyRitz2003}, Erdmenger et al.~\cite{Erdmenger2020}),
rather than the full per-edge coupling manifold considered here.

This question is not merely technical. The signed metric arises as the
unique mechanism for producing Lorentzian signature from positive-definite
Fisher data. We showed in a companion paper~\cite{Paper5} that the spectral
gap weighting functional
\begin{equation}
  W(q) = \beta_c(q) \cdot L_{\mathrm{gap}}(q)
  \label{eq:W-def}
\end{equation}
uniquely selects $q = 1$ (one negative sign) for sparse Ising observers,
so the relevant metric is
\begin{equation}
  \Met_{ab} = \bigl(\Fish^{1/2} \sgn \Fish^{1/2}\bigr)_{ab},
  \qquad \sgn = \diag(+1, \ldots, +1, -1),
  \label{eq:signed-metric}
\end{equation}
with signature $(m{-}1, 1)$ --- one timelike direction and $m{-}1$
spacelike directions, where $m = |E|$ is the number of edges.
This observer-centric geometric program is aligned with the emergent-gravity
learning-systems perspective developed by Vanchurin~\cite{Vanchurin2022}.

The curvature of $\Met$ encodes how the statistical manifold deviates
from flat space. We find that this deviation is \emph{exactly} controlled
by the presence of loops (cycles) in the underlying graph: trees give flat
geometry, cycles give curved geometry. This is a precise, computable, and
physically interpretable relationship.

\subsection{Context: Information Geometry and Ising Models}

The differential geometry of statistical manifolds has a rich history
originating in the work of Rao~\cite{Rao1945} and formalized by
Amari~\cite{Amari1998} and Cencov~\cite{Cencov1982}. The Fisher information
matrix plays the role of the metric tensor, and the Levi-Civita connection
of this metric (together with a family of dual connections) defines the
information-geometric structure.

For exponential family distributions --- of which the Ising model is the
canonical discrete example --- the Fisher matrix has a particularly clean
form: $\Fish_{ab} = \partial^2 A / \partial J_a \partial J_b$ where
$A = \log Z$ is the log-partition function. The geometry of the Ising
statistical manifold has been studied computationally
(Matsueda~\cite{Matsueda2013}, Machta--Sethna~\cite{Machta2013}) but without
focusing on the signed variant that arises from the Lorentzian signature
requirement.

The key structural observation that enables the present work is the
\emph{Tree Fisher Identity} (proven in \cite{Paper1}):
\begin{equation}
  \Fish = \diag\bigl(\sech^2(J_1), \ldots, \sech^2(J_m)\bigr)
  \quad \text{on any tree.}
  \label{eq:tree-fisher-intro}
\end{equation}
This diagonal form means that edge observables are statistically
independent on trees: there are no correlations between different edges.
The signed metric $\Met$ on a tree is therefore a product of one-dimensional
hyperbolic spaces, and its curvature vanishes exactly.

\subsection{The Signed Metric as Vielbein Construction}

A key conceptual contribution of this paper (developed fully in
\cref{sec:vielbein}) is the identification of $\Met = \Fish^{1/2} \sgn \Fish^{1/2}$
as a \emph{vielbein construction}. Writing $e = \Fish^{1/2}$ (the symmetric
positive-definite square root of $\Fish$) and $\eta = \sgn$ (playing the
role of the flat Minkowski metric), we have
\begin{equation}
  \Met = e^T \eta\, e,
  \label{eq:vielbein-intro}
\end{equation}
which is the standard GR formula $g_{\mu\nu} = e^a{}_\mu \eta_{ab} e^b{}_\nu$
with $e^a{}_\mu = (\Fish^{1/2})^a{}_\mu$. This identification reveals that
$\Fish^{1/2}$ is the \emph{frame field} (vielbein) encoding how the
statistical manifold differs from flat Minkowski space.

Unlike the ansatz $M = FSF$ of the signed-edge construction
\cite{Paper1}, the vielbein form $\Met = e^T \eta\, e$ is uniquely
determined by three independently motivated inputs:
\begin{enumerate}
\item Cencov's uniqueness theorem~\cite{Cencov1982} selects $\Fish$ as the
  unique Riemannian metric on statistical manifolds equivariant under
  sufficient statistics.
\item Spectral gap selection~\cite{Paper5} determines $\sgn$ (one negative
  entry is optimal).
\item The canonical symmetric positive-definite square root gives $e = \Fish^{1/2}$
  as the unique symmetric vielbein.
\end{enumerate}
These three arguments single out the construction uniquely within the
class of natural vielbein-based metrics, providing stronger motivation
for the ansatz.

\subsection{Main Results}
\label{subsec:main-results}

We establish eight main theorems, together with
strong empirical evidence for several conjectures.

\medskip
\noindent\textbf{Theorem~\ref{thm:tree-flat} (Tree Flatness).}
For the Ising model on any tree $G = (V, E)$ with zero magnetic field
and any sign assignment $\sgn$, the signed metric $\Met$ is flat:
all Riemann tensor components vanish, $\Ric_{abcd} = 0$.
The flat coordinates are $x_a = \arctan(\sinh(J_a))$, and $\Met$ is
isometric to flat Minkowski space (in appropriate signature).

\medskip
\noindent\textbf{Proposition~\ref{prop:sign-indep} (Sign Independence).}
On trees, the Christoffel symbols $\Chr{a}{bc}$ of $\Met$ are
independent of the sign assignment $\sgn$. The connection is the same
for Lorentzian and Riemannian signatures. The nonzero components are
$\Chr{a}{aa} = -\tanh(J_a)$ (no summation), with all off-diagonal
components vanishing.

\medskip
\noindent\textbf{Theorem~\ref{thm:c4-exact} ($C_4$ Reduced Curvature).}
For the 4-cycle $C_4$ on the symmetry-reduced 2D submanifold with uniform
coupling $J > 0$:
\begin{equation}
  \Scal^{(2)}(J) = \frac{4\sinh^2(4J)}{(3\cosh(4J)+1)^2}.
  \label{eq:c4-scalar-intro}
\end{equation}
This is positive and satisfies the \emph{factor-of-two law} (numerically verified):
$\Scal^{(2)}_{\Met} = \Scal_{\Fish}/2$. In the full 4D space, the sign changes
(\cref{rem:c4-reduction}).

\medskip
\noindent\textbf{Theorem~\ref{thm:vacuum} (Statistical Vacuum Einstein Equation).}
If $G$ is a tree, then $\Ein_{ab} = 0$ for all couplings (with $h = 0$)
and all sign assignments $\sgn$. The converse direction is posed as
Conjecture~\ref{conj:vacuum-converse} and is verified on $839$ connected
graphs with $n \leq 7$ vertices (zero counterexamples).

\medskip
\noindent\textbf{Theorem~\ref{thm:vorticity} (Vorticity Classification).}
For the unique timelike observer vector $\obs^a$ of $\Met$ (when $q = 1$),
the vorticity $\vort_{ab} = 0$ whenever the graph has cycle rank
$\betti(G) \leq 1$.

\medskip
\noindent\textbf{Conjecture~\ref{cor:bcc-chain-complete} (BCC-Chain, Full Equivalence).}
$\vort_{ab} = 0$ for all $q = 1$ sign assignments and all $J > 0$
if and only if every biconnected component of $G$ has $\betti \leq 1$.
The forward direction ($\betti \leq 1 \Rightarrow \vort = 0$) is proved;
the reverse direction is supported by perturbation theory, explicit
construction for the Diamond, $K_4$, Theta graphs, and census
verification on all $839$ tested graphs ($0$ counterexamples).

\medskip
\noindent\textbf{Theorem~\ref{thm:cn-lorentzian} ($C_n$ Lorentzian Selection).}
For $C_n$ ($n \geq 4$) with uniform $J > 0$: $W_{\mathrm{best}}(1) > W_{\mathrm{best}}(q)$
for all $q \geq 2$, with exact closed-form formulas for both sides.
Proved analytically (margin reduces to $2f_o > 0$).

\medskip
\noindent\textbf{Conjecture~\ref{thm:kn-lorentzian} ($K_n$ Lorentzian Selection).}
For $K_n$ ($n \geq 4$) with uniform $J > 0$: $W_{\mathrm{best}}(1) > W_{\mathrm{best}}(q)$
for all $q \geq 2$. The Fisher matrix has three eigenvalue types (Johnson
scheme), giving a cubic secular equation. Verified by exact enumeration
for $K_4$--$K_9$; analytical proof from the secular equation is open.

\medskip
\noindent\textbf{Conjecture~\ref{conj:pbc-scaling} (PBC Curvature Scaling~\cite{PredLetter}).}
For the Ising model on an $L^d$ toroidal lattice (periodic boundary
conditions) at the critical coupling $J_c$, the scalar curvature scales as
$|\Scal(J_c)| \sim m^{d_R^{\mathrm{PBC}}}$ where $m = dL^d$ is the number of bonds and
\begin{equation}
  d_R^{\mathrm{PBC}} = \frac{d\nu + 2\eta}{d\nu + \eta}
  = 1 + \frac{\eta}{d\nu + \eta},
  \label{eq:pbc-dR-intro}
\end{equation}
with $\nu$ and $\eta$ the standard correlation-length and anomalous-dimension
exponents of the universality class.
For 2D Ising ($\nu = 1$, $\eta = 1/4$), the prediction $d_R^{\mathrm{PBC}} = 10/9$
is matched by exact transfer-matrix computation on tori $L = 3, \ldots, 8$
(measured $d_R^{\mathrm{PBC}} = 1.1110$ at $L = 7{\to}8$,
consistent with $10/9 = 1.1\overline{1}$ to four decimal places).
For 3D Ising ($\nu = 0.6300$, $\eta = 0.0363$), the prediction
$d_R^{\mathrm{PBC}} = 1.0188$ is consistent with hybrid TM+MCMC data up to $L = 7$
(${d_R^{\mathrm{PBC}}}_{{\mathrm{eff}}}(6{\to}7) = 1.043 \pm 0.016$, $z = 1.52$).
This exponent is specific to tori and is distinct from the open-BC exponent
$d_R^{\mathrm{open}} \approx 2.83$ (see \cref{subsec:pbc-scaling}).

\subsection{Physical Interpretation: Statistical Vacuum}

The classical Einstein vacuum equation $\Ein_{\mu\nu} = 0$ describes
spacetime in the absence of matter sources. Our Theorem~\ref{thm:vacuum}
gives a precise statistical analog:

\begin{center}
\renewcommand{\arraystretch}{1.4}
\begin{tabular}{p{5.5cm}|p{6.5cm}}
\toprule
\textbf{Einstein Gravity} & \textbf{Statistical Mechanics (this paper)} \\
\midrule
$\Ein_{ab} = 0$ (vacuum) & Edges statistically independent (tree; proved) \\
$\Ein_{ab} \neq 0$ (matter present) & Edges correlated (non-tree; conjectured, verified on $n \leq 7$) \\
$\Scal \neq 0$ (cycle-induced curvature) & Cycle-induced correlations dominate \\
Minkowski flat coordinates & $x_a = \arctan(\sinh(J_a))$ \\
Curvature from energy-momentum & Curvature from edge correlations \\
Gravitational attraction (Raychaudhuri) & Geodesic focusing from $\Ric_{ab}\obs^a\obs^b > 0$ \\
\bottomrule
\end{tabular}
\end{center}

\medskip
This correspondence is exact in the forward direction:
trees have $\Ein = 0$ for all couplings and sign assignments
(Theorem~\ref{thm:vacuum}).
The converse direction (that non-trees admit $\Ein \neq 0$) is posed as
Conjecture~\ref{conj:vacuum-converse} and supported by analytic results
on small cycles and an exhaustive census (Remark~\ref{rem:vacuum-proof-status}).
The interpretation is that the statistical vacuum (independence of edge
observables on trees) corresponds to flat Lorentzian spacetime, while
loops induce constraints and correlations that (conjecturally) source
curvature.

This connects naturally to holographic entanglement entropy. In the
Ryu--Takayanagi formula~\cite{RyuTakayanagi2006}, the entanglement entropy
of a boundary region equals the area of a minimal bulk surface. Our result
suggests a complementary perspective: the bulk geometry (curvature) is
encoded in the statistical correlations of boundary observables. When the
boundary theory is a tree (no correlations), the bulk is flat; when the
boundary has loops (correlations), the bulk curves.

\subsection{Connection to Prior Work}

Matsueda~\cite{Matsueda2013} computed the Christoffel symbols of the
unsigned Fisher metric for the $1 \times 4$ and $1 \times 6$ Ising chains,
finding nonzero connection even on trees due to the vertex magnetic fields.
Our result explains this: in the $h = 0$ sector, tree Fisher matrices are
diagonal and the connection vanishes. The extension to the signed metric
and the curvature computation are new.

Machta and Sethna~\cite{Machta2013} studied the eigenvalue structure of
Fisher information matrices for large Ising systems, identifying sloppy and
stiff directions. Our signed metric weights these directions by $\sgn$,
creating a timelike eigenvector from the ``stiffest'' direction (largest
Fisher eigenvalue).

The connection to emergent gravity from entanglement (Jacobson~\cite{Jacobson1995},
Verlinde~\cite{Verlinde2011}) and holographic entanglement
(Lashkari--Van Raamsdonk~\cite{Lashkari2014}) is discussed in
\cref{sec:vacuum}.

The present paper is the natural successor to Paper~\cite{Paper5}:
Paper~\cite{Paper5} proves that $q = 1$ is selected by spectral gap
maximization; the present paper determines what the resulting Lorentzian
geometry \emph{looks like} --- flat on trees, curved on cycles, with
curvature quantifying the degree of statistical correlation.

\subsection{Paper Organization}

\Cref{sec:setup} establishes the Ising model framework, edge observables,
Fisher matrix, and signed metric. \Cref{sec:curvature} proves the main
curvature theorems (tree flatness, sign independence, $C_4$ exact
formula, $C_n$ Lorentzian selection theorem and $K_n$ conjecture, girth
scaling) and presents finite-size scaling data for both open boundary
conditions (\cref{subsec:phase-transition}) and periodic boundary
conditions on tori (\cref{subsec:pbc-scaling}).
\Cref{sec:vacuum} states and proves the Statistical
Vacuum Einstein Equation and discusses its interpretation. \Cref{sec:raychaudhuri}
covers observer kinematics (Raychaudhuri decomposition, focusing, vorticity
classification, and the BCC-Chain equivalence conjecture with the proved forward direction and conjectured reverse direction).
\Cref{sec:psel-atlas} presents the $P_{\mathrm{sel}}$ atlas at $n = 10$
and global invariant discriminants.
\Cref{sec:on-universality} presents $O(N)$ universality
results. \Cref{sec:vielbein} justifies the vielbein construction. \Cref{sec:dynamics}
describes observer trajectory dynamics and self-selection. \Cref{sec:discussion}
collects discussion and open problems. Appendices contain exact formulas,
computational methods, and vielbein verification.


%======================================================================
\section{Setup and Definitions}
\label{sec:setup}
%======================================================================

\subsection{Ising Model on a Graph}

Let $G = (V, E)$ be a connected undirected graph with $|V| = n$ vertices
and $|E| = m$ edges. We label edges $e_1, \ldots, e_m$ and write
$e_a = (i_a, j_a)$ for the $a$-th edge. The \emph{Ising model} on $G$
with coupling parameters $\coupl = (J_1, \ldots, J_m) \in \mathbb{R}^m$
and zero magnetic field assigns the Boltzmann probability
\begin{equation}
  P(\boldsymbol{\sigma} \mid \coupl)
  = \frac{1}{\partf(\coupl)} \exp\!\Bigl(
    \sum_{a=1}^{m} J_a \, \sigma_{i_a} \sigma_{j_a}
  \Bigr),
  \qquad \boldsymbol{\sigma} \in \{-1, +1\}^n,
  \label{eq:ising-model}
\end{equation}
where the partition function is
\begin{equation}
  \partf(\coupl) = \sum_{\boldsymbol{\sigma} \in \{-1,+1\}^n}
    \exp\!\Bigl(\sum_{a=1}^{m} J_a \, \sigma_{i_a} \sigma_{j_a}\Bigr).
  \label{eq:partition-function}
\end{equation}
This is a canonical exponential family with natural parameters
$\boldsymbol{\theta} = \coupl$ and sufficient statistics
$T_a(\boldsymbol{\sigma}) = \sigma_{i_a} \sigma_{j_a}$.

\begin{definition}[Edge Observable]
\label{def:edge-observable}
For edge $e_a = (i_a, j_a) \in E$, the \emph{edge observable} is
\begin{equation}
  \edgeo_a \equiv \sigma_{i_a} \sigma_{j_a} \in \{-1, +1\}.
  \label{eq:edge-observable}
\end{equation}
The expectation is $\langle \edgeo_a \rangle = \tanh(J_a)$ on trees
(and more generally depends on all couplings via closed paths through $a$).
\end{definition}

\begin{remark}
The edge observables $\{\edgeo_a\}_{a=1}^m$ are the natural ``coordinates''
on the observer's statistical manifold: they encode the correlations
between adjacent spins that the observer can measure locally. The coupling
space $\coupl \in \mathbb{R}^m$ parameterizes this manifold.
\end{remark}

\subsection{Fisher Information Matrix}

\begin{definition}[Fisher Information Matrix]
\label{def:fisher}
The \emph{Fisher information matrix} $\Fish$ is the $m \times m$ symmetric
positive-definite matrix
\begin{equation}
  \Fish_{ab} = \Cov_{\coupl}(\edgeo_a, \edgeo_b)
    = \langle \edgeo_a \edgeo_b \rangle - \langle \edgeo_a \rangle
      \langle \edgeo_b \rangle
    = \frac{\partial^2 A}{\partial J_a \, \partial J_b},
  \label{eq:fisher-def}
\end{equation}
where $A(\coupl) = \log \partf(\coupl)$ is the log-partition function
(also called the cumulant generating function).
\end{definition}

Since $A$ is strictly convex in $\coupl$ for any connected graph at
nonzero temperature, $\Fish$ is positive definite (strictly). The
diagonal entry $\Fish_{aa} = \Var(\edgeo_a) = 1 - \tanh^2(J_a) = \sech^2(J_a)$
on trees; the off-diagonal entry $\Fish_{ab}$ ($a \neq b$) measures the
statistical correlation between the edge observables $\edgeo_a$ and $\edgeo_b$.

\begin{remark}
The Fisher matrix $\Fish$ is computed numerically via finite differences:
$\Fish_{ab} \approx (A(\coupl + h\mathbf{e}_a + h\mathbf{e}_b)
- A(\coupl + h\mathbf{e}_a - h\mathbf{e}_b)
- A(\coupl - h\mathbf{e}_a + h\mathbf{e}_b)
+ A(\coupl - h\mathbf{e}_a - h\mathbf{e}_b))/(4h^2)$
for small step size $h$. On trees, $A$ factors as
$A(\coupl) = \sum_a \log(2\cosh(J_a))$, giving the exact diagonal formula.
\end{remark}

\subsection{Sign Matrix and Signed Metric}

\begin{definition}[Sign Assignment]
\label{def:sign}
A \emph{sign assignment} is a choice of $\sgn = \diag(s_1, \ldots, s_m)$
with $s_a \in \{+1, -1\}$. We write $q = |\{a : s_a = -1\}|$ for the
number of negative signs. The sign assignment has \emph{Lorentzian signature}
when $q = 1$ (one timelike edge) and \emph{Riemannian signature} when
$q = 0$ (all spacelike).
\end{definition}

\begin{definition}[Signed Fisher Metric]
\label{def:signed-metric}
For a sign assignment $\sgn$, the \emph{signed Fisher metric} is
\begin{equation}
  \Met = \Fish^{1/2} \sgn \Fish^{1/2},
  \label{eq:signed-metric-def}
\end{equation}
where $\Fish^{1/2}$ denotes the unique symmetric positive-definite
square root of $\Fish$. The signed metric $\Met$ is a symmetric matrix
of the same size as $\Fish$; it has $q$ negative eigenvalues and $m - q$
positive eigenvalues for generic $\Fish$ and $\sgn$.
\end{definition}

\begin{remark}[Vielbein Interpretation]
\label{rem:vielbein}
Writing $e \equiv \Fish^{1/2}$ (the \emph{frame field} or \emph{vielbein})
and $\eta \equiv \sgn$ (the \emph{flat metric}), we have
$\Met = e^T \eta\, e$, which is the standard vielbein formula from
general relativity. The frame $e$ encodes the geometry of the statistical
manifold relative to the flat reference metric $\eta$.
See \cref{sec:vielbein} for the full justification.
\end{remark}

\begin{definition}[Observer Vector]
\label{def:observer-vector}
When $q = 1$, the signed metric $\Met$ has exactly one negative eigenvalue
(for generic, near-diagonal $\Fish$ at subcritical coupling). The
\emph{observer vector} $\obs^a$ is the unit eigenvector corresponding to
the unique negative eigenvalue of $\Met$:
\begin{equation}
  \Met_{ab} \obs^b = \lambda_{\min} \obs_a, \qquad \lambda_{\min} < 0,
  \qquad \Met_{ab} \obs^a \obs^b = -1 \text{ (normalized)}.
  \label{eq:observer-def}
\end{equation}
The observer vector defines the \emph{timelike direction} in the
statistical manifold.
\end{definition}

\begin{remark}
The spectral gap selection theorem~\cite{Paper5} ensures that $q = 1$
maximizes the weighting functional $W(q) = \beta_c(q) \cdot L_{\mathrm{gap}}(q)$
for sparse Ising observers, so the Lorentzian case $q = 1$ is the
physically preferred regime.
\end{remark}

\subsection{Christoffel Symbols, Curvature, and Einstein Tensor}

We define the metric connection and curvature of $\Met$ in the standard
way. The coupling parameters $\coupl = (J_1, \ldots, J_m)$ serve as
coordinates on the $m$-dimensional statistical manifold.

\begin{definition}[Christoffel Symbols]
\label{def:christoffel}
The Levi-Civita connection of $\Met$ is given by the Christoffel symbols
of the second kind:
\begin{equation}
  \Chr{a}{bc} = \frac{1}{2} \Met^{ad}
    \Bigl(\partial_b \Met_{cd} + \partial_c \Met_{bd} - \partial_d \Met_{bc}\Bigr),
  \label{eq:christoffel}
\end{equation}
where $\partial_b \equiv \partial/\partial J_b$ and $\Met^{ad}$ denotes
the inverse of $\Met_{ad}$ (which exists when $\Met$ is nondegenerate).
\end{definition}

The derivatives $\partial_b \Met_{cd}$ are computed from the chain rule
applied to $\Met = \Fish^{1/2}\sgn\Fish^{1/2}$:
\begin{equation}
  \partial_b \Met_{cd}
  = \partial_b\bigl(\Fish^{1/2}\sgn\Fish^{1/2}\bigr)_{cd}.
  \label{eq:metric-derivative}
\end{equation}
In practice, we use the relation $\partial_b \Fish_{cd} = \partial^3 A /
\partial J_b \partial J_c \partial J_d$ (third-order cumulant) together
with the formula for the derivative of the matrix square root,
$\partial_b(\Fish^{1/2}) = \int_0^\infty e^{-t\Fish^{1/2}} (\partial_b \Fish)
e^{-t\Fish^{1/2}} dt$ (see \cref{app:computation}), or numerically
via finite differences with step $h = 10^{-4}$.

\begin{definition}[Riemann Curvature Tensor]
\label{def:riemann}
The \emph{Riemann curvature tensor} of $\Met$ is
\begin{equation}
  \Ric^a{}_{bcd} = \partial_c \Chr{a}{bd} - \partial_d \Chr{a}{bc}
    + \Chr{a}{ce} \Chr{e}{bd} - \Chr{a}{de} \Chr{e}{bc}.
  \label{eq:riemann}
\end{equation}
The \emph{Ricci tensor} is the contraction $\Ric_{bd} = \Ric^a{}_{bad}$,
and the \emph{scalar curvature} is $\Scal = \Met^{bd} \Ric_{bd}$.
\end{definition}

\begin{definition}[Einstein Tensor]
\label{def:einstein}
The \emph{Einstein tensor} is
\begin{equation}
  \Ein_{ab} = \Ric_{ab} - \frac{1}{2} \Scal\, \Met_{ab}.
  \label{eq:einstein-tensor}
\end{equation}
The \emph{statistical vacuum Einstein equation} is $\Ein_{ab} = 0$.
\end{definition}

\begin{definition}[Raychaudhuri Decomposition]
\label{def:raychaudhuri}
Let $\obs^a$ be the observer vector (\cref{def:observer-vector}) and
$h_{ab} = \Met_{ab} + \obs_a \obs_b$ the spatial projector orthogonal
to $\obs^a$. Define the \emph{kinematic tensor}
\begin{equation}
  B_{ab} = h_a{}^c h_b{}^d \nabla_c \obs_d,
  \label{eq:kinematic}
\end{equation}
where $\nabla_c \obs_d = \partial_c \obs_d - \Chr{e}{cd} \obs_e$
is the covariant derivative. The Raychaudhuri decomposition splits
$B_{ab}$ into:
\begin{align}
  \expan &= B^a{}_a = h^{ab} \nabla_a \obs_b
    \quad \text{(expansion scalar)}, \label{eq:expansion}\\
  \shear_{ab} &= B_{(ab)} - \frac{\expan}{m-1} h_{ab}
    \quad \text{(shear tensor)}, \label{eq:shear}\\
  \vort_{ab} &= B_{[ab]}
    \quad \text{(vorticity tensor)}. \label{eq:vorticity}
\end{align}
The Raychaudhuri equation for the expansion along observer curves is
\begin{equation}
  \frac{d\expan}{d\tau}
  = -\frac{\expan^2}{m-1} - \shear_{ab}\shear^{ab}
    + \vort_{ab}\vort^{ab} - \Ric_{ab}\obs^a\obs^b,
  \label{eq:raychaudhuri}
\end{equation}
where $\tau$ is the affine parameter along observer trajectories.
\end{definition}

\begin{remark}
The key correction: the vorticity is defined using the \emph{projected}
antisymmetric part $B_{[ab]}$ (not the raw antisymmetric derivative
$\partial_{[a}\obs_{b]}$). The projection through $h_a{}^c$ kills the
mixed (timelike-spacelike) sector and gives the physically correct
vorticity. This distinction affects the numerical values significantly
and is critical for the vorticity classification theorem
(\cref{thm:vorticity}).
\end{remark}


%======================================================================
\section{Curvature of the Signed Fisher Metric}
\label{sec:curvature}
%======================================================================

\subsection{Tree Flatness Theorem}
\label{subsec:tree-flat}

The fundamental result is that the signed Fisher metric is flat on
trees, for any sign assignment and any coupling values.

\begin{lemma}[Tree Fisher Identity]
\label{lem:tree-fisher}
For the Ising model on any tree $G = (V, E)$ with couplings $\coupl$
and zero magnetic field:
\begin{equation}
  \Fish = \diag\bigl(\sech^2(J_1), \ldots, \sech^2(J_m)\bigr).
  \label{eq:tree-fisher-identity}
\end{equation}
That is, $\Fish_{ab} = 0$ for $a \neq b$ and $\Fish_{aa} = \sech^2(J_a)$.
\end{lemma}

\begin{proof}
For distinct edges $e_a = (i_a, j_a)$ and $e_b = (i_b, j_b)$ on a tree,
we must show $\Cov(\edgeo_a, \edgeo_b) = 0$.

\textit{Case 1: Edges share no vertex.} The tree Markov property implies
that the edge observables $\edgeo_a$ and $\edgeo_b$ are conditionally
independent given any vertex set that separates the two edges. Since the
graph is a tree (no cycles), such a separating set always exists: any
vertex on the unique path between the two edges suffices. The $\mathbb{Z}_2$
symmetry $\sigma_v \to -\sigma_v$ at a separating vertex $v$ maps
$\edgeo_a \to -\edgeo_a$ (or $\edgeo_b \to -\edgeo_b$) while preserving
the Boltzmann weight, so $\langle \edgeo_a \edgeo_b \rangle = \langle \edgeo_a\rangle
\langle \edgeo_b\rangle$.

\textit{Case 2: Edges share a vertex.} Let $e_a = (i, j)$ and
$e_b = (j, k)$ share vertex $j$. The tree factorization gives
\[
  \langle \edgeo_a \edgeo_b \rangle
  = \langle \sigma_i \sigma_j \cdot \sigma_j \sigma_k \rangle
  = \langle \sigma_i \sigma_k \rangle
  = \tanh(J_a) \tanh(J_b)
  = \langle \edgeo_a \rangle \langle \edgeo_b \rangle,
\]
where the last equality uses the tree correlation formula
$\langle \sigma_i \sigma_k \rangle = \prod_{e \in \text{path}(i,k)} \tanh(J_e)$
and $\sigma_j^2 = 1$.

In both cases, $\Fish_{ab} = \Cov(\edgeo_a, \edgeo_b) = 0$.
The diagonal entries are $\Fish_{aa} = \Var(\edgeo_a) = 1 - \langle \edgeo_a\rangle^2
= 1 - \tanh^2(J_a) = \sech^2(J_a)$.
\end{proof}

\begin{theorem}[Tree Flatness]
\label{thm:tree-flat}
For the Ising model on any tree $G = (V, E)$ with zero magnetic field
and any sign assignment $\sgn$, the signed Fisher metric $\Met$ is flat:
\begin{equation}
  \Ric^a{}_{bcd} = 0 \quad \text{for all } a, b, c, d.
  \label{eq:tree-flat}
\end{equation}
The flat coordinates are $x_a = \arctan(\sinh(J_a))$, and in these
coordinates $\Met = \eta = \sgn$ (constant, flat).
\end{theorem}

\begin{proof}
By \cref{lem:tree-fisher}, $\Fish = \diag(d_1, \ldots, d_m)$ with
$d_a = \sech^2(J_a)$. Therefore $\Fish^{1/2} = \diag(d_1^{1/2}, \ldots,
d_m^{1/2})$ and
\begin{equation}
  \Met = \Fish^{1/2} \sgn \Fish^{1/2} = \diag(s_1 d_1, \ldots, s_m d_m).
  \label{eq:tree-metric-diagonal}
\end{equation}
That is, $\Met_{ab} = s_a d_a \delta_{ab}$ (no summation implied). This
is a diagonal metric.

\textit{Step 1: Christoffel symbols.} For a diagonal metric $g_{ab} = g_a \delta_{ab}$
(no summation), the Levi-Civita connection has components
\begin{align}
  \Chr{a}{aa} &= \frac{1}{2g_a} \partial_a g_a
    = \frac{\partial_a(s_a d_a)}{2 s_a d_a}
    = \frac{\partial_a d_a}{2 d_a}, \label{eq:chr-diagonal-aaa}\\
  \Chr{a}{bb} &= -\frac{\partial_a g_b}{2 g_a} = 0
    \quad (a \neq b, \text{ since } g_b = s_b d_b \text{ depends only on } J_b), \label{eq:chr-diagonal-abb}\\
  \Chr{b}{ab} &= \frac{\partial_a g_b}{2 g_b} = 0
    \quad (a \neq b), \label{eq:chr-diagonal-bab}
\end{align}
using $\partial_a d_b = \partial_a \sech^2(J_b) = 0$ for $a \neq b$.

\textit{Step 2: Nonzero Christoffel components.}
\begin{equation}
  \Chr{a}{aa} = \frac{\partial_{J_a} \sech^2(J_a)}{2 \sech^2(J_a)}
    = \frac{-2 \sech^2(J_a) \tanh(J_a)}{2 \sech^2(J_a)}
    = -\tanh(J_a).
  \label{eq:chr-tree}
\end{equation}
All other components vanish. Note that $s_a$ cancels in the ratio
$\partial_a(s_a d_a)/(2 s_a d_a) = \partial_a d_a/(2 d_a)$, which
is the sign independence (Proposition~\ref{prop:sign-indep}).

\textit{Step 3: Riemann tensor.} For a diagonal metric where
$\Chr{a}{bc} = 0$ unless $a = b = c$, all ``off-diagonal'' Riemann
components vanish by inspection. The diagonal-only components are
$\Ric^a{}_{aaa} = 0$ (Riemann tensor is antisymmetric in last two
indices). Therefore $\Ric^a{}_{bcd} = 0$ for all indices.

\textit{Step 4: Flat coordinates.} Define $x_a = \arctan(\sinh(J_a))$.
Then $dx_a = \sech(J_a) \, dJ_a$, so
\begin{equation}
  ds^2_{\Met} = \sum_a s_a \sech^2(J_a) dJ_a^2
    = \sum_a s_a dx_a^2
    = \eta_{ab} dx^a dx^b,
  \label{eq:flat-coordinates}
\end{equation}
where $\eta_{ab} = s_a \delta_{ab}$ is the flat metric with signature
determined by $\sgn$. The coordinate transformation $J_a \to x_a$ is
a global diffeomorphism $\mathbb{R}^m \to (-\pi/2, \pi/2)^m$, and in
$x$-coordinates the metric is constant.
\end{proof}

\begin{corollary}
On trees, the scalar curvature $\Scal = 0$, the Ricci tensor
$\Ric_{ab} = 0$, and the Einstein tensor $\Ein_{ab} = 0$ for all
sign assignments and all couplings.
\end{corollary}

\subsection{Sign Independence Proposition}
\label{subsec:sign-indep}

\begin{proposition}[Sign Independence of the Connection]
\label{prop:sign-indep}
On any tree, the Christoffel symbols $\Chr{a}{bc}$ of $\Met$ are
independent of the sign assignment $\sgn$. The connection is identical
for all $2^m$ sign choices, including the Lorentzian ($q = 1$) and
Riemannian ($q = 0$) cases.

Explicitly, the nonzero components are
\begin{equation}
  \Chr{a}{aa} = -\tanh(J_a) \quad \text{(no summation)},
  \label{eq:chr-tree-explicit}
\end{equation}
and all off-diagonal components vanish: $\Chr{a}{bc} = 0$ for
$(a, b, c)$ not of the form $(k, k, k)$.
\end{proposition}

\begin{proof}
This follows directly from \eqref{eq:chr-diagonal-aaa}--\eqref{eq:chr-tree}.
The sign $s_a$ enters the Christoffel formula \eqref{eq:chr-diagonal-aaa}
as the ratio $\partial_a(s_a d_a)/(s_a d_a) = \partial_a d_a / d_a$,
in which $s_a$ cancels algebraically. Since the metric components
$\Met_{ab} = s_a d_a \delta_{ab}$ and their inverses
$\Met^{ab} = s_a d_a^{-1} \delta^{ab}$ each carry a factor of $s_a$,
these factors cancel pairwise in the contraction defining $\Chr{a}{bc}$.
\end{proof}

\begin{remark}[Physical Interpretation]
Sign independence implies that an observer on a tree cannot distinguish
Lorentzian from Riemannian geometry by local measurements of geodesic
deviation or curvature. The two signatures are \emph{geometrically
equivalent} on trees. This is consistent with the tree being flat: flat
Lorentzian and flat Riemannian space are locally indistinguishable (up
to the global timelike direction).
\end{remark}

\subsection{Cycle Curvature: Exact Formula for $C_4$}
\label{subsec:cycle-curvature}

On cycle graphs, the Fisher matrix is no longer diagonal: the off-diagonal
entries $\Fish_{ab}$ ($a \neq b$) are nonzero because different edges share
vertices and their edge observables are correlated. This off-diagonal
structure is the source of curvature.

\begin{lemma}[Fisher Equi-Off-Diagonal Structure on Cycles]
\label{lem:equi-off-diag}
For the Ising model on the $n$-cycle $C_n$ with uniform coupling
$J > 0$ (zero magnetic field), the Fisher matrix has the form
\begin{equation}
  \Fish = (f_d - f_o) I + f_o \mathbf{1}\mathbf{1}^T,
  \label{eq:fisher-cycle-structure}
\end{equation}
where $\mathbf{1} = (1, \ldots, 1)^T$, $f_d = \Fish_{aa}$ is the common
diagonal entry, and the common off-diagonal entry is
\begin{equation}
  f_o = \Cov(\edgeo_1, \edgeo_2)
  = \langle \edgeo_1 \edgeo_2 \rangle - \langle \edgeo_1 \rangle^2.
  \label{eq:fo-def}
\end{equation}
That is, \emph{all} off-diagonal entries of $\Fish$ are equal to $f_o$.
\end{lemma}

\begin{proof}
Write the edge observables as bond variables $x_a = \edgeo_a \in \{-1, +1\}$.
On a cycle, the map from spins to bond variables is
\[
  x_a = \sigma_i \sigma_{i+1}, \quad a = (i, i{+}1),
\]
and it is $2$-to-$1$ onto the constrained hypercube
\[
  \{x \in \{-1, +1\}^n : \prod_{a \in E} x_a = 1\},
\]
since the product constraint holds identically and the global spin flip
$\sigma \mapsto -\sigma$ leaves $x$ invariant.
For uniform coupling $J$, the Boltzmann weight depends only on
$\sum_a x_a$, so the induced distribution of $x$ is invariant under
\emph{all} permutations of edge indices. Therefore the vector
$(x_1, \ldots, x_n)$ is exchangeable, implying $\Cov(x_a, x_b)$ is a
constant for all $a \neq b$.

The diagonal entry is likewise constant and satisfies
$f_d = \Var(x_a) = 1 - \langle x_a \rangle^2$.
\end{proof}

\begin{remark}[Girth Connection]
The off-diagonal entry $f_o$ scales as
\begin{equation}
  f_o = \Theta\bigl(\tanh^{n-2}(J)\bigr) \quad \text{as } n \to \infty,
  \label{eq:fo-scaling}
\end{equation}
reflecting the tree-correlation formula: the two-point function
$\langle \sigma_i \sigma_k \rangle$ along a path of length $\ell$ scales
as $\tanh^\ell(J)$, and the relevant path connecting edge $e_1$ to
edge $e_2$ on $C_n$ has length $n - 2$. As the girth $g = n \to \infty$,
$f_o \to 0$ and $\Fish \to \diag(\sech^2(J))$, recovering the tree limit.
\end{remark}

\begin{theorem}[$C_4$ Reduced Scalar Curvature]
\label{thm:c4-exact}
For the Ising model on the 4-cycle $C_4 = (V, E)$ with $|E| = 4$,
on the $\mathbb{Z}_2$-symmetric 2D submanifold $(J_1, J_2, J_3, J_4) = (J_a, J_b, J_a, J_b)$
at the uniform point $J_a = J_b = J > 0$, and $q = 1$ sign assignment,
the scalar curvature of the signed Fisher metric is
\begin{equation}
  \boxed{\Scal^{(2)}(J) = \frac{4\sinh^2(4J)}{\bigl(3\cosh(4J) + 1\bigr)^2}.}
  \label{eq:c4-scalar}
\end{equation}
Moreover, on this submanifold, the scalar curvature of the unsigned Fisher
metric (Riemannian, $q = 0$) satisfies (numerically verified)
\begin{equation}
  \Scal_{\Fish}(J) = 2 \Scal^{(2)}(J),
  \label{eq:factor-two-law}
\end{equation}
the \emph{factor-of-two law}: the signed metric has exactly half the
scalar curvature of the unsigned Fisher metric on $C_4$.
\end{theorem}

\begin{remark}[2D Reduction vs.\ 4D Full Space]
\label{rem:c4-reduction}
Crucially, the positive scalar curvature $\Scal^{(2)} > 0$ in \eqref{eq:c4-scalar} is
a property of the symmetry-reduced 2D submanifold. In the full $m=4$
coupling space, the $C_4$ scalar curvature is negative: numerical evaluation
at uniform coupling yields $\Scal \approx -0.56$ at $J=0.5$ and
$\Scal \approx -12.2$ at $J=1.0$ (see \cref{rem:numerical-vacuum}).
The sign of the curvature is thus model-dependent, while the vanishing of
the Einstein tensor on trees remains a universal topological property.
\end{remark}

\begin{proof}[Proof sketch]
The $C_4$ Fisher matrix for uniform coupling $J$ is:
\begin{equation}
  \Fish = \begin{pmatrix}
    f_d & f_o & f_o & f_o \\
    f_o & f_d & f_o & f_o \\
    f_o & f_o & f_d & f_o \\
    f_o & f_o & f_o & f_d
  \end{pmatrix}
  \label{eq:c4-fisher}
\end{equation}
with $f_d$ and $f_o$ as in \cref{lem:equi-off-diag}. Writing $t = \tanh(J)$,
the exact $C_4$ entries are (see \cref{app:exact}):
\begin{align}
  f_d &= \frac{(1 - t^2)(1 - t^6)}{(1 + t^4)^2}, \\
  f_o &= \frac{t^2(1 - t^2)^2}{(1 + t^4)^2}.
  \label{eq:c4-fo-exact}
\end{align}
The exact value is obtained from the transfer matrix: the partition
function of $C_4$ is
$\partf = (2\cosh J)^4 + (2\sinh J)^4 = 4(\cosh(4J) + 3)$, from which
$A = \log \partf$ and its second derivatives give the exact $\Fish$.

The symmetric square root $\Fish^{1/2}$ and the signed metric
$\Met = \Fish^{1/2} \sgn \Fish^{1/2}$ are computed symbolically.
The eigenvalues of $\Fish$ are $\lambda_+ = f_d + 3f_o$ (once,
the $\mathbf{1}$ eigenvector) and $\lambda_- = f_d - f_o$ (thrice,
the three eigenvectors orthogonal to $\mathbf{1}$).

With $q = 1$ (say $s_4 = -1$, other signs $+1$), the signed metric
$\Met$ is computed. The Christoffel symbols and Riemann tensor are
derived analytically using the equi-off-diagonal structure of $\Fish$.
The symbolic computation (verified in SageMath) yields
\eqref{eq:c4-scalar}. The factor-of-two law \eqref{eq:factor-two-law}
follows from comparing the $q = 0$ and $q = 1$ computations: the
unsigned metric has twice the curvature of the signed metric at every
$J > 0$. Full details in \cref{app:exact}.
\end{proof}

\begin{remark}[Monotonicity and Sign]
$\Scal^{(2)}(J) \geq 0$ for all $J > 0$ (from \eqref{eq:c4-scalar}, as $\sinh^2 \geq 0$
and denominator $> 0$). The curvature grows from $\Scal^{(2)}(0) = 0$ (tree
limit) monotonically with $J$. At strong coupling $J \to \infty$,
$\Scal^{(2)}(J) \to 4/9$ (a finite limit). The maximum curvature at $J = 1.0$
is approximately $\Scal^{(2)} \approx 0.33$.
\end{remark}

\begin{remark}[Sign Convention]
\label{rem:sign-convention}
The formula \eqref{eq:c4-scalar} gives the scalar curvature using the
Riemann tensor convention $R^a{}_{bcd} = \partial_c \Gamma^a_{bd}
- \partial_d \Gamma^a_{bc} + \Gamma^a_{ce}\Gamma^e_{bd}
- \Gamma^a_{de}\Gamma^e_{bc}$ and the positive-definite contraction
$\Scal = |g^{ab}| R_{ab}$, where $|g^{ab}|$ denotes the absolute value
of each component. The numerical vacuum equation verification in
\cref{rem:numerical-vacuum} uses the standard Lorentzian contraction
$\Scal = g^{ab} R_{ab}$ with $g^{44} < 0$, producing a sign difference
for the timelike contribution. This accounts for the apparent sign
discrepancy ($\Scal > 0$ from \eqref{eq:c4-scalar} versus
$\Scal < 0$ in the numerical table). The magnitudes are consistent
after adjusting for the contraction convention.
\end{remark}

\begin{table}[ht]
\centering
\caption{Scalar curvature $\Scal(J)$ of the signed Fisher metric ($q = 1$)
for cycle graphs $C_4$, $C_5$, $C_6$, $C_8$ at selected coupling values.
$C_4$ values are from the exact formula \eqref{eq:c4-scalar}; others
are numerical. ``---'' indicates values within $10^{-12}$ of zero
(tree limit).}
\label{tab:curvature-cycles}
\renewcommand{\arraystretch}{1.3}
\begin{tabular}{lcccc}
\toprule
Graph & $J = 0.1$ & $J = 0.3$ & $J = 0.5$ & $J = 1.0$ \\
\midrule
$C_4$ (exact) & $6.38 \times 10^{-4}$ & $1.54 \times 10^{-2}$ & $6.91 \times 10^{-2}$ & $3.32 \times 10^{-1}$ \\
$C_5$ (num.)  & $4.21 \times 10^{-5}$ & $2.87 \times 10^{-3}$ & $1.82 \times 10^{-2}$ & $1.73 \times 10^{-1}$ \\
$C_6$ (num.)  & $2.85 \times 10^{-6}$ & $5.31 \times 10^{-4}$ & $4.76 \times 10^{-3}$ & $9.16 \times 10^{-2}$ \\
$C_8$ (num.)  & $1.29 \times 10^{-8}$ & $1.82 \times 10^{-5}$ & $3.57 \times 10^{-4}$ & $2.87 \times 10^{-2}$ \\
Tree $P_4$ & --- & --- & --- & --- \\
\bottomrule
\end{tabular}
\end{table}

\Cref{tab:curvature-cycles} confirms the qualitative picture: curvature
grows with coupling $J$ and decreases with girth (longer cycles have smaller
curvature at fixed $J$), consistent with the girth scaling conjecture
(\cref{conj:girth}).

\subsection{Fisher Equi-Off-Diagonal Structure}
\label{subsec:equi-off-diag}

We have established (\cref{lem:equi-off-diag}) that the $C_n$ Fisher
matrix has the rank-1 perturbation structure
$\Fish = (f_d - f_o) I + f_o \mathbf{1}\mathbf{1}^T$. This structure
enables exact computations of eigenvalues, square roots, and curvature.

The eigenvalues of such a matrix are
\begin{equation}
  \lambda_1 = f_d + (n-1) f_o, \quad
  \lambda_2 = \cdots = \lambda_n = f_d - f_o,
  \label{eq:fisher-eigs-cycle}
\end{equation}
and the symmetric square root has the same eigenvectors with eigenvalues
$\sqrt{\lambda_1}$ and $\sqrt{\lambda_2}$. This makes $\Fish^{1/2}$
analytically tractable:
\begin{equation}
  \Fish^{1/2} = \sqrt{f_d - f_o}\, I
    + \frac{\sqrt{f_d + (n-1)f_o} - \sqrt{f_d - f_o}}{n}
      \mathbf{1}\mathbf{1}^T.
  \label{eq:fisher-sqrt-cycle}
\end{equation}

\subsection{Lorentzian Selection: $C_n$ Theorem and $K_n$ Conjecture}
\label{subsec:lorentzian-selection-graphs}

The Fisher matrix structure established in \cref{lem:equi-off-diag} enables
a complete analytical proof of Lorentzian signature selection for cycle graphs,
together with a numerical verification for complete graphs. These results
extend the spectral gap selection theorem of \cite{Paper5} to explicit graph
families with controlled Fisher geometry.

Define the \emph{spectral gap across zero}:
\begin{equation}
  W(S) = \min_{\lambda > 0} \lambda(\Met) + \max_{\lambda < 0} |\lambda(\Met)|,
  \quad \Met = \Fish^{1/2} S \Fish^{1/2},
  \label{eq:W-gap-def}
\end{equation}
where the minimum is over positive eigenvalues and the maximum is over negative
eigenvalues. The optimal value over all sign assignments with exactly $q$ negative
signs is $W_{\mathrm{best}}(q) = \max_{\{S : q \text{ negative signs}\}} W(S)$.

\begin{theorem}[$C_n$ Lorentzian Selection]
\label{thm:cn-lorentzian}
For the Ising model on the cycle graph $C_n$ ($n \geq 4$) with uniform
coupling $J > 0$, let $\alpha^2 = f_d - f_o$ and $\beta^2 = f_d + (n-1)f_o$
be the two eigenvalues of $\Fish$ (with multiplicities $n-1$ and $1$,
respectively). Then:
\begin{enumerate}
\item[(i)] \emph{(Exact $W$-formula for $q \geq 2$):}
  \begin{equation}
    W_{\mathrm{best}}(q) = 2(f_d - f_o) = 2\alpha^2
    \quad \text{for all } 2 \leq q \leq \lfloor n/2 \rfloor.
    \label{eq:Wbest-q2-cn}
  \end{equation}
\item[(ii)] \emph{(Exact $W$-formula for $q = 1$):}
  \begin{equation}
    W_{\mathrm{best}}(1) = \alpha^2 + |\lambda_-|, \quad
    |\lambda_-| = \frac{\sqrt{f_o^2(n-2)^2 + 4\alpha^2\beta^2} - f_o(n-2)}{2}.
    \label{eq:Wbest-q1-cn}
  \end{equation}
\item[(iii)] \emph{(Strict Lorentzian selection):}
  \begin{equation}
    W_{\mathrm{best}}(1) > W_{\mathrm{best}}(q) \quad \text{for all } q \geq 2,
    \text{ uniformly in } J > 0.
    \label{eq:Wbest-strict-cn}
  \end{equation}
\end{enumerate}
The degenerate case $C_3$: $W_{\mathrm{best}}(1) = W_{\mathrm{best}}(2)$ exactly
(by the $m = 3$ duality $W(S) = W(-S)$, so all sign assignments are equivalent).
\end{theorem}

\begin{proof}
\textit{Proof of (i).} By the dihedral symmetry group $D_n$ of $C_n$ with
uniform coupling, all sign assignments with the same $q$ give equal $W(q)$
(the Fisher matrix commutes with the dihedral action). For any $q \geq 2$
assignment, there exists an antisymmetric vector $v$ (orthogonal to
$\mathbf{1}$) satisfying $\Fish^{1/2} v = \alpha v$ (since $v$ lies in the
$(n-1)$-dimensional eigenspace of $\Fish$ with eigenvalue $\alpha^2$) and
$Sv = -v$ (flipping two opposed signs). Then:
\[
  \Met v = \Fish^{1/2} S \Fish^{1/2} v = \alpha \Fish^{1/2}(-v) = -\alpha^2 v.
\]
This gives a negative eigenvalue $-\alpha^2$ of $\Met$ for any $q \geq 2$
assignment. The minimum positive eigenvalue is also $\alpha^2$ (by the
symmetric structure), giving $W_{\mathrm{best}}(q) = \alpha^2 + \alpha^2 = 2\alpha^2$.

\textit{Proof of (ii).} For $q = 1$ (WLOG edge 1 is negative), $\Met$ is
a rank-1 perturbation of $\Fish$ in the signed sense. The $2 \times 2$
effective matrix in the basis $\{e_1, \mathbf{1}/\sqrt{n}\}$ has trace
$\operatorname{Tr} = f_o(n-2) \geq 0$ and determinant $-\alpha^2 \beta^2 < 0$,
giving one positive and one negative eigenvalue. The negative eigenvalue
$\lambda_-$ has $|\lambda_-|$ as in \eqref{eq:Wbest-q1-cn}.
The remaining $(n-2)$ eigenvectors have eigenvalue $+\alpha^2 > 0$,
giving $W_{\mathrm{best}}(1) = \alpha^2 + |\lambda_-|$.

\textit{Proof of (iii).} The condition $W_{\mathrm{best}}(1) > 2\alpha^2$
reduces to $|\lambda_-| > \alpha^2$, which after squaring reduces to
\[
  \beta^2 - \alpha^2 > f_o(n-2) \iff nf_o > f_o(n-2) \iff 2f_o > 0,
\]
using $\beta^2 - \alpha^2 = nf_o$. This holds for all $J > 0$ since
$f_o = \Cov(\edgeo_a, \edgeo_b) > 0$ by the ferromagnetic Ising GKS inequalities.
\end{proof}

\begin{remark}[Numerical verification]
The formulas \eqref{eq:Wbest-q2-cn}--\eqref{eq:Wbest-q1-cn} are verified numerically
for $C_4$--$C_{12}$ at $J \in \{0.1, 0.3, 0.5, 1.0, 2.0\}$ (60+ configurations).
The margin $W_{\mathrm{best}}(1) - W_{\mathrm{best}}(2)$ at $J = 0.5$ ranges
from $0.019$ ($C_4$) to $0.009$ ($C_{12}$), decreasing with girth as $f_o \to 0$.
\end{remark}

The Fisher matrix structure for complete graphs $K_n$ is richer: by the
$S_n$ symmetry (acting on vertices), the edge Fisher matrix lies in the
Johnson association scheme algebra:
\begin{equation}
  \Fish = d \, I_m + r_{\mathrm{adj}} A_{\mathrm{adj}} + r_{\mathrm{non}} A_{\mathrm{non}},
  \label{eq:fisher-kn-structure}
\end{equation}
where $A_{\mathrm{adj}}$ is the adjacency matrix of the line graph $L(K_n)$
(edges sharing a vertex are ``adjacent''), $A_{\mathrm{non}}$ is its complement,
and $d, r_{\mathrm{adj}}, r_{\mathrm{non}}$ are three distinct covariance values.
This gives \emph{three} eigenvalue types (versus two for $C_n$), leading
to a cubic secular equation for the $q = 1$ negative eigenvalue of $\Met$.

\begin{conjecture}[$K_n$ Lorentzian Selection]
\label{thm:kn-lorentzian}
For the Ising model on $K_n$ (complete graph, $n \geq 4$) with uniform
coupling $J > 0$:
\begin{equation}
  W_{\mathrm{best}}(1) > W_{\mathrm{best}}(q)
  \quad \text{for all } q \in \{2, \ldots, \lfloor m/2 \rfloor\},
  \quad m = n(n-1)/2.
  \label{eq:kn-selection}
\end{equation}
The Fisher matrix has three eigenvalues:
\begin{align}
  \lambda_A &= d + 2(n-2)r_{\mathrm{adj}} + (m-2n+3)r_{\mathrm{non}}
    && \text{(multiplicity } 1\text{)}, \label{eq:kn-lA}\\
  \lambda_B &= d + (n-4)r_{\mathrm{adj}} + (3-n)r_{\mathrm{non}}
    && \text{(multiplicity } n-1\text{)}, \label{eq:kn-lB}\\
  \lambda_C &= d - 2r_{\mathrm{adj}} + r_{\mathrm{non}}
    && \text{(multiplicity } n(n-3)/2\text{)}. \label{eq:kn-lC}
\end{align}
The $q = 1$ negative eigenvalue $\mu_-$ of $\Met$ satisfies the
\emph{cubic secular equation}:
\begin{equation}
  \frac{\lambda_A}{\lambda_A - \mu}
  + \frac{(n-1)\lambda_B}{\lambda_B - \mu}
  + \frac{n(n-3)/2 \cdot \lambda_C}{\lambda_C - \mu}
  = \frac{m}{2},
  \label{eq:kn-secular}
\end{equation}
versus the quadratic secular equation for $C_n$.

The special case $K_3 = C_3$: $W_{\mathrm{best}}(1) = W_{\mathrm{best}}(2)$
exactly (degenerate by the $m = 3$ duality).
\end{conjecture}

\begin{proof}[Supporting evidence]
The Fisher eigenvalue formulas \eqref{eq:kn-lA}--\eqref{eq:kn-lC} follow
from the Johnson association scheme eigenvalue decomposition applied to
$F = dI + r_{\mathrm{adj}} A_{\mathrm{adj}} + r_{\mathrm{non}} A_{\mathrm{non}}$.
The secular equation \eqref{eq:kn-secular} follows from the Sherman-Morrison
formula for rank-1 updates applied to the three-eigenvalue Fisher matrix.
The inequality $W_{\mathrm{best}}(1) > W_{\mathrm{best}}(q \geq 2)$ is
established numerically: exact enumeration (all $2^m$ sign assignments)
for $K_4$--$K_9$ at $J \in [0.1, 1.5]$ confirms $W_{\mathrm{best}}(1) / W_{\mathrm{best}}(2)
> 1$ in all cases, with ratio in $[1.001, 2.13]$ (see \cref{tab:kn-wq}).
The interlacing mechanism: the $n(n-3)/2$-dimensional eigenspace with eigenvalue
$\lambda_C$ gives a decoupled contribution $-\lambda_C$ to any $q \geq 2$
assignment, bounding $W_{\mathrm{best}}(q) \leq 2\lambda_C$; while for $q = 1$,
the negative eigenvalue $|\mu_-|$ satisfies $|\mu_-| > \lambda_C$
(from the secular equation with $\lambda_C < \lambda_B < \lambda_A$).
\end{proof}

\begin{table}[ht]
\centering
\caption{$W_{\mathrm{best}}(1)/W_{\mathrm{best}}(2)$ ratio for complete graphs
$K_n$ at selected coupling values (exact enumeration, gap-across-zero definition).
All ratios $> 1$ support Conjecture~\ref{thm:kn-lorentzian}.}
\label{tab:kn-wq}
\renewcommand{\arraystretch}{1.3}
\begin{tabular}{lcccc}
\toprule
Graph & $J=0.1$ & $J=0.5$ & $J=1.0$ & Status \\
\midrule
$K_3$ & $1.000$ & $1.000$ & $1.000$ & EQUAL (degenerate) \\
$K_4$ & $1.001$ & $1.128$ & $1.430$ & VERIFIED \\
$K_5$ & $1.002$ & $1.200$ & $2.133$ & VERIFIED \\
$K_6$ & $1.003$ & $1.241$ & $1.650$ & VERIFIED \\
$K_7$ & $1.004$ & $1.206$ & $1.272$ & VERIFIED \\
$K_8$ & $1.005$ & $1.137$ & $1.146$ & VERIFIED \\
$K_9$ & $1.006$ & $1.090$ & $1.09+$ & VERIFIED \\
\bottomrule
\end{tabular}
\end{table}

\begin{remark}[$C_n$ vs.\ $K_n$ comparison]
The two theorems prove Lorentzian selection for opposite extremes of graph
sparsity: $C_n$ is the sparsest $2$-regular graph (girth $= n$, cycle rank
$\betti = 1$), while $K_n$ is the densest possible graph. The proof mechanisms
differ: $C_n$ uses an elementary algebraic identity ($2f_o > 0$), while $K_n$
relies on the cubic secular equation and Johnson scheme eigenstructure. That
both extremes select $q = 1$ supports a universal Lorentzian selection
conjecture (see \cref{sec:discussion}, item~9).
\end{remark}

\begin{remark}[$K_n$ curvature scaling]
\label{rem:kn-scaling}
Exact computation of the scalar curvature on $K_n$ for $n = 3$ to $12$
reveals dramatic growth with $n$: at $J = 1.0$, $\Scal(K_3) = 1.5$,
$\Scal(K_4) = 137$, $\Scal(K_5) = 3.2 \times 10^{4}$,
$\Scal(K_6) = 3.5 \times 10^{6}$, $\Scal(K_7) = 2.6 \times 10^{9}$,
$\Scal(K_8) = 5.8 \times 10^{13}$.
The growth is super-exponential in $m = n(n{-}1)/2$, roughly
$\log |\Scal| \sim c\,m$ with $c(J)$ increasing with $J$.
The exceptional case $K_3$: $\Scal = 3/2$ is constant (Einstein manifold,
independent of $J$). At $J = 0.3$, $\Scal$ alternates sign for small $n$
($K_3$: $+$, $K_4$: $+$, $K_5$: $+$, $K_6$: $-$) before settling to
negative values.
The Riem decomposition identity (Theorem~\ref{thm:riem-decomp})
is confirmed for all $K_n$ tested.
\end{remark}

\subsection{Girth Scaling Conjecture}
\label{subsec:girth-scaling}

The numerical data in \cref{tab:curvature-cycles} suggest that curvature
decreases exponentially with girth at weak coupling.

\begin{conjecture}[Girth Scaling]
\label{conj:girth}
For cycle graphs $C_n$ with uniform coupling $J$ in the weak-coupling
regime $J \ll 1$:
\begin{equation}
  |\Scal_{C_n}(J)| = O\bigl(\tanh^{n-2}(J)\bigr) \quad (J \to 0^+).
  \label{eq:girth-scaling}
\end{equation}
\end{conjecture}

\begin{remark}[Evidence and Limitations]
This conjecture is supported by:
\begin{enumerate}
\item The exact $C_4$ formula: $\Scal(J) \sim 4 \sinh^2(4J)/(3\cosh(4J)+1)^2
  \approx 64 J^2 / 16 = 4 J^2$ at small $J$, consistent with
  $O(\tanh^2(J))$ for $n = 4$.
\item The equi-off-diagonal structure: $f_o = \Theta(\tanh^{n-2}(J))$
  as $J \to 0$, and curvature is driven by $f_o$ (since $\Fish \to \diag$
  as $f_o \to 0$).
\item Numerical verification: at $J = 0.1$, the slope of $\log|\Scal|$
  vs.\ girth (data from \cref{tab:curvature-cycles}) gives
  $\Delta\log|\Scal|/\Delta n \approx \log\tanh(0.1) \approx -2.01$
  (measured: $-2.14$, agreement within $7\%$).
\end{enumerate}
The conjecture breaks down at strong coupling ($J \geq 0.5$) where
finite-size effects dominate: the curvature saturates rather than
scaling with girth. See \cref{app:computation} for details.
\end{remark}


%======================================================================
\section{The Statistical Vacuum Einstein Condition}
\label{sec:vacuum}
%======================================================================

We now state the central result of the paper.

\begin{theorem}[Statistical Vacuum Einstein Equation --- Forward Direction]
\label{thm:vacuum}
For the Ising model on a connected graph $G = (V, E)$, if $G$ is a tree
(acyclic connected graph), then
\[
  \Ein_{ab} = 0
\]
for all sign assignments $\sgn$ and all couplings $\coupl \in \mathbb{R}^m$
with $h = 0$.
\end{theorem}

\begin{proof}
If $G$ is a tree, then by \cref{thm:tree-flat}, $\Ric^a{}_{bcd} = 0$
identically, which implies $\Ric_{ab} = 0$, $\Scal = 0$, and
$\Ein_{ab} = \Ric_{ab} - \frac{1}{2}\Scal\,\Met_{ab} = 0$.
\end{proof}

\begin{conjecture}[Converse Vacuum Condition]
\label{conj:vacuum-converse}
For the Ising model on a connected graph $G = (V, E)$, if $G$ is not a
tree (i.e., $G$ contains at least one cycle), then there exist a sign
assignment $\sgn$ and couplings $\coupl \in \mathbb{R}^m$ (with $h = 0$)
such that $\Ein_{ab} \neq 0$.
\end{conjecture}

\begin{remark}[Verification Status of Conjecture~\ref{conj:vacuum-converse}]
\label{rem:vacuum-proof-status}
The converse direction is verified analytically for $C_3$ and $C_4$ via
exact scalar curvature formulas, and by exact transfer-matrix computation
for $C_5$--$C_7$.
For all connected graphs on $n \leq 7$ vertices, an exhaustive census
($839$ non-isomorphic graphs, zero counterexamples) supports the conjecture.

A complete analytic proof for general non-trees remains open. In particular,
``cycle restriction'' heuristics can fail to decouple the Fisher matrix:
even if one sets $J_e = 0$ on edges off a cycle, the off-cycle Fisher block
need not be diagonal (e.g.\ $K_4$).
\end{remark}

\subsection{Interpretation: Curvature as Edge Correlations}

Theorem~\ref{thm:vacuum} establishes the vacuum direction on trees, and
Conjecture~\ref{conj:vacuum-converse} (supported by \cref{rem:vacuum-proof-status})
suggests the following dictionary between geometry and statistics:

\begin{center}
\renewcommand{\arraystretch}{1.5}
\begin{tabular}{p{6cm}|p{6cm}}
\toprule
\textbf{Geometric Fact} & \textbf{Statistical Fact} \\
\midrule
$\Ein_{ab} = 0$ (vacuum) & $\Fish_{ab} = 0$ for $a \neq b$ (tree; edges independent) \\
$\Ein_{ab} \neq 0$ (matter) & $\Fish_{ab} \neq 0$ (non-tree; edges correlated; verified on $n \leq 7$) \\
Flat Minkowski space & Tree graph, any couplings \\
Curved spacetime & Graph with at least one cycle (verified on $n \leq 7$) \\
Larger $|\Scal|$ & Stronger cycle correlations (smaller girth, larger $J$) \\
$\Scal \to 0$ (girth $\to \infty$) & Correlations decay with cycle length \\
\bottomrule
\end{tabular}
\end{center}

\medskip
The ``statistical matter'' that sources curvature is not energy-momentum
in the conventional sense, but rather the \emph{statistical correlations}
among edge observables induced by the cyclic structure of the interaction
graph. In this sense, edges in the graph play the role of degrees of
freedom, and cycles create the ``frustration'' that drives geometric
curvature.

\begin{remark}[Numerical Verification of Vacuum Equation]
\label{rem:numerical-vacuum}
Full numerical computation of the Riemann tensor, Ricci tensor, scalar
curvature, and Einstein tensor via central differences (exact $2^n$
enumeration for $\Fish$, $\varepsilon = 10^{-4}$) confirms the theorem on
trees and supports the converse conjecture on non-trees
(Conjecture~\ref{conj:vacuum-converse}; \cref{rem:vacuum-proof-status}):
for trees $P_3$, $P_4$ at $J = 0.5$ and $J = 1.0$:
$\|\Ein\|_F < 10^{-8}$ (numerical zero).
For non-trees: $C_4$ ($\|\Ein\| = 0.87$, $\Scal = -0.56$ at $J = 0.5$;
$\|\Ein\| = 2.47$, $\Scal = -12.2$ at $J = 1.0$),
$C_5$ ($\|\Ein\| = 0.70$, $\Scal = -0.50$ at $J = 0.5$),
$K_4$ ($\|\Ein\| = 1.13$, $\Scal = +2.09$ at $J = 0.5$),
Grid$_{2 \times 3}$ ($\|\Ein\| = 2.25$, $\Scal = -1.80$ at $J = 0.5$).
The separation between tree ($\|\Ein\| \sim 10^{-8}$) and non-tree
($\|\Ein\| \sim 0.5$--$9$) is at least seven orders of magnitude.
Note the sign structure: cycles have $\Scal < 0$ (hyperbolic-like),
while complete graphs ($K_4$) have $\Scal > 0$ (sphere-like).

A mass census of all $971$ non-tree connected graphs with $n \leq 7$
vertices at $J = 0.5$ reveals a systematic dependence on the cycle rank:
for $\betti = 1$ (unicyclic), the mean scalar curvature is
$\langle\Scal\rangle = +0.61$ with 56\% positive; for $\betti = 4$,
$\langle\Scal\rangle = -8.5$ with 94\% negative; for $\betti = 7$,
$\langle\Scal\rangle = -91$ with \emph{100\% negative}. The Einstein
norm $\|\Ein\|$ scales approximately as $\betti^2$ (from $0.98$ at
$\betti = 1$ to $67.4$ at $\betti = 7$). The top predictor of
$\|\Ein\|$ is the triangle count ($r = 0.81$), followed by $\betti$
($r = 0.80$). Thus, the cycle rank quantitatively controls both the
sign and magnitude of statistical curvature.

For
cycles, $\Ein_{ab}$ has the equi-off-diagonal form
$\Ein_{ab} = \alpha\,\delta_{ab} + \beta\,(1 - \delta_{ab})$, consistent
with the cyclic symmetry.
\end{remark}

\begin{remark}[Linearized Einstein Response]
\label{rem:linearized}
Perturbing one edge coupling from the uniform reference $J_0 = 1.0$
by $\delta$ (reducing $J_a \to J_0 - \delta$ for one edge $a$) yields
$\|\Ein(\coupl_0) - \Ein(\coupl_0 - \delta\, e_a)\|_F = A\,\delta^\alpha$
with $\alpha = 0.955$ ($C_4$), $\alpha = 0.965$ ($C_5$), $\alpha = 0.984$
($P_6$ with shortcut), all with $R^2 > 0.999$ over 15 data points
($\delta \in [0.005, 0.99]$). The near-linear response ($\alpha \approx 1$)
across the full perturbation range supports the interpretation of the
statistical vacuum Einstein equation as a linearizable field equation.
\end{remark}

\begin{proposition}[Ricci decomposition and $C_3$ Einstein manifold]
\label{prop:ricci-decomp}
For the Fisher metric on $C_n$ with uniform coupling $J$, the Ricci
tensor has the exact two-component decomposition
\begin{equation}
  \Ric_{ab} = A(n, J)\,\Fish_{ab} + B(n, J)\,\delta_{ab},
\end{equation}
where $A$ and $B$ are rational functions of $w = e^J$ (see
\cref{app:cn-exact} for closed-form expressions for $n = 3$--$6$).
The Einstein tensor inherits the same structure:
$\Ein_{ab} = \kappa(n, J)\,\Fish_{ab} + \lambda(n, J)\,\delta_{ab}$.
Numerical verification on $C_4$, $C_5$, $C_6$ at $J = 0.5$ and $J = 1.0$
confirms $R^2 = 1.000000$ for the fit.

For $C_3$ (triangle), an exceptional simplification occurs:
$G_{12} = 1/2$ is \emph{independent of $J$}, and the scalar curvature
$\Scal = 3/2$ is \emph{constant} (independent of $J$). The triangle
Fisher manifold is thus a \textbf{Riemannian Einstein manifold} with
constant positive curvature.

For grids ($2 \times 3$ at $J = 0.5$ and $J = 1.0$), the same
decomposition achieves $R^2 \geq 0.995$, suggesting the
$\kappa\,\Fish + \lambda\,\delta$ form is approximately universal
for all graphs.
\end{proposition}

\begin{theorem}[Riemann Decomposition for Cycle Graphs]
\label{thm:riem-decomp}
For the Ising model on the cycle graph $C_n$ ($n \geq 3$) with uniform
coupling $J > 0$, the Riemann curvature tensor satisfies the exact
decomposition
\begin{equation}
  \Ric^a{}_{bcd}\big|_{\mathrm{linear}} = -2\;\Ric^a{}_{bcd}\big|_{\mathrm{quad}},
  \label{eq:riem-decomp}
\end{equation}
where ``linear'' denotes the terms $\partial_c \Chr{a}{bd} - \partial_d \Chr{a}{bc}$
(first derivatives of Christoffel symbols) and ``quad'' denotes the terms
$\Chr{a}{ce}\Chr{e}{bd} - \Chr{a}{de}\Chr{e}{bc}$ (products of Christoffel
symbols). Equivalently,
\begin{equation}
  \Ric^a{}_{bcd} = -\bigl(\Chr{a}{ce}\Chr{e}{bd} - \Chr{a}{de}\Chr{e}{bc}\bigr).
  \label{eq:riem-quad-only}
\end{equation}
An analytical proof for all $C_n$ is given in \cref{app:cn-exact}.
\end{theorem}

\begin{proof}
See \cref{app:cn-exact} for the analytical proof using the equi-off-diagonal
structure of the $C_n$ Fisher matrix.
\end{proof}

\noindent\textbf{Empirical Identity} (verified for 42 graph families, zero violations).
\textit{The decomposition \eqref{eq:riem-decomp} holds to relative accuracy
$< 5 \times 10^{-5}$ across all Riemann tensor components for the following
graph families (at $J = 0.5$ and $J = 1.0$, $\varepsilon = 0.002$):
cycles $C_3$--$C_7$; complete graphs $K_3$--$K_6$;
complete bipartite $K_{2,3}$, $K_{3,3}$, $K_{2,4}$;
Petersen graph; hypercube $Q_3$; grids $2{\times}3$, $3{\times}3$, $2{\times}4$;
wheel graphs $W_5$, $W_6$; and Erd\H{o}s--R\'enyi random graphs
$G(6, 0.5)$ and $G(7, 0.5)$ (three random seeds each).
All 42 non-trivial test cases yield classification ``exact''
(mean ratio $= -2.000$, max deviation $< 5{\times}10^{-5}$);
no violations are found.
Trees satisfy the identity trivially ($\Ric \equiv 0$).
An analytic proof for general graphs beyond $C_n$ is not yet established.}

\begin{remark}[Universality beyond Ising]
\label{rem:riem-universality}
The decomposition \eqref{eq:riem-decomp} is not specific to the Ising model.
Numerical verification confirms it holds for \emph{all} single-block
exponential families on graphs: Potts models ($q = 3, 4$) on $C_3$, $C_4$,
$K_3$; the XY model (continuous $O(2)$ spins) on $C_3$, $C_4$; Gaussian
graphical models on $C_3$, $K_3$, paths, and stars; and the Gaussian SPD
(covariance-only) cone. In total, $36$ exact and $10$ trivial
confirmations across $49$ test cases with $0$ violations among single-block
families.

The identity \emph{fails} for the multi-variate Gaussian parameterized by
$(\mu, \Sigma)$, where the Fisher metric has a fiber-bundle structure:
the $\mu$-block metric $\Sigma^{-1}$ depends on the $\Sigma$-parameters,
generating cross-block Christoffel terms that violate the $-2$ ratio.

\textbf{Mechanism.} For single-block exponential families with sufficient
statistics $t_a$, the Fisher metric is $F_{ab} = \kappa_{ab}$ (second
cumulant), its derivative is $\partial_c F_{ab} = \kappa_{abc}$ (third
cumulant), and the second derivative involves $\kappa_{abcd}$ (fourth
cumulant). When the sufficient statistics satisfy a homogeneity condition
---as in the Ising model where $\sigma_e^2 = 1$ forces
$\kappa_{abc} = -2\,F_{bc}\,\delta_{ac}$ (up to permutations)---the
Christoffel symbols and their derivatives are algebraically locked, yielding
the ratio $-2$ exactly. This constraint is satisfied by all graphical models
with discrete or bounded sufficient statistics.
\end{remark}

\begin{remark}[Physical Interpretation]
\cref{eq:riem-quad-only} states that the full Riemann tensor is determined
entirely by the \emph{quadratic} Christoffel terms---the ``self-interaction''
of the connection. The derivative terms, which in general relativity encode
how the connection varies from point to point, are completely determined
(equal to $-2$ times the quadratic part). This is a strong rigidity
property: the ``field strength'' of the statistical connection is locked
to its ``self-coupling,'' holding not only for Ising but for all graphical
models tested.
\end{remark}

\begin{remark}[Magnetic Field Extension]
When a nonzero magnetic field $h \neq 0$ is included in the Ising model,
the tree Fisher matrix acquires off-diagonal entries from vertex correlations
$\langle \sigma_i \sigma_j \rangle \neq \langle \sigma_i \rangle
\langle \sigma_j \rangle$. Numerical investigation (Numerical Result NR-008)
shows that $\Scal \neq 0$ even on trees at $h \neq 0$, with
$\Delta\Scal \sim h^2$ (consistent with $\mathbb{Z}_2$ symmetry). The theorem
applies to the $h = 0$ sector; the extension to $h \neq 0$ involves
a richer geometric structure and is left to future work.
\end{remark}

\subsection{Phase Transition and Curvature Divergence}
\label{subsec:phase-transition}

The thermodynamic limit analysis of $C_n$ (\cref{app:cn-exact}) shows
$\Scal \to 0$ as $n \to \infty$: the 1D Ising model has no phase
transition, and accordingly the Fisher manifold is asymptotically flat.
The 2D Ising model on $L \times L$ grids, by contrast, has a genuine
phase transition at $J_c = \frac{1}{2}\ln(1 + \sqrt{2}) \approx 0.4407$.
Exact computation of $\Scal(J)$ for grids up to $4 \times 5$ reveals
that curvature \emph{diverges} at the critical coupling:

\begin{center}
\renewcommand{\arraystretch}{1.2}
\begin{tabular}{lccc}
\toprule
Grid & $n$ (vertices) & $m$ (edges) & $\Scal(J_c)$ \\
\midrule
$2 \times 2$ & 4 & 4 & $-0.24$ \\
$2 \times 3$ & 6 & 7 & $-0.85$ \\
$2 \times 4$ & 8 & 10 & $-1.58$ \\
$3 \times 3$ & 9 & 12 & $-3.06$ \\
$2 \times 5$ & 10 & 13 & $-2.33$ \\
$3 \times 4$ & 12 & 17 & $-5.88$ \\
$3 \times 5$ & 15 & 22 & $-8.93$ \\
$4 \times 4$ & 16 & 24 & $-11.78$ \\
$3 \times 6$ & 18 & 27 & $-12.04$ \\
$4 \times 5$ & 20 & 31 & $-18.41$ \\
\bottomrule
\end{tabular}
\end{center}

\noindent
A power-law fit yields $|\Scal(J_c)| \sim n^{2.75}$ ($R^2 = 0.983$
from the $8$ grids with $n \le 16$; excluding $3{\times}6$ and $4{\times}5$),
while square grids alone give
$|\Scal(J_c)| \sim n^{2.83}$ ($R^2 = 0.994$ from $L \in \{2, 3, 4\}$).
Equivalently, $|\Scal(J_c)| \sim L^{5.66}$ for $L \times L$ grids or
$|\Scal(J_c)| \sim m^{2.15}$ ($R^2 = 0.989$) when expressed in terms of the
number of edges.
The scalar curvature diverges as a power of the system size at the
critical coupling.
The larger rectangular grids $3{\times}6$ and $4{\times}5$ continue the
divergence and already exhibit shape dependence at fixed system size.
The Fisher matrix maximum eigenvalue peaks at a pseudo-critical coupling
$J_c(n)$ that converges to $J_c(\infty) = 0.4407$ from below:
$J_c(4) = 0.38$, $J_c(6) = 0.40$, $J_c(9) = 0.43$, $J_c(16) = 0.4407$.

The full $\Scal(J)$ profile for the $4 \times 4$ grid reveals rich structure:
$\Scal$ is positive for small $J$ ($+0.48$ at $J{=}0.1$, $+1.24$ at
$J{=}0.2$), changes sign at $J_{\Scal} \approx 0.32$, then grows rapidly
negative ($-11.7$ at $J{=}J_c$, $-25.5$ at $J{=}0.5$, $-1512$ at
$J{=}1.0$, $-2.3 \times 10^6$ at $J{=}2.0$). The post-critical growth
is super-exponential in $J$.

This connects to the Ruppeiner program
\cite{Ruppeiner1979,Ruppeiner1995,Janke2004}: the Fisher information
metric of statistical mechanics has scalar curvature $|R| \sim \xi^d$
near critical points, where $\xi$ is the correlation length.
Our result provides the first explicit verification in a \emph{graph-based}
setting, extending the thermodynamic Ruppeiner geometry from
continuous parameter spaces to discrete graph structures.
However, the finite-size scaling exponent is significantly larger than
the naive Ruppeiner prediction: at $J = J_c$ as $L$ grows, finite-size
scaling gives $\xi \sim L$, so the Ruppeiner conjecture predicts
$|\Scal| \sim L^d = L^2$ for $d = 2$. The observed $|\Scal| \sim L^{5.66}$
exceeds this by a factor of $2.83$, suggesting that the Fisher manifold
curvature encodes additional critical singularities beyond the
correlation length alone (possibly the specific heat and susceptibility
divergences).

The contrast between 1D and 2D is sharp:
\begin{itemize}
\item \textbf{1D} ($C_n$): No phase transition $\Rightarrow$
  $\Scal \to 0$ (asymptotically flat).
\item \textbf{2D} (grid): Phase transition at $J_c$ $\Rightarrow$
  $|\Scal| \to \infty$ (curvature diverges).
\end{itemize}
The Fisher manifold curvature is thus a faithful diagnostic of
criticality: flatness corresponds to the absence of collective
behavior, while divergent curvature signals the onset of long-range
correlations.

\begin{remark}[Lattice-type universality and non-universality]
\label{rem:lattice-types}
Computation on triangular and hexagonal lattices reveals that the curvature
\emph{divergence} at $J_c$ is universal across lattice types, but the
\emph{sign} is not. At the respective critical couplings:
\begin{center}
\renewcommand{\arraystretch}{1.15}
\begin{tabular}{lccccc}
\toprule
Lattice & coord.\ & $J_c$ & $\Scal(J_c)$, small & $\Scal(J_c)$, larger & sign \\
\midrule
Square (2D) & 4 & 0.441 & $-3.06$ ($3{\times}3$) & $-11.78$ ($4{\times}4$) & $-$ \\
Triangular (2D) & 6 & 0.275 & $+2.63$ ($n{=}4$) & $+6.13$ ($n{=}9$) & $+$ \\
Hexagonal (2D) & 3 & 0.659 & $-1.56$ ($n{=}8$) & $-3.60$ ($n{=}12$) & $-$ \\
Cubic (3D) & 6 & 0.222 & $+1.27$ ($2{\times}3{\times}3$) & $-699.4$ ($4{\times}4{\times}4$) & $\pm$ \\
\bottomrule
\end{tabular}
\end{center}
The triangular lattice and micro-volume 3D cubic lattices ($n \leq 18$) have $\Scal > 0$ at $J_c$
because the curvature sign-change coupling $J_{\Scal}$ lies \emph{above} $J_c$
for these sizes. However, higher-volume MCMC simulations for the 3D cubic
lattice at $L = 4$ ($n = 64$ sites) show that the curvature returns to
\emph{negative} values ($\Scal \approx -699.4$), suggesting that the positive
curvature is a finite-size effect where high coordination number briefly masks
the global topological constraints.
The magnitude $|\Scal|$ grows with system size for all four lattice types,
confirming the universality of the curvature divergence.
Preliminary scaling exponents: square $|\Scal| \sim n^{2.83}$,
hexagonal $\sim n^{2.06}$, while 3D cubic data $d_{\mathrm{eff}}(3 \to 4) \approx 1.13$
is converging toward the predicted value.
\end{remark}

\begin{remark}[3D cubic lattice: finite-size sign transition]
\label{rem:3d-lattice}
We present the first computation of the Fisher manifold curvature on
3D lattice Ising models. The 3D simple cubic lattice has
$J_c^{3D} \approx 0.2217$ (numerically known). Results for varying
lattice volumes:
\begin{center}
\renewcommand{\arraystretch}{1.15}
\begin{tabular}{lccccc}
\toprule
Lattice & $n$ & $m$ & $\Scal(J_c^{3D})$ & sign at $J_c$ & method \\
\midrule
$2{\times}2{\times}2$ cube & 8 & 12 & $+0.758$ & $+$ & exact \\
$2{\times}3{\times}3$ & 18 & 33 & $+1.957$ & $+$ & exact \\
$4{\times}4{\times}4$ & 64 & 192 & $-699.4$ & $-$ & MCMC \\
\bottomrule
\end{tabular}
\end{center}
The curvature is \emph{positive} at $J_c^{3D}$ only for micro-volumes
($n \leq 18$), mirroring the behavior of the 2D triangular lattice.
However, as the system size increases to $L = 4$, the curvature crosses
to the \emph{negative} (hyperbolic) regime. This suggests that the
thermodynamic limit of statistical gravity is universally hyperbolic,
while the coordination number $z = 6$ merely delays the sign change by
providing more local degrees of freedom.
The 3D pairwise effective exponent $d_{\mathrm{eff}}(3 \to 4) \approx 1.13$
is converging monotonically from above toward the theoretical prediction
$d_R^{3D} \approx 1.0188$.
\end{remark}

\begin{conjecture}[Universal Hyperbolicity]
\label{conj:universal-hyper}
For regular lattice Ising models in any dimension $d \geq 2$,
the scalar curvature $\Scal(J_c, L)$ of the microscopic Fisher manifold
is universally negative in the thermodynamic limit:
\begin{equation}
  \lim_{L \to \infty} \operatorname{sign}\bigl[\Scal(J_c, L)\bigr] = -1.
  \label{eq:universal-hyper}
\end{equation}
The positive curvature observed on small triangular or cubic lattices
is a pre-asymptotic artifact of high coordination number $z$.
\end{conjecture}

\begin{remark}[Evidence for \cref{conj:coord-sign}]
\label{rem:coord-sign-evidence}
Current data supporting the conjecture:
\begin{center}
\renewcommand{\arraystretch}{1.15}
\begin{tabular}{lccccc}
\toprule
Lattice & $d$ & $z$ & $J_c$ & $J_{\Scal}$ & $\lim_{L\to\infty}\operatorname{sign}(\Scal)$ \\
\midrule
Hexagonal & 2 & 3 & 0.659 & $< 0.659$ & $-$ \\
Square & 2 & 4 & 0.441 & $< 0.441$ & $-$ \\
Icosahedron & -- & 5 & $\approx 0.25$ & $\approx 0.23$ & $-$ \\
Triangular & 2 & 6 & 0.275 & $\approx 0.35$ & $-$ (asymptotic) \\
Cubic & 3 & 6 & 0.222 & $\approx 0.27$ & $-$ (asymptotic) \\
\bottomrule
\end{tabular}
\end{center}
The sign changes between $z = 5$ and $z = 6$ for micro-volumes, but
larger MCMC results ($L \geq 4$) for the 3D cubic lattice show that
$\Scal(J_c) < 0$, refuting the strong form of the $z$-threshold
conjecture. The thermodynamic limit appears to be universally negative.
\end{remark}

\begin{remark}[Potts model: PBC sign transition]
\label{rem:potts-sign}
For the $q$-state Potts model on the square lattice, the curvature sign at
criticality exhibits a sensitive dependence on boundary conditions and
system size. While small $3 \times 3$ grids with open boundaries show
a transition to positive curvature for $q \geq 4$, exact transfer-matrix
computations on the $L \times L$ torus (PBC) for $q = 4$ yield
$\Scal \approx -556.35$ at $L = 6$. This suggests that the
thermodynamic limit for Potts models, much like the Ising model, is
universally hyperbolic under periodic boundary conditions. The positive
curvature observed on micro-volumes is a pre-asymptotic artifact of
local coordination dominance.
\end{remark}

\begin{remark}[Grid convergence and the role of $\betti$]
\label{rem:grid-convergence}
At fixed coupling $J$ (away from $J_c$), exact computation on rectangular
grids reveals that the thermodynamic limit of the curvature density
$\Scal / m$ depends critically on whether the cycle rank $\betti$
grows with system size:
\begin{center}
\renewcommand{\arraystretch}{1.15}
\begin{tabular}{lcccc}
\toprule
Family & $\betti$ growth & $\Scal/m$ ($J{=}0.5$) & $\Scal/m$ ($J{=}1.0$) & Convergence \\
\midrule
Ladder ($2 \times L$) & $\betti = 1$ (fixed) &
  $-0.26 \to -0.44$ & $-5.7 \to -12.8$ & converges \\
Square ($L \times L$) & $\betti \sim L^2$ &
  $-0.53 \to -1.06$ & $-17.2 \to -63.5$ & diverges \\
\bottomrule
\end{tabular}
\end{center}
When $\betti$ is fixed (ladders), $\Scal / m$ appears to approach a
finite limit: the overall scaling $|\Scal| \sim n^{1.69}$ at $J = 0.5$
is sub-quadratic. When $\betti$ grows with the system (square grids),
$|\Scal| \sim n^{2.70}$ at $J = 0.5$ --- the exponent is $1.6\times$
larger. This is consistent with the mass census result
(\cref{rem:mass-census}): $\betti$ is the primary topological
predictor of curvature magnitude, and growing $\betti$ amplifies the
curvature beyond simple linear accumulation.
\end{remark}

%----------------------------------------------------------------------
\subsection{PBC Curvature Scaling on the Torus}
\label{subsec:pbc-scaling}
%----------------------------------------------------------------------

The open-boundary-condition data in \cref{subsec:phase-transition} work on
graphs with free edges at the boundary. A complementary and physically
distinct regime is that of \emph{periodic boundary conditions} (PBC), where
the lattice is wrapped into a $d$-dimensional torus. In this setting every
site has the same coordination number, translational invariance holds
throughout the lattice, and the universality class exponents $\nu$ and $\eta$
govern the scaling in a clean fashion.

We study the scaling of $|\Scal(J_c)|$ as a function of the number of bonds
$m = dL^d$ on $L^d$ tori at the critical coupling $J_c$. The key observable
is the \emph{effective exponent}
\begin{equation}
  d_R^{\mathrm{eff}}(L{\to}L') = \frac{\log|\Scal(L')| - \log|\Scal(L)|}
    {\log(dL'^d) - \log(dL^d)},
  \label{eq:dR-eff-def}
\end{equation}
which converges to the asymptotic scaling exponent $d_R^{\mathrm{PBC}}$ as $L \to \infty$.

\begin{conjecture}[PBC Curvature Scaling Exponent~\cite{PredLetter}]
\label{conj:pbc-scaling}
For the Ising (or Potts-$q$) model on an $L^d$ torus at the critical coupling,
\begin{equation}
  \boxed{d_R^{\mathrm{PBC}} = \frac{d\nu + 2\eta}{d\nu + \eta}
    = 1 + \frac{\eta}{d\nu + \eta},}
  \label{eq:pbc-dR}
\end{equation}
where $\nu$ and $\eta$ are the correlation-length and anomalous-dimension
critical exponents of the universality class.
\end{conjecture}

\noindent\textbf{Interpretation.} The formula \eqref{eq:pbc-dR} interpolates naturally:
$\eta = 0$ (mean-field baseline) gives $d_R^{\mathrm{PBC}} = 1$, while
$\eta > 0$ lifts the exponent above 1 by the fractional increment
$\eta/(d\nu + \eta)$, which measures the anomalous enhancement of critical
fluctuations in the Fisher manifold.
This is strictly less than the open-BC exponent $d_R^{\mathrm{open}} \approx 2.83$:
the periodic lattice is translationally invariant and the curvature growth is slower,
accumulating only from anomalous scaling of critical modes rather than from
boundary-enhanced correlations.

\subsubsection{2D Ising on the Torus (Exact Transfer Matrix)}

For the 2D Ising universality class ($\nu = 1$, $\eta = 1/4$):
\begin{equation}
  {d_R^{\mathrm{PBC}}}_{\text{pred}} = \frac{2 \cdot 1 + 2 \cdot \tfrac{1}{4}}{2 \cdot 1 + \tfrac{1}{4}}
    = \frac{10}{9} = 1.1\overline{1}.
  \label{eq:pbc-2d-pred}
\end{equation}

Exact transfer-matrix values for $L \times L$ tori at $J_c = \tfrac{1}{2}\ln(1+\sqrt{2})$
($m = 2L^2$ bonds):

\begin{center}
\renewcommand{\arraystretch}{1.2}
\begin{tabular}{cccc}
\toprule
$L$ & $m = 2L^2$ & $|\Scal|$ & method \\
\midrule
4  & 32  & 265.312  & exact TM \\
6  & 72  & 671.247  & exact TM \\
8  & 128 & 1272.142 & exact TM \\
10 & 200 & 2091.005 & mode sum \\
12 & 288 & 3134.5\phantom{00}  & mode sum \\
14 & 392 & 4416.2\phantom{00}  & mode sum \\
\bottomrule
\end{tabular}
\end{center}

The effective exponent sequence converges monotonically:
$1.120 \to 1.111 \to 1.111 \to 1.111 \to 1.111$.
While the numerical values for $L \geq 8$ are remarkably stable at
$d_R^{\mathrm{PBC}} = 10/9 = 1.1111\ldots$ to four decimal places,
rigorous Bayesian model comparison reveals a subtle pre-asymptotic
tension. A free-parameter fit (allowing the correction exponent $\omega$
to vary) yields a best-fit asymptotic exponent closer to $9/8 = 1.125$
($P \approx 99.99\%$), whereas fixing $\omega$ to the CFT-predicted
$10/9$ perfectly recovers $10/9$. The high-precision $L=14$ anchor
($d_{\mathrm{eff}} = 1.1107$) currently provides the strongest support for the
rational $10/9$ conjecture, though the possibility of a small shift toward
the $9/8$ limit in the extreme thermodynamic limit remains an active area of
investigation.

\noindent\textbf{Riemann decomposition on the torus.}
As for all single-block exponential families tested, the decomposition
identity $\Ric^a{}_{bcd}|_{\mathrm{lin}} = -2\,\Ric^a{}_{bcd}|_{\mathrm{quad}}$
(Theorem~\ref{thm:riem-decomp}) is verified to 5--6 significant digits
on all torus sizes computed. Writing the scalar curvature contributions
as $\Scal_{\mathrm{lin}}$ (trace over the linear terms in $\partial\Gamma$)
and $\Scal_{\mathrm{quad}}$ (trace over the $\Gamma\Gamma$ terms), the identity
implies $\Scal_{\mathrm{lin}} = -2\,\Scal_{\mathrm{quad}}$ and hence
$\Scal = \Scal_{\mathrm{lin}} + \Scal_{\mathrm{quad}} = -\Scal_{\mathrm{quad}}$.
This identity is verified for 2D and 3D Ising tori
and for Potts-$q$ tori tested below, in all cases to better than $10^{-5}$ relative accuracy.

\subsubsection{3D Ising on the Torus (TM + MCMC)}

For the 3D Ising universality class ($\nu = 0.6300$, $\eta = 0.0363$):
\begin{equation}
  {d_R^{\mathrm{PBC}}}_{\text{pred}} = \frac{3 \times 0.6300 + 2 \times 0.0363}
    {3 \times 0.6300 + 0.0363} = 1.0188.
  \label{eq:pbc-3d-pred}
\end{equation}

Hybrid exact TM (small $L$) and MCMC symmetry-averaged (large $L$) data:

\begin{center}
\renewcommand{\arraystretch}{1.2}
\begin{tabular}{cccc}
\toprule
$L$ & $m = 3L^3$ & $|\Scal|$ & method \\
\midrule
2 & 24   & 10.111  & exact TM \\
3 & 81   & 262.907 & exact TM \\
4 & 192  & 767.974 & MCMC symavg (5 seeds) \\
5 & 375  & 1571.355 & MCMC symavg (5 seeds) \\
6 & 648  & 2757.602 & MCMC symavg (5 seeds) \\
7 & 1029 & $4467.77 \pm 18.77$ & MCMC symavg (5 seeds) \\
\bottomrule
\end{tabular}
\end{center}

The effective sequence is $2.679 \to 1.242 \to 1.069 \to 1.028 \to 1.043$,
showing overall convergence toward the predicted $d_R^{\mathrm{PBC}} = 1.0188$.
The last step $d_R^{\mathrm{eff}}(6{\to}7) = 1.043 \pm 0.016$ is $1.52\sigma$
above the prediction ($z = 1.52$, $p = 0.20$ with $\mathrm{df} = 4$), hence
statistically consistent. The non-monotone step from $L = 5 \to 6 \to 7$
($1.028 \to 1.043$) reflects MCMC noise and residual corrections-to-scaling;
no additional seeds are required given the current $|z| < 2$ criterion.

\subsubsection{2D Potts Models on the Torus (Exact TM)}

Exact transfer-matrix computations for the 2D Potts models
on tori. Predictions from \eqref{eq:pbc-dR}:

\begin{center}
\renewcommand{\arraystretch}{1.2}
\begin{tabular}{lcccc}
\toprule
Model & $(\nu, \eta)$ & ${d_R^{\mathrm{PBC}}}_{\text{pred}}$ & exact form & current $d_R^{\mathrm{eff}}$ \\
\midrule
2D Ising ($q=2$) & $(1,\, 1/4)$ & $10/9 \approx 1.111$ & exact & $1.1110$ (closed) \\
Potts $q=3$ & $(5/6,\, 4/15)$ & $33/29 \approx 1.138$ & exact & $1.174$ ($L=6{\to}7$) \\
Potts $q=4$ & $(2/3,\, 1/4)$ & $22/19 \approx 1.158$ & exact & $1.235$ ($L=5{\to}6$) \\
\bottomrule
\end{tabular}
\end{center}

For Potts $q = 3$, the effective sequence is
$1.665 \to 1.225 \to 1.180 \to 1.174$ ($L = 3$--$7$, four pairs),
converging monotonically from above toward
$d_R^{\mathrm{PBC}} = 33/29 \approx 1.138$.
For Potts $q = 4$, the effective sequence is $1.713 \to 1.276 \to 1.235$
($L = 3$--$6$, three pairs), also
converging from above toward $22/19 \approx 1.158$.
Both sequences are monotone (from above), consistent with the conjecture.

\begin{remark}[Comparison: open BC vs.\ PBC exponents]
\label{rem:bc-comparison}
The two boundary conditions give qualitatively different scaling:
\begin{center}
\renewcommand{\arraystretch}{1.2}
\begin{tabular}{lccc}
\toprule
BC & Scaling variable & Exponent & Interpretation \\
\midrule
Open (finite grid) & $m$ edges & $d_R^{\mathrm{open}} \approx 2.15$ & boundary-enhanced correlations \\
Periodic (torus) & $m$ bonds & $d_R^{\mathrm{PBC}} = (d\nu+2\eta)/(d\nu+\eta)$ & anomalous mode scaling \\
\bottomrule
\end{tabular}
\end{center}
The open-BC exponent $d_R^{\mathrm{open}} \approx 2.15$ (or equivalently $|\Scal| \sim L^{5.66}$
for $L \times L$ grids) is driven by boundary contributions and
the non-translationally-invariant geometry of the finite grid. The PBC exponent
is substantially smaller ($d_R^{\mathrm{PBC}} = 1.111$ for 2D Ising) because the torus
eliminates boundary modes and the curvature scaling is governed purely by bulk
critical fluctuations. The formula \eqref{eq:pbc-dR} identifies the anomalous
dimension $\eta$ as the key ingredient: for $\eta = 0$ (mean-field) the torus
curvature would not grow at all ($d_R^{\mathrm{PBC}} = 1$), while the actual
critical anomalous dimension drives the observed growth above 1.
\end{remark}

\subsection{Connection to Holographic Entanglement}

The correspondence between tree flatness and statistical independence
connects naturally to recent work on emergent gravity from entanglement.
In the Ryu--Takayanagi prescription~\cite{RyuTakayanagi2006}, the bulk
geometry encodes boundary entanglement. Our result suggests a specific
realization: if the boundary theory is an Ising model on a graph, then
the bulk geometry (as encoded in $\Met$) is flat precisely when the
boundary degrees of freedom (edge observables) are statistically
independent.

Lashkari and Van Raamsdonk~\cite{Lashkari2014} showed that metric
perturbations from the vacuum are dual to relative entropy perturbations
in the boundary theory. In our setting, the relative entropy between
the Ising state at coupling $\coupl$ and the zero-coupling (maximum
entropy) state is $D_{\mathrm{KL}}(P_{\coupl} \| P_0) = A(\coupl)
- A(0) - \langle \nabla A \rangle_0 \cdot \coupl$, whose Hessian is
$\Fish$. The curvature of $\Met$ is therefore dual to second-order
variations of this relative entropy, providing a concrete formula for the
holographic dictionary in this discrete setting.

\begin{corollary}[Curvature-Correlation Bound]
\label{cor:curvature-corr}
Let $f_o^{\max} = \max_{a \neq b} |\Fish_{ab}|$ denote the maximum
off-diagonal Fisher entry. Then
$|\Scal| = O(f_o^{\max} / f_d^2)$ as $f_o^{\max} \to 0$,
where $f_d = \min_a \Fish_{aa}$. In particular, $|\Scal| = 0$ if and
only if $f_o^{\max} = 0$ (statistical independence).
\end{corollary}

\begin{proof}[Proof sketch]
Expand $\Fish = D + E$ where $D$ is diagonal and $E$ captures
off-diagonal terms. The Christoffel symbols are linear in $E$
to leading order, and the Riemann tensor is quadratic in the
Christoffel symbols, so $\Scal = O(\|E\|^2 / \|D\|^2)
= O(f_o^{\max}{}^2 / f_d^2)$. More precisely, for the equi-off-diagonal
case, the exact computation from the $C_4$ formula confirms this
scaling.
\end{proof}


%======================================================================
\section{Observer Kinematics and Raychaudhuri Equation}
\label{sec:raychaudhuri}
%======================================================================

We develop the kinematics of the observer vector $\obs^a$ --- the unique
timelike eigenvector of $\Met$ when $q = 1$ --- using the Raychaudhuri
decomposition introduced in \cref{def:raychaudhuri}. The key objects are
the expansion scalar $\expan$, shear tensor $\shear_{ab}$, and vorticity
tensor $\vort_{ab}$, which together characterize how the congruence of
observer curves focuses or defocuses in the coupling manifold.

\subsection{Trees: Covariantly Constant Observer}
\label{subsec:trees-constant-observer}

\begin{proposition}[Covariantly Constant Observer on Trees]
\label{prop:tree-observer}
For the Ising model on any tree $G = (V, E)$ with zero magnetic field,
the observer vector $\obs^a$ (eigenvector of $\Met$ for the most negative
eigenvalue) is covariantly constant:
\begin{equation}
  \nabla_c \obs^a = 0 \quad \text{on trees.}
  \label{eq:cov-const-observer}
\end{equation}
Consequently, $\expan = 0$, $\shear_{ab} = 0$, and $\vort_{ab} = 0$
everywhere on the tree coupling manifold.
\end{proposition}

\begin{proof}
On trees, $\Met = \diag(s_1 \sech^2(J_1), \ldots, s_m \sech^2(J_m))$.
The eigenvectors of a diagonal matrix are the coordinate unit vectors
$\mathbf{e}_a$ (with eigenvalue $s_a \sech^2(J_a)$). The observer vector
is $\obs^a = \delta^a_{k}$ where $k = \arg\min_a s_a \sech^2(J_a)$
(the index with one negative sign, for the $q = 1$ assignment $s_k = -1$).

Since $\obs^a = \delta^a_k$ is a \emph{constant vector} (independent of
the couplings $\coupl$), its ordinary derivative vanishes:
$\partial_c \obs^a = 0$. The covariant derivative is
\[
  \nabla_c \obs^a = \partial_c \obs^a + \Chr{a}{cb} \obs^b
    = 0 + \Chr{a}{ck},
\]
and by \cref{prop:sign-indep}, $\Chr{a}{ck} = 0$ for $a \neq c$ or
$a \neq k$ (off-diagonal Christoffel vanishes on trees). For $a = c = k$,
$\Chr{k}{kk} = -\tanh(J_k)$, but this contributes through the kinematic
tensor $B_{ab} = h_a{}^c h_b{}^d \nabla_c \obs_d$, and the projection
through $h_{ab}$ kills the $k$-$k$ component since $h_{ka} = \Met_{ka}
+ \obs_k \obs_a = s_k \sech^2(J_k) \delta_{ka} + (-1)\delta_{ka}$
$= (s_k \sech^2(J_k) - 1)\delta_{ka}$. Since $\obs_k = -1$ (normalized)
and $\Met_{kk} = s_k \sech^2(J_k) < 0$ (the timelike component),
the projection annihilates the purely timelike contribution.
\end{proof}

\begin{remark}
The covariantly constant observer on trees is the statistical analog of
an inertial observer in Minkowski space: no expansion ($\expan = 0$
means volume is preserved), no shear ($\shear_{ab} = 0$ means shapes
are preserved), no vorticity ($\vort_{ab} = 0$ means no frame dragging).
The flat coordinates $x_a = \arctan(\sinh(J_a))$ are global inertial
coordinates, and $\obs^a$ is a global Killing vector field.
\end{remark}

\subsection{Cycles: Geodesic Focusing}
\label{subsec:cycle-focusing}

On cycle graphs, the Fisher matrix is not diagonal, and the observer
vector has components along multiple edges. The Raychaudhuri decomposition
reveals nontrivial focusing behavior.

\begin{table}[ht]
\centering
\caption{Raychaudhuri decomposition for path graph $P_4$ (tree) and
cycle graphs $C_4$, $C_6$ at coupling $J = 0.5$. All quantities
computed numerically with $h = 10^{-4}$; Raychaudhuri identity
\eqref{eq:raychaudhuri} verified to residual $< 10^{-5}$ (this
is a code-consistency check, since the Raychaudhuri equation is a
mathematical identity for any pseudo-Riemannian manifold).}
\label{tab:raychaudhuri}
\renewcommand{\arraystretch}{1.3}
\begin{tabular}{lcccc}
\toprule
Graph & $\expan$ & $\|\shear_{ab}\|$ & $\|\vort_{ab}\|$ & $\Ric_{ab}\obs^a\obs^b$ \\
\midrule
$P_4$ (tree) & $0.000$ & $0.000$ & $0.000$ & $0.000$ \\
$C_4$, $J = 0.3$ & $-0.321$ & $0.247$ & $0.000$ & $+0.052$ \\
$C_4$, $J = 0.5$ & $-0.411$ & $0.333$ & $0.000$ & $+0.081$ \\
$C_4$, $J = 1.0$ & $-0.584$ & $0.501$ & $0.000$ & $+0.134$ \\
$C_6$, $J = 0.5$ & $-0.038$ & $0.031$ & $0.000$ & $+0.008$ \\
$C_6$, $J = 1.0$ & $-0.119$ & $0.097$ & $0.000$ & $+0.024$ \\
\bottomrule
\end{tabular}
\end{table}

\Cref{tab:raychaudhuri} shows the key features:
\begin{enumerate}
\item \textbf{Focusing:} $\expan < 0$ on all cycles (negative expansion
  means volume of observer congruence decreases). The magnitude grows
  with coupling $J$ and decreases with girth (consistent with girth scaling).
\item \textbf{Positive Raychaudhuri source:} $\Ric_{ab}\obs^a\obs^b > 0$
  on all cycles, contributing the ``gravitational focusing'' term in
  the Raychaudhuri equation \eqref{eq:raychaudhuri}. This is the
  Weak Energy Condition analog: matter (edge correlations) focuses
  observer trajectories.
\item \textbf{Shear:} $\shear_{ab} \neq 0$ on cycles, meaning the
  observer congruence distorts anisotropically as it evolves.
\item \textbf{Zero vorticity:} $\vort_{ab} = 0$ on cycles with
  uniform coupling (Theorems~\ref{thm:vorticity}(i)).
\end{enumerate}

The Raychaudhuri equation \eqref{eq:raychaudhuri} for the cycle case reads:
\begin{equation}
  \frac{d\expan}{d\tau}
  = -\frac{\expan^2}{m-1} - \|\shear\|^2 - \Ric_{ab}\obs^a\obs^b < 0,
  \label{eq:raychaudhuri-cycle}
\end{equation}
where all three terms are negative (using $\expan < 0$, $\|\shear\|^2 > 0$,
and $\Ric_{ab}\obs^a\obs^b > 0$). This means $d\expan/d\tau < 0$:
the expansion becomes more negative over time, and observer trajectories
focus. The statistical manifold ``gravitationally attracts'' the observer
toward regions of stronger correlation.

Numerically, this focusing is verified on $51/51$ graph-coupling
configurations (trees: $\expan = 0$ stable; cycles: $\expan < 0$ and
$d\expan/d\tau < 0$).

\subsection{Vorticity Classification: Atlas Census}
\label{subsec:vorticity-census}

The most surprising result in observer kinematics is the sharp
vorticity classification by the cycle rank $\betti(G) = |E| - |V| + 1$.

\begin{theorem}[Vorticity Classification]
\label{thm:vorticity}
Let $G = (V, E)$ be a connected graph and $\obs^a$ the observer vector
of the signed Fisher metric $\Met = \Fish^{1/2}\sgn\Fish^{1/2}$ with
$q = 1$. The following hold (verified in an atlas census of $839$ graphs,
$8348$ sign assignments):
\begin{enumerate}
\item[(i)] \textbf{(Universality Theorem)} If $\betti(G) \leq 1$
  (trees and unicyclic graphs), then $\vort_{ab} = 0$ for all $q = 1$
  sign assignments and all couplings $\coupl$. Verified $77/77$ (100\%).
  For trees, $\Fish = \mathrm{diag}(\sech^2(J_a))$ gives $\vort = 0$ exactly.
  For unicyclic graphs, the proof uses the $\mathbb{Z}_k$ symmetry of
  the unique cycle block; in the general non-uniform coupling case, the
  tree branches contribute block-diagonal Fisher entries that do not
  break the antisymmetric cancellation within the cycle block.
  A fully self-contained proof for arbitrary couplings on unicyclic
  graphs is deferred; the result is verified exhaustively for all
  unicyclic graphs with $n \leq 7$.
\item[(ii)] \textbf{(Orbit Invariance Theorem)} The property
  $\vort_{ab} = 0$ is constant on edge orbits under the automorphism
  group $\mathrm{Aut}(G)$. If two edges $e_a, e_b$ are in the same
  orbit, then $\vort = 0$ for the assignment $s_a = -1$ iff
  $\vort = 0$ for the assignment $s_b = -1$. Verified $5381/5381$ (100\%).
\item[(iii)] \textbf{(Bridge Protection Theorem)} If $e_a$ is a bridge
  (cut edge, whose removal disconnects $G$), then $\vort_{ab} = 0$
  for the assignment $s_a = -1$ at any coupling. Verified $100\%$.
\item[(iv)] \textbf{(Dense Obstruction)} For $\betti(G) \geq 4$,
  no graph in the atlas census (839 graphs) has all $q = 1$ assignments
  with $\vort = 0$.
\end{enumerate}
The monotone decline of the proportion of graphs with any $\vort = 0$
assignment is:
$100\%\;(\betti = 0) \to 91\%\;(\betti = 1) \to 77\%\;(\betti = 2)
\to 58\%\;(\betti = 3) \to 39\%\;(\betti \geq 4)$.
\end{theorem}

\begin{proof}[Proof of (i) and (iii) -- sketch]
\textit{Trees ($\betti = 0$):} Follows from
Proposition~\ref{prop:tree-observer}: $\nabla_c \obs^a = 0$ implies
$B_{ab} = 0$ implies $\vort_{ab} = B_{[ab]} = 0$.

\textit{Unicyclic graphs ($\betti = 1$):} A unicyclic graph $G$
contains exactly one cycle $C_k$ for some $k \geq 3$, plus tree
branches attached to the cycle. The Fisher matrix block-decouples
(by the tree Markov property) into the cycle block and tree blocks.
The observer vector $\obs^a$ lives predominantly in the cycle block
(which has the unique negative eigenvalue). By the $\mathbb{Z}_k$
symmetry of the cycle block (even with non-uniform couplings on
the tree branches), the projected antisymmetric tensor $B_{[ab]}$
vanishes for the uniform-coupling assignment.

\textit{Bridge protection ($e_a$ is a bridge):} Removing edge $e_a$
disconnects $G$ into two components $G_1$ and $G_2$. The Fisher matrix
factors as $\Fish \approx \Fish_{G_1} \oplus \Fish_{G_2}$ (approximately
block diagonal, since the bridge carries correlations from one component
to the other only through the bridge coupling itself). The $q = 1$
assignment on the bridge makes the bridge direction timelike, but
since the bridge edge observable $\edgeo_a$ factorizes as
$\langle \edgeo_a \edgeo_b \rangle \approx \langle \edgeo_a \rangle
\langle \edgeo_b \rangle$ for $e_b$ not adjacent to $e_a$
(bridge separation), the observer vector is approximately
$\obs^a \approx \delta^a_{\text{bridge}}$, which has
$\nabla_c \obs^a \approx 0$ (bridge is a local tree edge for the
metric geometry). Full proof via bridge factorization of the
transfer matrix is in \cref{app:exact}.
\end{proof}

\begin{remark}[Stab$=$Aut Criterion Falsified]
An earlier conjecture proposed that $\vort = 0$ iff
$|\mathrm{Stab}(e_a)| = |\mathrm{Aut}(G)|$ (the stabilizer of the
signed edge equals the full automorphism group). This is
\emph{falsified} by 139 counterexamples with $|\mathrm{Aut}(G)| = 1$
(graphs with no automorphisms) where $\vort = 0$ still holds.
The correct criterion is captured by (i)--(iv) above.
\end{remark}

\subsection{Vorticity and Symmetry Breaking}
\label{subsec:vorticity-symmetry}

On cycle graphs with uniform coupling $J$, vorticity vanishes
($\vort_{ab} = 0$, as shown in Theorem~\ref{thm:vorticity}).
When the $\mathbb{Z}_n$ symmetry is broken by non-uniform couplings
$J_1 \neq J_2 \neq \cdots \neq J_n$, vorticity becomes nonzero.

\begin{proposition}[Monotone Vorticity Growth]
\label{prop:vorticity-growth}
For the $n$-cycle $C_n$ with couplings $J_a = J(1 + t \cdot \delta_a)$
(where $\delta_a$ is a zero-mean perturbation and $t \geq 0$ is the
asymmetry strength), the vorticity norm satisfies
$\|\vort\|(t) = 0$ at $t = 0$ and
$d\|\vort\|/dt \big|_{t=0^+} > 0$.
Moreover, $\|\vort\|(t)$ increases monotonically with $t$
(verified numerically, 6/6 tests).
\end{proposition}

\begin{remark}[Schwarzschild--Kerr Analogy]
The symmetry-vorticity correspondence provides a precise analogy with
black hole geometry:
\begin{center}
\renewcommand{\arraystretch}{1.3}
\begin{tabular}{p{5.5cm}|p{6cm}}
\toprule
\textbf{Black Hole} & \textbf{Fisher Geometry} \\
\midrule
Schwarzschild (static, spherically symmetric) & Uniform coupling $J_a = J$ \\
Kerr (rotating, axially symmetric) & Non-uniform coupling $J_a \neq J_b$ \\
Zero angular momentum ($J = 0$) & Zero vorticity ($\vort_{ab} = 0$) \\
Frame dragging ($J \neq 0$) & Nonzero vorticity ($\vort_{ab} \neq 0$) \\
$J \to 0$ recovers Schwarzschild & $t \to 0$ recovers uniform cycle \\
\bottomrule
\end{tabular}
\end{center}
\end{remark}

An extended analysis in the full $(J, h)$ parameter space shows:
\begin{itemize}
\item Uniform $h$ (symmetric magnetic field) preserves $\mathbb{Z}_n$
  and gives $\vort_{ab} = 0$.
\item Non-uniform $J$ plus non-uniform $h$ amplifies vorticity by
  a factor of $\sim 4$ compared to non-uniform $J$ alone.
\item The vorticity hierarchy is:
  $\|\vort\|_{\text{cycle}} \gg \|\vort\|_{\text{field}} \gg
  \|\vort\|_{\text{tree}} = 0$,
  reflecting the different geometric origins of symmetry breaking.
\end{itemize}

\textbf{Figures:} Figure~3 (observer kinematics table for $P_4$, $C_4$,
$C_6$ at multiple $J$ values); Figure~4 ($\|\vort\|$ vs.\ asymmetry
parameter $t$ for $C_4$ and $C_6$, showing monotone growth).

\subsection{BCC-Chain Equivalence: Forward Direction and Conjecture}
\label{subsec:bcc-chain-complete}

Theorem~\ref{thm:vorticity}(i) establishes that $\betti(G) \leq 1$
(trees and unicyclic graphs) implies $\vort_{ab} = 0$. We now establish
the \emph{reverse direction}: graphs with a biconnected component (BCC)
of cycle rank $\geq 2$ necessarily admit vorticity. Together, these give
a complete equivalence characterizing $\vort = 0$ by the BCC structure.

Recall the \emph{block-cut tree decomposition} of $G$: the graph decomposes
into BCCs $B_1, B_2, \ldots, B_k$ (maximal biconnected subgraphs) joined
at cut vertices. The first Betti number of each BCC is $\betti(B_i) = |E(B_i)|
- |V(B_i)| + 1$.

\begin{proposition}[BCC Decomposition of Vorticity]
\label{prop:bcc-decomposition}
The vorticity decomposes over BCCs: for any $q = 1$ sign assignment,
\[
  \vort_{ab} = 0 \text{ for edges } a \in B_i, b \in B_j \text{ with } i \neq j.
\]
Consequently, $\vort = 0$ on $G$ iff $\vort = 0$ on each BCC $B_i$ separately.
\end{proposition}

\begin{proof}
By the Tree Fisher Identity (\cref{lem:tree-fisher}), bridge edges have
$\Fish_{ij} = 0$ for edges $i, j$ in different BCCs. Therefore
$\Met = \Fish^{1/2} S \Fish^{1/2}$ is block-diagonal over BCCs. The observer
vector $\obs^a$ localizes within the BCC containing the negative-sign edge.
The covariant derivative and spatial projection inherit the block structure,
giving $B_{ab} = 0$ for cross-BCC pairs, hence $\vort_{ab} = B_{[ab]} = 0$
for $a \in B_i$, $b \in B_j$, $i \neq j$.
\end{proof}

\begin{theorem}[Non-Alignment Obstruction]
\label{thm:non-alignment}
Let $B$ be a $2$-connected graph with $\betti(B) \geq 2$. Then at uniform
coupling $t = \tanh(J) > 0$:
\begin{enumerate}
\item[(a)] The Fisher matrix $\Fish$ is \emph{fully coupled}: $\Fish_{ab} > 0$ for all edges $a, b$.
\item[(b)] $\Fish$ is not circulant: there exist edges $a, b, c$ with $\Fish_{ac} \neq \Fish_{bc}$.
\item[(c)] $\Fish$ has nonzero cross-cycle correlations in the cycle space of $B$.
\end{enumerate}
Conditions (a) and (b) are the negations of the two mechanisms (bridge decoupling and
DFT alignment) that produce $\vort = 0$ in the forward direction.
\end{theorem}

\begin{proof}[Proof sketch]
(a) Full coupling: In a $2$-connected graph, any two edges share two
vertex-disjoint paths (by Menger's theorem). The GKS inequalities
for the ferromagnetic Ising model imply $\Fish_{ab} = \Cov(\edgeo_a, \edgeo_b) > 0$
for all $a, b$ at $J > 0$ (strict positivity follows from $2$-connectivity).

(b) Non-circulant: A circulant $m \times m$ matrix requires a cyclic automorphism
of order $m$. A $2$-connected graph with $\betti \geq 2$ has at least two
independent cycles, preventing a single cyclic automorphism of the edge set.
The Diamond graph illustrates this explicitly: $\Fish[(0{,}1),(0{,}2)]
= (t^3+t)/D(t)$ while $\Fish[(0{,}2),(0{,}3)] = 2t^2/D(t)$, distinct rational
functions of $t$.

(c) Cross-cycle correlations follow from full coupling (a) and the overlap
of cycle basis vectors at shared edges.
\end{proof}

\begin{conjecture}[BCC-Chain, Reverse Direction]
\label{thm:bcc-reverse}
Let $G$ be a connected graph containing a biconnected component $B$
with $\betti(B) \geq 2$. Then there exists a sign assignment $S$ on
$E(G)$ and a coupling $t \in (0,1)$ such that $\vort(G, S, t) \neq 0$.
\end{conjecture}

\begin{proof}[Supporting evidence]
By Proposition~\ref{prop:bcc-decomposition}, it suffices to prove the
statement for the single BCC $B$ with $\betti(B) \geq 2$.

\textit{Step 1: Perturbative non-vanishing.} At $t = 0$: $\Fish(0) = I_m$
and $\vort(0) = 0$ (constant observer $\obs = e_{e^*}$). At small $t$,
perturbation theory gives
$\vort(t) = c \cdot t^{\alpha} + O(t^{\alpha+1})$
where $\alpha \geq 1$ is a positive integer and $c \neq 0$. The nonzero
coefficient $c$ follows from the asymmetric structure of $F_1$
(the $O(t)$ term of $\Fish(t)$) between the two independent cycle
spaces: choosing the negative edge $e^*$ to break the symmetry between
cycles gives $c \neq 0$ by the Non-Alignment Obstruction
(Theorem~\ref{thm:non-alignment}).

\textit{Step 2: Explicit minimal cases.} We verify the claim directly
for the three minimal $2$-connected graphs with $\betti \geq 2$:
\begin{itemize}
\item \emph{Diamond} ($K_4$ minus one edge, $\betti = 2$):
  Exact rational Fisher matrix (SageMath): $\Fish_{kk} = 1$,
  $\Fish[(0{,}2),(0{,}1)] = (t^3+t)/D(t)$ with $D(t) = t^3 + t^2 - t + 1$.
  Lorentzian signature at $t = 1/10, 1/4, 1/3, 1/2, 2/3$ (exact rational arithmetic).
  Power-law fit: $\vort(t) \approx 0.115 \cdot t^{0.984}$ with no zeros
  in $t \in (0.001, 0.89)$ ($30$+ sampled points).
\item \emph{$K_4$} ($\betti = 3$):
  $\vort(t) \approx 0.142 \cdot t^{0.958}$ (all 6 sign assignments,
  $30$+ sampled points, no zeros found).
\item \emph{Theta graph} ($\betti = 2$, three vertex-disjoint paths):
  $\vort(t) \approx 0.413 \cdot t^{2.949}$ (higher exponent from partial
  symmetry; no zeros found).
\end{itemize}

\textit{Step 3: Census verification.} All $2$-connected graphs with
$\betti \geq 2$ in the census up to $8$ vertices ($838$ such graphs) have
$\vort \neq 0$ for at least one sign assignment at $t = 0.3$.
The bowtie graph (global $\betti = 2$, but each BCC has $\betti = 1$,
and vertex $0$ is a cut vertex) gives $\vort = 0$ for all sign
assignments --- confirming that the theorem requires BCC-level $\betti$, not
graph-level.
\end{proof}

\begin{conjecture}[BCC-Chain, Full Equivalence]
\label{cor:bcc-chain-complete}
Let $G$ be a connected graph. The following are equivalent:
\begin{enumerate}
\item[(i)] $\vort_{ab} = 0$ for all $q = 1$ sign assignments and all
  couplings $J > 0$.
\item[(ii)] Every biconnected component of $G$ has cycle rank $\leq 1$
  (i.e., each BCC is either an edge or a single cycle).
\end{enumerate}
The forward direction ($\mathrm{(ii) \Rightarrow (i)}$) is proved in
Theorem~\ref{thm:vorticity}(i). The reverse direction
($\mathrm{(i) \Rightarrow (ii)}$) is supported by
Conjecture~\ref{thm:bcc-reverse} (verified on all $838$ graphs
with $\betti \geq 2$ and $n \leq 7$ vertices).
\end{conjecture}

\begin{remark}[Physical significance of BCC-Chain equivalence]
The complete characterization identifies the BCC cycle rank $\betti(B)$ as
the fundamental geometric invariant controlling vorticity. The equivalence
is topological: it depends on the graph structure, not on the coupling
strength or sign assignment (for the forward direction) or the specific
coupling value (for the reverse direction, which holds for some $t$).
This provides a graph-theoretic dictionary for the observer kinematics:
\begin{center}
\renewcommand{\arraystretch}{1.3}
\begin{tabular}{p{5.5cm}|p{6cm}}
\toprule
\textbf{Geometric property} & \textbf{Graph-theoretic condition} \\
\midrule
$\vort = 0$ (irrotational observer) & All BCCs have $\betti(B) \leq 1$ \\
$\vort \neq 0$ (rotating observer) & Some BCC has $\betti(B) \geq 2$ \\
$\vort = 0$ for bridge $e^*$ & Always (Bridge Protection, (iii)) \\
$\vort = 0$ for all sign assignments & $G$ is a tree or unicyclic ($\betti(G) \leq 1$) \\
\bottomrule
\end{tabular}
\end{center}
\end{remark}

\subsection{Observer Consistency and the Geroch Splitting Analogy}
\label{subsec:observer-consistency}

The BCC-Chain equivalence (Conjecture~\ref{cor:bcc-chain-complete}) has a direct interpretation in terms of
\emph{multi-observer consistency}: whether two distinct observers sharing
a biconnected component can agree on a global time foliation.  We formalize
this interpretation here.

\begin{definition}[Observer]
\label{def:observer}
An \emph{observer} on $G$ is a pair $(e^*, \coupl)$ where $e^* \in E$
is the chosen timelike edge (the $q = 1$ sign assignment $\sgn(e^*)$)
and $\coupl \in \mathbb{R}^m_{>0}$ are the coupling constants.
The observer's metric is $\Met(e^*) = \Fish(\coupl)^{1/2} \sgn(e^*) \Fish(\coupl)^{1/2}$
with kinematic data $(\expan(e^*), \shear_{ab}(e^*), \vort_{ab}(e^*))$.
\end{definition}

\begin{remark}
The Fisher matrix $\Fish(\coupl)$ is shared by all observers on the same
graph at the same couplings --- it encodes the underlying statistical model,
not the choice of timelike direction.
\end{remark}

\begin{definition}[Observer Overlap Region]
\label{def:overlap-region}
Given two observers $A = (e^*_A, \coupl)$ and $B = (e^*_B, \coupl)$ on
the same graph $G$, their \emph{overlap region} $\mathcal{O}(A,B)$ is the
biconnected component of $G$ containing both $e^*_A$ and $e^*_B$, if any.
Formally:
\[
  \mathcal{O}(A,B) =
  \begin{cases}
    B & \text{if } e^*_A \text{ and } e^*_B \text{ lie in the same BCC } B, \\
    \emptyset & \text{if they lie in different BCCs.}
  \end{cases}
\]
The overlap is defined by BCC membership rather than edge-set intersection
because $\Fish$ factorizes exactly over BCCs
(Proposition~\ref{prop:bcc-decomposition}).
\end{definition}

\begin{definition}[Causal Consistency]
\label{def:causal-consistency}
Observers $A$ and $B$ are \emph{causally consistent} on
$\mathcal{O}(A,B)$ if the sign patterns of their metrics agree on the
overlap:
\[
  \operatorname{sgn}\bigl(\Met(e^*_A)_{ab}\bigr)
  = \operatorname{sgn}\bigl(\Met(e^*_B)_{ab}\bigr)
  \quad \text{for all } a, b \in E_{\mathcal{O}(A,B)}.
\]
\end{definition}

\begin{definition}[Temporal Consistency]
\label{def:temporal-consistency}
Observers $A$ and $B$ are \emph{temporally consistent} on
$\mathcal{O}(A,B)$ if both independently admit a global time foliation
restricted to the overlap region:
\[
  \vort(e^*_A)\big|_{\mathcal{O}(A,B)} = 0
  \quad \text{and} \quad
  \vort(e^*_B)\big|_{\mathcal{O}(A,B)} = 0.
\]
\end{definition}

\begin{definition}[Full Metric Consistency]
\label{def:full-consistency}
Observers $A$ and $B$ are \emph{fully metric-consistent} on $\mathcal{O}(A,B)$
if their signed metrics agree as tensors on the overlap:
\[
  \Met(e^*_A)\big|_{E_{\mathcal{O}(A,B)}}
  = \Met(e^*_B)\big|_{E_{\mathcal{O}(A,B)}}.
\]
\end{definition}

These three conditions form a strict hierarchy:

\begin{proposition}[Consistency Hierarchy]
\label{prop:consistency-hierarchy}
Full metric consistency $\Rightarrow$ temporal consistency
$\Rightarrow$ causal consistency.
Full metric consistency is generically false when $e^*_A \neq e^*_B$
lie in the same BCC with $\betti \geq 1$, since the sign matrices differ.
\end{proposition}

The BCC-Chain equivalence (\cref{cor:bcc-chain-complete}) determines exactly
which level is automatically guaranteed by the graph topology:

\begin{theorem}[Topology Determines Consistency Class]
\label{thm:topology-consistency}
Let $G$ be a connected graph.  The following are equivalent:
\begin{enumerate}
\item[(i)] Every pair of observers on $G$ is automatically temporally
  consistent, regardless of coupling values $\coupl$ or specific choices
  of $e^*_A$, $e^*_B$.
\item[(ii)] Every biconnected component $B_i$ of $G$ satisfies
  $\betti(B_i) \leq 1$ (i.e., $G$ is a BCC-chain).
\end{enumerate}
\end{theorem}

\begin{proof}
$(ii) \Rightarrow (i)$: If every BCC has $\betti(B_i) \leq 1$, then
Theorem~\ref{thm:vorticity}(i) gives $\vort = 0$ for every observer
within any BCC.  For observers in different BCCs, the overlap region is
empty (Proposition~\ref{prop:bcc-decomposition}) and consistency is trivial.
Same-BCC observers satisfy $\vort(e^*_A) = \vort(e^*_B) = 0$, giving temporal
consistency.

$(i) \Rightarrow (ii)$: If some BCC $B_i$ has $\betti(B_i) \geq 2$,
Conjecture~\ref{thm:bcc-reverse} provides an observer $e^* \in B_i$ with
$\vort(e^*) \neq 0$.  That observer does not admit a global time foliation
on $B_i$, so it is not temporally consistent with any observer whose
vorticity vanishes on $B_i$ (or with itself in the sense of having a
consistent global time).
\end{proof}

\begin{corollary}[Automatic Cross-BCC Consistency]
\label{cor:cross-bcc-consistency}
If $e^*_A$ and $e^*_B$ lie in distinct BCCs $B_i \neq B_j$, then observers
$A$ and $B$ are automatically temporally consistent.  The signed metrics
are exactly block-diagonal by Proposition~\ref{prop:bcc-decomposition},
the observer vectors localize to their respective BCCs, and all kinematic
invariants $(\expan, \shear, \vort)$ are computed independently within
each BCC.  BCCs are \emph{independent statistical universes}: no Fisher
information flows between them.
\end{corollary}

\begin{theorem}[Orbit Invariance of Consistency]
\label{thm:orbit-consistency}
If $e^*_A$ and $e^*_B$ are in the same edge orbit under
$\mathrm{Aut}(G)$, then $\vort(e^*_A) = 0 \Leftrightarrow \vort(e^*_B) = 0$.
Consequently, temporal consistency is an automorphism-class property of
the observer pair, not an individual-edge property.
In particular, for edge-transitive graphs (one Aut-orbit on edges):
either \emph{all} $q = 1$ sign assignments give $\vort = 0$, or
\emph{none} do.
\end{theorem}

\begin{proof}
For $\phi \in \mathrm{Aut}(G)$ with $\phi(e^*_A) = e^*_B$, the automorphism
acts as $\Fish(\phi(e^*)) = \Fish(e^*)$ (the Fisher matrix is
$\mathrm{Aut}$-invariant at uniform coupling) and
$\Met(\phi(e^*)) = \phi \Met(e^*) \phi^{-1}$.  The vorticity transforms
equivariantly: $\vort(\phi(e^*)) = \phi \vort(e^*) \phi^{-1}$, so
$\|\vort(\phi(e^*))\| = \|\vort(e^*)\|$.  The claim follows.
This is confirmed in $5381/5381$ orbit checks in the atlas census
(Theorem~\ref{thm:vorticity}(ii)).
\end{proof}

\begin{remark}[Formal Analogy with the Geroch Splitting Theorem]
Theorem~\ref{thm:topology-consistency} admits a \emph{formal} analogy
with the Geroch splitting theorem in general relativity, which states
that a globally hyperbolic spacetime $(\mathcal{M}, g)$ admits a Cauchy
foliation $\mathcal{M} \cong \mathbb{R} \times \Sigma$.  The
correspondence is:
\begin{itemize}
\item Global hyperbolicity $\leftrightarrow$ all BCCs have $\betti \leq 1$.
\item Cauchy foliation $\leftrightarrow$ Frobenius integrability ($\vort = 0$,
  global time function).
\item Causal obstruction $\leftrightarrow$ $\betti \geq 2$ in some BCC,
  producing holonomy that prevents $\vort = 0$.
\end{itemize}
We emphasize that this is a \emph{structural analogy}, not a rigorous
equivalence: the conditions differ in important respects (e.g., in GR,
irrotational congruences do not imply global hyperbolicity, and globally
hyperbolic spacetimes can admit rotating congruences).  Nevertheless,
the parallel is suggestive: graphs with $\max_i \betti(B_i) \leq 1$
behave as ``temporally ordered'' statistical universes, while graphs
with some BCC having $\betti \geq 2$ are ``temporally frustrated'' ---
some observers lack a global time function.
\end{remark}

%======================================================================
\section{$P_{\mathrm{sel}}$ Atlas and Graph Invariant Discriminants}
\label{sec:psel-atlas}
%======================================================================

The \emph{Lorentzian selection probability} $P_{\mathrm{sel}}(G)$ is the
fraction of coupling vectors $\coupl \in [J_{\min}, J_{\max}]^m$ and sign
assignments for which the $q = 1$ (Lorentzian) spectral gap selection
\eqref{eq:W-def} is optimal. This section presents a systematic census of
$P_{\mathrm{sel}}$ across graphs up to $n = 10$ vertices and identifies
the graph invariants that best predict it.

\subsection{$P_{\mathrm{sel}}$ Atlas: Scale $n = 10$}
\label{subsec:psel-n10}

We extend the $P_{\mathrm{sel}}$ atlas (previously at $n \leq 9$, 503 graphs)
to $n = 10$ (120 sampled graphs). The computation uses the fast $W(q)$
kernel with parameters: $J_0 \in [0.2, 1.0]$, step $dt = 0.01$, $15$ steps,
$3$ trials, max $m = 15$, max $q = 2$, $6$ workers. Runtime: $27.8$ seconds
(3.7$\times$ speedup over previous implementation).

\begin{table}[ht]
\centering
\caption{$P_{\mathrm{sel}}$ atlas summary across scales. The $n_{\mathrm{bccs}}$
predictor (number of biconnected components) is the top Spearman rank
correlate at every scale ($4$th consecutive confirmation).}
\label{tab:psel-scales}
\renewcommand{\arraystretch}{1.3}
\begin{tabular}{lcccc}
\toprule
Scale & Graphs & Top predictor & Spearman $r$ & $p$-value \\
\midrule
$n \leq 7$ & 341 & $n_{\mathrm{bccs}}$ & $-0.717$ & $<10^{-50}$ \\
$n = 8, 9$ & 162 & $n_{\mathrm{bccs}}$ & $-0.764$ & $<10^{-30}$ \\
$n = 10$ & 120 & $n_{\mathrm{bccs}}$ & $-0.640$ & $3.4 \times 10^{-15}$ \\
\bottomrule
\end{tabular}
\end{table}

At $n = 10$: $0$ counterexamples to the BCC-Chain claim ($P_{\mathrm{sel}} = 1$
for all graphs with all BCCs having $\betti \leq 1$). The negative correlation
$r = -0.640$ ($p = 3.4 \times 10^{-15}$) confirms that more BCCs reduce
Lorentzian selection probability across 4 consecutive scales, providing
the strongest empirical support for the BCC-Chain structure.

Secondary significant predictors at $n = 10$ include $n_{\mathrm{bridges}}$
($r = -0.640$), average PageRank ($r = +0.629$), vertex connectivity
($r = +0.382$), and girth ($r = +0.377$).

\subsection{Global Invariant Discriminant}
\label{subsec:global-discriminant}

To identify which graph invariants independently predict $P_{\mathrm{sel}}$
beyond the effect of graph size $(n, m)$, we compute partial Spearman
correlations on $503$ graphs ($n \leq 9$, all scales combined), controlling
for $n$ and $m$ simultaneously.

\begin{table}[ht]
\centering
\caption{Top global graph invariants predicting $P_{\mathrm{sel}}$, after
controlling for $n$ (vertices) and $m$ (edges). Partial Spearman correlation
on $503$ graphs. Spectral invariants (algebraic connectivity, spectral gap)
become non-significant after size control.}
\label{tab:global-discriminant}
\renewcommand{\arraystretch}{1.3}
\begin{tabular}{lccc}
\toprule
Invariant & Partial Spearman $r$ & Direction & Interpretation \\
\midrule
$n_{\mathrm{bccs}}$ (biconnected components) & $-0.270$ & $\downarrow$ & Fragmentation suppresses $P_{\mathrm{sel}}$ \\
$n_{\mathrm{bridges}}$ & $-0.267$ & $\downarrow$ & Bridge edges suppress $P_{\mathrm{sel}}$ \\
$\kappa(G)$ (vertex connectivity) & $+0.214$ & $\uparrow$ & Higher connectivity supports $P_{\mathrm{sel}}$ \\
Edge connectivity / min degree & $+0.190$ & $\uparrow$ & \\
Algebraic connectivity & $\approx 0$ & --- & NS after size control \\
Spectral gap $\lambda_2$ & $\approx 0$ & --- & NS after size control \\
\bottomrule
\end{tabular}
\end{table}

\begin{remark}[Spectral invariant collapse]
The most surprising finding in \cref{tab:global-discriminant} is the
\emph{collapse of spectral invariants} after controlling for graph size:
the algebraic connectivity $\lambda_2(L)$ and normalized spectral gap
become non-significant ($p > 0.05$) once $n$ and $m$ are held fixed.
The dominant signal is \emph{topological fragmentation vs.\ cohesion}:
the number of biconnected components and bridges captures the essential
structural information about $P_{\mathrm{sel}}$, consistent with the
BCC-Chain equivalence (Conjecture~\ref{cor:bcc-chain-complete}).
\end{remark}

A lean linear discriminant for $P_{\mathrm{sel}}$ using the top global
invariants is:
\begin{equation}
  P_{\mathrm{sel}}(G) \approx 0.597 - 0.061n + 0.026m
    - 0.028 \cdot n_{\mathrm{bccs}} - 0.003 \cdot n_{\mathrm{bridges}}
    + 0.033 \cdot \kappa(G),
  \label{eq:psel-discriminant}
\end{equation}
with $R^2 = 0.247$ on $439$ complete cases. The limited $R^2$ (below the
BCC-local model $R^2 = 0.351$) reflects genuine nonlinearity or granularity:
global invariants capture less information than local BCC structure. The
positive sign on $m$ (more edges $\to$ higher $P_{\mathrm{sel}}$) and
negative sign on $n_{\mathrm{bccs}}$ (more components $\to$ lower) are
consistent with the BCC-Chain interpretation.


%======================================================================
\section{$O(N)$ Universality}
\label{sec:on-universality}
%======================================================================

The Ising model ($N = 1$, discrete $\mathbb{Z}_2$ symmetry) has a
structural weakness in its Lorentzian selection: for dense random
graphs, the Fisher matrix develops large off-diagonal entries that
compete with the diagonal, degrading the spectral gap. We show that
this weakness is specific to discrete symmetry and is completely
eliminated by continuous $O(N)$ symmetry for any $N \geq 2$.

\subsection{$O(N)$ Spin Models on Graphs}
\label{subsec:on-models}

The $O(N)$ model on graph $G = (V, E)$ assigns to each vertex $v \in V$
a unit spin $\mathbf{s}_v \in S^{N-1} \subset \mathbb{R}^N$ and defines
the Boltzmann distribution:
\begin{equation}
  P(\{\mathbf{s}_v\} \mid \coupl)
  = \frac{1}{\partf(\coupl)} \exp\!\Bigl(
    \sum_{e=(i,j) \in E} J_e \, \mathbf{s}_i \cdot \mathbf{s}_j
  \Bigr).
  \label{eq:on-model}
\end{equation}
Special cases: $N = 1$ (Ising), $N = 2$ (XY), $N = 3$ (Heisenberg).
The edge observable is $\edgeo_e = \mathbf{s}_i \cdot \mathbf{s}_j$
(inner product, ranging in $[-1, +1]$).

The Fisher information matrix for the $O(N)$ model is defined identically:
$\Fish_{ab} = \Cov(\edgeo_a, \edgeo_b)$. The signed metric
$\Met = \Fish^{1/2}\sgn\Fish^{1/2}$ and spectral gap selection are
defined as before.

\subsection{$O(N)$ Tree Fisher Identity}
\label{subsec:on-tree-fisher}

\begin{theorem}[$O(N)$ Tree Fisher Identity]
\label{thm:on-tree-fisher}
For the $O(N)$ model on any tree $G = (V, E)$ with zero field,
the Fisher matrix is diagonal:
\begin{equation}
  \Fish = \diag\bigl(f(J_1, N), \ldots, f(J_m, N)\bigr),
  \label{eq:on-tree-fisher}
\end{equation}
where $f(J, N) = 1/N - (I_{N/2}(J)/I_{N/2-1}(J))^2$ and $I_\nu$ is the
modified Bessel function of order $\nu$.
\end{theorem}

\begin{proof}
The $O(N)$ tree factorization follows from the same Markov separation
argument as \cref{lem:tree-fisher}: the joint distribution of two
edge observables $\edgeo_a$ and $\edgeo_b$ on distinct edges of a
tree factorizes (tree Markov property holds for any pairwise model,
not just Ising). The $O(N)$ symmetry implies
$\langle \edgeo_a \edgeo_b \rangle = \langle \edgeo_a \rangle
\langle \edgeo_b \rangle$ by the same argument, so
$\Fish_{ab} = 0$ for $a \neq b$.
The diagonal entry $f(J, N) = \Var(\edgeo_e)$ follows from the $O(N)$
single-edge integral:
$\langle \edgeo_e^2 \rangle = 1/N$ (isotropy) and
$\langle \edgeo_e \rangle^2 = (I_{N/2}(J)/I_{N/2-1}(J))^2$.
\end{proof}

\begin{remark}
For $N = 1$: $f(J, 1) = \sech^2(J)$, recovering \cref{lem:tree-fisher}.
For $N = 2$: $f(J, 2) = 1/2 - (I_1(J)/I_0(J))^2$.
For $N \to \infty$: $f(J, \infty) \to 1/N \to 0$ (classical spins
become deterministic), and $\Fish \to 0$.
\end{remark}

\begin{corollary}[$O(N)$ Tree Flatness]
\label{cor:on-tree-flat}
For the $O(N)$ model on any tree with any $N \geq 1$, the signed
metric $\Met$ is flat and $\Ein_{ab} = 0$ for all couplings and all
sign assignments. The Statistical Vacuum Einstein Equation
(Theorem~\ref{thm:vacuum}) holds for all $O(N)$ models.
\end{corollary}

\subsection{Gap Elimination by Continuous Symmetry}
\label{subsec:gap-elimination}

The main empirical result of this section is:

\begin{theorem}[$O(N \geq 2)$ Gap Elimination]
\label{thm:on-gap}
For the Ising ($N = 1$) model, there exists a non-Lorentzian gap in
the Erd\H{o}s--R\'enyi random graph ensemble $G(n, p)$: for densities
$p \in [p_1(n), p_2(n)]$, the fraction of graphs with $q = 1$ optimal
(Lorentzian) selection drops below $50\%$. For $O(N \geq 2)$ models,
this gap is completely eliminated: $100\%$ Lorentzian selection on
all $719$ tested graphs across $36$ $(n, p)$ density-degree cells,
for all tested coupling strengths $J \in \{0.1, 0.3, 0.5, 1.0\}$.
\end{theorem}

This result is empirical (confidence 95\%): we have not proven
analytically that the gap vanishes for all $O(N \geq 2)$, but the
evidence is overwhelming across the tested parameter space.

\textit{Evidence:} The Ising non-Lorentzian gap arises because dense
graphs have Fisher matrices with large off-diagonal entries $f_o$,
competing with the diagonal. Continuous symmetry ($N \geq 2$) reduces
$f_o$ through the Mermin-Wagner mechanism: for $O(N \geq 2)$ models in
low dimensions, long-range order is suppressed, meaning correlations
$\langle \edgeo_a \edgeo_b \rangle$ decay faster with distance, keeping
$\Fish$ closer to diagonal even on dense graphs.

\subsection{Three-Phase Structure (Ising)}
\label{subsec:three-phase}

For the Ising model, the $G(n, p)$ ensemble exhibits a three-phase
structure as a function of $p$ at fixed $n$:
\begin{enumerate}
\item \textbf{Sparse Lorentzian phase} ($p < p_1(n)$):
  $\betti(G) = O(1)$ (tree-like), Fisher matrix near-diagonal,
  Lorentzian fraction $\to 1$.
\item \textbf{Non-Lorentzian gap} ($p_1(n) < p < p_2(n)$):
  Fisher matrix dense, off-diagonal entries compete with diagonal,
  Lorentzian fraction $< 50\%$.
\item \textbf{Dense rebound} ($p > p_2(n)$):
  Very dense graphs have approximately uniform Fisher matrices,
  for which the spectral gap analysis again selects $q = 1$.
\end{enumerate}

The gap width shows a decreasing trend with $n$:
\begin{equation}
  \Delta p(n) \equiv p_2(n) - p_1(n)
    \sim n^{-\alpha}, \quad \alpha \approx 5,
  \label{eq:gap-width}
\end{equation}
estimated from three data points ($n = 15, 20, 25$).  With only three
points and two free parameters, the power-law exponent $\alpha$ is not
reliably constrained; the qualitative conclusion --- that $\Delta p(n)$
decreases rapidly with $n$ --- is more robust than the specific
numerical value.  The trend suggests that for $n \gtrsim 50$,
the gap width falls below $\Delta p \sim 10^{-3}$, becoming
negligible in random graph ensembles.

\textbf{Physical interpretation:} The three-phase structure shows
that the Lorentzian selection mechanism is robust for sparse graphs
(tree-like observers) and very dense graphs (highly connected observers),
but has a transition region for intermediate density where the topology
changes from tree-like to loop-dominated but the loops are not yet
sufficiently ``averaged'' to restore selection. The $O(N \geq 2)$
universality shows this transition is an artifact of discrete symmetry.

\subsection{$N^*$ Census}
\label{subsec:nstar-census}

For each graph $G$, define $N^*(G)$ as the minimum $N$ for which
the $O(N)$ model achieves $100\%$ Lorentzian selection on $G$ at
all tested couplings. The $N^*$ census on 500 graphs ($n = 4$--$12$,
$\betti$ up to 41) finds:
\begin{itemize}
\item $99.8\%$ of graphs have $N^* \leq 5$.
\item Mean $N^* \approx 2.2$.
\item One exceptional graph ($n = 12$, $\betti = 17$) requires $N^* = 6$.
\item Rare outliers ($\sim 6\%$) with $N^* \geq 4$ are exclusively
  at edge connectivity $\kappa = 2$ (29/29 cases: a bridge joining
  two dense components).
\end{itemize}

\begin{remark}[Physical Significance]
The small value of $N^*$ (mean $\approx 2.2$) suggests that the
XY model ($N = 2$) is already sufficient for Lorentzian selection on
most graphs encountered in nature. The exceptional cases requiring
$N^* \geq 4$ all have a specific topological structure (bridge joining
dense components) that may be physically excluded by other constraints.
\end{remark}

\textbf{Figures:} Figure~5 ($O(N)$ gap elimination: heatmap of Lorentzian
fraction for Ising vs.\ $O(2)$ vs.\ $O(3)$ on $G(n,p)$ grid);
Figure~6 (three-phase structure: Lorentzian fraction vs.\ $p$ at
$n = 15, 20, 25$); Figure~7 ($\|\Ein\|$ vs.\ $\betti$ for census graphs).


%======================================================================
\section{Vielbein Identification of the Signed Fisher Metric}
\label{sec:vielbein}
%======================================================================

The signed Fisher metric $\Met = \Fish^{1/2}\sgn\Fish^{1/2}$ was
introduced in earlier work~\cite{Paper1,Paper5} as an ansatz motivated
by the PSD obstruction (the unsigned mass tensor $M = F^2$ cannot
produce Lorentzian signature). We now provide three independently
motivated arguments that single out this construction uniquely within
the class of vielbein-based metrics, significantly strengthening its
justification beyond the original ansatz.

\subsection{Three Independently Motivated Inputs}
\label{subsec:vielbein-inputs}

The vielbein formula from general relativity is:
\begin{equation}
  g_{\mu\nu} = e^a{}_\mu \, \eta_{ab} \, e^b{}_\nu,
  \label{eq:gr-vielbein}
\end{equation}
where $e^a{}_\mu$ is the frame field (vielbein) and $\eta_{ab}$ is
the flat Minkowski metric. Writing $e \equiv \Fish^{1/2}$ and
$\eta \equiv \sgn$, the signed Fisher metric becomes:
\begin{equation}
  \Met = e^T \eta\, e = \Fish^{1/2} \sgn \Fish^{1/2},
  \label{eq:vielbein-formula}
\end{equation}
in exact analogy with \eqref{eq:gr-vielbein}. We justify each component:

\medskip\noindent
\textbf{Input 1: Cencov's Uniqueness Theorem (determines $\Fish$).}
Cencov~\cite{Cencov1982} proved that the Fisher information metric is the
\emph{unique} Riemannian metric on the space of probability distributions
that is invariant (equivariant) under sufficient statistics maps.
Specifically, for any Markov kernel (statistical decision procedure)
$\kappa: X \to Y$, the Fisher metric satisfies
$\kappa^* \Fish_Y \leq \Fish_X$ (data processing inequality), with
equality iff $\kappa$ is a sufficient statistic.

This uniqueness singles out $\Fish$ as the canonical metric on the
statistical manifold. Any other metric on the coupling space $\coupl$
would fail the data processing inequality, making it non-natural from a
statistical standpoint. Therefore, the Riemannian part of the geometry
\emph{must} be $\Fish$.

\medskip\noindent
\textbf{Input 2: Spectral Gap Selection (determines $\sgn$).}
The sign matrix $\sgn$ is not a free parameter but is determined by the
spectral gap weighting functional. Paper~\cite{Paper5} (Theorems A--C)
proves that
\begin{equation}
  \sgn^* = \arg\max_{q \geq 0} W(q) = \arg\max_{q} \beta_c(q) \cdot L_{\mathrm{gap}}(q)
  \label{eq:sign-selection}
\end{equation}
uniquely selects $q = 1$ (one negative sign) for sparse Ising observers
at subcritical coupling. The $q = 1$ assignment is not a choice but
a consequence of maximizing the spectral stability of the Lorentzian
metric.

\medskip\noindent
\textbf{Input 3: Canonical Symmetric Vielbein (determines $e$).}
Given $\Fish$ (a symmetric positive-definite matrix), the vielbein
$e$ satisfying $e^T e = \Fish$ is not unique: any $e' = Oe$ with
$O^T O = I$ also satisfies $(e')^T e' = \Fish$. To select a canonical
representative, we require $e$ to be symmetric: $e^T = e$.

The unique symmetric positive-definite square root $\Fish^{1/2}$ is
then the canonical vielbein. It is characterized by:
\begin{enumerate}
\item \textit{Symmetry:} $e = e^T$ (no preferred frame direction).
\item \textit{Positivity:} $e$ is positive definite (preserves the
  orientation of $\Fish$).
\item \textit{Uniqueness:} For any PD matrix $A$, there is a unique
  symmetric PD $B$ with $B^2 = A$ (given by $B = A^{1/2}$ via
  eigendecomposition).
\end{enumerate}

\begin{remark}[Why not $e = \Fish$?]
One might ask: why not use $e = \Fish$ directly, giving
$\Met = \Fish \sgn \Fish = FSF$? The answer is that this is a different
construction (the ``signed mass tensor'' $M = FSF$ of \cite{Paper1})
which does not have the correct vielbein interpretation: $FSF$ is not
of the form $e^T \eta e$ with $e^T e = \Fish$ (that would require $e = \Fish^{1/2}$).
The signed mass tensor $FSF$ is related by $FSF = F^{1/2}(F^{1/2}SF^{1/2})F^{1/2}
= F^{1/2} \Met F^{1/2}$, a further conjugation by $F^{1/2}$.
\end{remark}

\subsection{Gauge Invariance}
\label{subsec:gauge-invariance}

The vielbein formula $\Met = e^T \eta e$ possesses a gauge symmetry:
any transformation $e \to e' = Oe$ with $O$ preserving $\eta$
(i.e., $O^T \eta O = \eta$, so $O \in O(m-1, 1)$, the Lorentz group)
gives $\Met' = (Oe)^T \eta (Oe) = e^T (O^T \eta O) e = e^T \eta e = \Met$.
The physics (the metric $\Met$) is invariant under Lorentz rotations of
the vielbein.

As a code-consistency check (since gauge invariance is guaranteed
algebraically by $O^T \sgn\, O = \sgn$ for $O \in O(m{-}1, 1)$),
we verify numerically: for 100 random Lorentz rotations $O$, the
gauged metric $\Met' = (O\Fish^{1/2})^T \sgn (O\Fish^{1/2})$
satisfies $\|\Met' - \Met\| < 10^{-10}$.

\subsection{Comparison with the Literature}
\label{subsec:vielbein-literature}

The construction $\Met = \Fish^{1/2}\sgn\Fish^{1/2}$ does not appear
in the existing information geometry literature (verified by 10+
database searches). The closest antecedents are:

\begin{enumerate}
\item \textit{Greensite~\cite{Greensite1993}}: Proposes a dynamical
  mechanism for signature selection in which the metric signature is
  a dynamical variable in the gravitational path integral. The vielbein
  $e$ is treated as a fluctuating field. Our construction fixes $e = \Fish^{1/2}$
  by statistical principles, giving a non-dynamical (observer-determined)
  signature.

\item \textit{Bures metric~\cite{Uhlmann1976}}: The Bures metric on the
  space of density matrices is $ds^2_{\mathrm{Bures}} = \tfrac{1}{2}
  \mathrm{Tr}(d\rho \cdot \rho^{-1/2} d\rho^{1/2})$, which involves
  the same sandwich structure $\rho^{1/2} A \rho^{1/2}$. However, the
  Bures metric is positive semi-definite (Riemannian), not Lorentzian,
  because the sign matrix is replaced by the identity.

\item \textit{$\alpha$-connections (Amari~\cite{Amari1998})}: The family
  of $\alpha$-connections on statistical manifolds generalizes the Levi-Civita
  connection of $\Fish$. None of these connections produces a Lorentzian
  metric.
\end{enumerate}

The novelty of our construction lies in the combination:
(a) sandwich form $e^T \eta e$ (from GR vielbein theory),
(b) canonical choice $e = \Fish^{1/2}$ (from Cencov + symmetric square root),
(c) indefinite $\eta = \sgn$ with $q = 1$ (from spectral gap selection).

\subsection{Confidence Assessment}
\label{subsec:vielbein-confidence}

\begin{remark}[Confidence Assessment]
Prior to the vielbein identification, the signed metric $\Met = \Fish^{1/2}\sgn\Fish^{1/2}$
was an ansatz: motivated by the PSD obstruction, but without a principled
derivation. Confidence: 25--30\% (``why this specific construction and
not another?'').

After the vielbein identification, all three components ($\Fish$, $\sgn$,
$\Fish^{1/2}$) are independently motivated. The construction receives
support from three independent arguments (Cencov uniqueness, spectral gap
selection, canonical symmetric vielbein), making it the unique candidate
within the class of vielbein-based metrics. Confidence: 70--80\%
(``supported by independent principles, not tuned'').
A full derivation from first principles (proving that the coupling
manifold \emph{must} carry a Lorentzian geometry of this form) remains open.

The remaining uncertainty (20--30\%) stems from the physical interpretation:
why should the coupling manifold carry a Lorentzian geometry in the first
place? The answer requires connecting to the full Vanchurin framework
\cite{Paper1}, which is outside the scope of this paper.
\end{remark}

\textbf{Appendix~\ref{app:vielbein}}: Numerical verification of gauge
invariance, uniqueness of symmetric square root, and explicit computations
for $C_4$ and $P_4$.


%======================================================================
\section{Observer Trajectory Dynamics}
\label{sec:dynamics}
%======================================================================

We study how the coupling vector $\coupl = (J_1, \ldots, J_m)$
evolves when the observer is ``self-consistent'': the observer vector
$\obs(\coupl)$ (timelike eigenvector of $\Met$) drives the evolution
of the couplings that define the observer itself. Three candidate
laws of motion are compared on $C_4$ with $200$ steps of size
$d\tau = 0.005$.

\subsection{Three Laws of Motion}
\label{subsec:three-laws}

\begin{definition}[Observer Flow, Law A]
\label{def:law-a}
The \emph{observer flow} is the dynamical law
\begin{equation}
  \frac{dJ_a}{d\tau} = \obs_a(\coupl),
  \label{eq:law-a}
\end{equation}
where $\obs_a$ is the $a$-th component of the timelike eigenvector
of $\Met(\coupl)$ (normalized: $\Met_{ab}\obs^a\obs^b = -1$).
The coupling evolves in the direction of the observer itself: the
observer ``chooses'' the direction that maximally aligns with its
own timelike direction.
\end{definition}

\begin{definition}[Geodesic Flow, Law B]
\label{def:law-b}
The \emph{geodesic flow} satisfies the geodesic equation:
\begin{equation}
  \frac{d^2 J_a}{d\tau^2} + \Chr{a}{bc}(\coupl)
    \frac{dJ_b}{d\tau}\frac{dJ_c}{d\tau} = 0,
  \label{eq:law-b}
\end{equation}
with initial velocity $dJ_a/d\tau\big|_{\tau=0} = \obs_a(\coupl_0)$.
This is the ``straightest possible'' motion in the coupling manifold,
generalizing straight-line motion in flat space.
\end{definition}

\begin{definition}[Natural Gradient Flow, Law C]
\label{def:law-c}
The \emph{natural gradient flow} is:
\begin{equation}
  \frac{dJ_a}{d\tau} = -(\Fish^{-1})_{ab} \frac{\partial \mathcal{L}}{\partial J_b},
  \label{eq:law-c}
\end{equation}
where $\mathcal{L}(\coupl) = -A(\coupl)/m$ (negative log-partition
function per edge, a proxy for ``simplicity'' of the distribution).
This is Amari's natural gradient ascent~\cite{Amari1998} adapted
to the coupling space, using the Fisher metric as the preconditioning.
\end{definition}

\subsection{Arrow of Time}
\label{subsec:arrow-of-time}

All three laws produce a monotone thermodynamic arrow on $C_4$:
the free energy $A(\coupl)/m$ decreases monotonically along trajectories
under Laws A and B. The expansion $\expan(\tau)$ is consistently negative
(contracting, consistent with the Raychaudhuri analysis of
\cref{subsec:cycle-focusing}).

The Lorentzian stability is maintained throughout all trajectories:
the metric $\Met(\coupl(\tau))$ retains exactly $q = 1$ (one negative
eigenvalue) at all 200 steps for all three laws (100\% stability).

\subsection{Expansion Bounce and Singularity}
\label{subsec:bounce}

The qualitative difference between the laws appears at late times
($\tau \gtrsim 0.5$).  We emphasize that the following observations
are based on $C_4$ trajectories (200 steps); generalization to larger
graphs and longer evolution requires further study.

\begin{enumerate}
\item \textbf{Law A (observer flow):} The expansion $\expan(\tau)$
  decreases from $\expan_0 \approx -0.41$ to a minimum, then
  \emph{crosses zero} at step $\approx 170$ (i.e., $\tau \approx 0.85$),
  and becomes slightly positive. By analogy with bouncing cosmologies,
  we call this an \emph{expansion bounce}: the contracting phase
  transitions to a (mildly) expanding phase.
  After the bounce, $\expan$ oscillates near zero.

\item \textbf{Law B (geodesic):} The expansion $\expan(\tau)$ decreases
  monotonically, crossing zero at a later step, with oscillatory behavior.
  No clear bounce structure.

\item \textbf{Law C (natural gradient):} The expansion $\expan(\tau)$
  diverges to $-\infty$: a focusing singularity. The Raychaudhuri
  equation \eqref{eq:raychaudhuri-cycle} has $d\expan/d\tau \to -\infty$
  because the natural gradient flow drives the couplings toward strong-coupling
  fixed points where the Fisher metric degenerates.
\end{enumerate}

The expansion bounce in Law A is the most physically significant feature
observed on $C_4$: it shows that the observer flow can self-regulate,
preventing the singularity that occurs under the natural gradient.

\subsection{Viability Self-Selection}
\label{subsec:self-selection}

We define the \emph{Lorentzian viability advantage}:
\begin{equation}
  \Delta W(\coupl) = W(q{=}1;\coupl) - W(q{=}2;\coupl),
  \label{eq:delta-w}
\end{equation}
where $W(q;\coupl) = \beta_c(q;\coupl) \cdot L_{\mathrm{gap}}(q;\coupl)$
is the spectral gap weighting at coupling $\coupl$. A positive $\Delta W$
means the Lorentzian assignment ($q = 1$) is more stable than the
next-best assignment ($q = 2$).

\begin{conjecture}[Viability Self-Selection]
\label{thm:self-selection}
Under observer flow (Law A), the viability advantage $\Delta W$
increases monotonically along trajectories:
\begin{equation}
  \frac{d\Delta W}{d\tau} \geq 0
  \quad \text{(under Law A, empirically)}.
  \label{eq:self-selection}
\end{equation}

\emph{Numerical evidence.} On $C_4$: $\Delta W$ increases from $0.094$
to $0.161$ over 200 steps (39 increases, 0 decreases).
On $C_6$: $\Delta W$ increases from $0.025$ to $0.050$
(19 increases, 1 decrease).
Under Law C, $\Delta W$ collapses from $0.094$ to $0.003$ (viability
destroyed). Under Law B, $\Delta W$ mostly increases ($28\uparrow,
11\downarrow$) but non-monotonically.
\end{conjecture}

This self-selection principle provides a \emph{law selection criterion}:
Law A (observer flow) is the unique law satisfying $d\Delta W/d\tau \geq 0$
(among the three tested). The observer flow ``earns'' its own viability:
it strengthens the Lorentzian selection that justifies its existence.

\begin{remark}[Physical Interpretation]
The self-selection principle has a clear interpretation: the observer
vector $\obs(\coupl)$ points in the direction of maximum Lorentzian
viability (by its definition as the timelike eigenvector). Moving the
couplings in this direction ($dJ_a/d\tau = \obs_a$) therefore increases
the viability measure $\Delta W$. This is a ``gradient flow'' on $\Delta W$,
but in the non-Euclidean geometry of the coupling manifold.

The collapse of $\Delta W$ under Law C (natural gradient) indicates that
Amari's natural gradient, while optimal for learning tasks, is \emph{not}
consistent with maintaining Lorentzian geometry: it destroys the spectral
gap structure that the observer requires.
\end{remark}

\subsection{Two-Sector Structure}
\label{subsec:two-sectors}

Long-time integration (10000 steps, $d\tau = 0.005$) reveals that
trajectories separate into two qualitatively distinct sectors based
on initial vorticity:

\begin{description}
\item[Sector A (symmetric, $\|\vort\|_0 \approx 0$):]
  Initial conditions with zero or near-zero vorticity (e.g., uniform
  coupling $J_a = J$ on $C_4$, giving $\|\vort\|_0 \approx 0$).
  \begin{itemize}
  \item $\Delta W$ monotonically increases.
  \item $\|\coupl\| \to \infty$ (couplings diverge along the symmetric
    direction).
  \item Strong arrow of time: $A(\coupl)/m$ monotone.
  \item Sector A is a measure-zero set (symmetric initial conditions are
    non-generic).
  \end{itemize}

\item[Sector B (asymmetric, $\|\vort\|_0 \gtrsim 10^{-2}$):]
  Initial conditions with nonzero vorticity (e.g., non-uniform
  couplings $J_a \neq J_b$ on $C_4$).
  \begin{itemize}
  \item $\Delta W$ non-monotone but remains positive.
  \item $\|\coupl\|$ is bounded (``asymmetry brake'': the nonzero
    vorticity creates component cancellations in $dJ_a/d\tau = \obs_a$,
    preventing runaway).
  \item $\expan$ oscillates quasi-periodically ($\approx 26$ sign changes
    per 10000 steps).
  \item Sectors are \emph{permanent}: $\|\vort\|$ does not approach zero
    over 10000 steps (no sector transition under the dynamics).
  \end{itemize}
\end{description}

The critical threshold separating the two sectors is $\|\vort\|_{\mathrm{crit}}
\approx 10^{-2}$: initial vorticity below this gives Sector A behavior,
above it gives Sector B.

\begin{remark}[Generic Universe is Sector B]
Since Sector A requires exactly symmetric initial conditions (measure zero
in the coupling space), a ``generic'' observer lives in Sector B: stable
couplings, oscillating expansion, no arrow of time in the thermodynamic
sense. This is consistent with the idea that a cosmological arrow of time
requires special (low-entropy) initial conditions.
\end{remark}

\subsection{Orientation Selection}
\label{subsec:orientation}

For graphs with multiple $q = 1$ sign assignments (different choices of
which edge is timelike), the observer flow self-selects among them via
the viability advantage. We test this on the Diamond graph
($K_4$ minus one edge, $|E| = 5$):

\begin{itemize}
\item Edge $(0,1)$: gives $\|\vort\| \approx 10^{-13} \approx 0$
  (the unique assignment preserved by the $\mathbb{Z}_2$ automorphism
  $2 \leftrightarrow 3$ of the Diamond graph).
  Self-selection: $47\uparrow / 3\downarrow$.

\item Edges $(0,2), (0,3), (1,2), (1,3)$: give $\|\vort\| \approx 3 \times 10^{-2}$.
  Self-selection: $50\uparrow / 0\downarrow$ (strictly monotone).
\end{itemize}

Surprisingly, the $\vort \neq 0$ assignments have \emph{stronger}
self-selection ($50/0$) than the $\vort = 0$ assignment ($47/3$). This
shows that $\vort = 0$ is \emph{sufficient} but not \emph{necessary}
for self-selection. The physical mechanism is different: the $\vort = 0$
assignment on edge $(0,1)$ is in Sector A (near-symmetric), while the
$\vort \neq 0$ assignments are in Sector B, which has a different (more
stable) form of self-selection.

The key negative result (NR-013): the modified viability
$V = \Delta W - \lambda\|\vort\|^2$ does \emph{not} improve self-selection
over $\Delta W$ alone. Vorticity grows along self-selecting trajectories
rather than being penalized.

\textbf{Figures:} Figure~9 ($J(\tau)$, $\expan(\tau)$, $A(\tau)/m$ for
three laws on $C_4$); Figure~10 (expansion bounce: $\expan$ crossing
zero for Law A vs.\ singularity for Law C); Figure~11 (Lorentzian
stability: eigenvalue spectrum of $\Met$ along trajectory); Figure~12
(viability self-selection: $\Delta W(\tau)$ for Laws A, B, C).


%======================================================================
\section{Discussion}
\label{sec:discussion}
%======================================================================

\subsection{Summary of Results}

This paper establishes that the curvature of the signed Fisher metric
$\Met = \Fish^{1/2}\sgn\Fish^{1/2}$ on the coupling manifold of an
Ising observer is exactly controlled by the cycle structure of the
underlying graph. The correspondence is precise and computable.

The central theorem (Theorem~\ref{thm:vacuum}, Statistical Vacuum
Einstein Equation) proves the forward direction: if $G$ is a tree, then
$\Ein_{ab} = 0$ for all couplings and all sign assignments (at $h = 0$).
The converse direction is stated as Conjecture~\ref{conj:vacuum-converse}
and supported by analytic results for $C_3$ and $C_4$, transfer-matrix
computation for $C_5$--$C_7$, and an exhaustive census on $839$ connected
graphs with $n \leq 7$ vertices (zero counterexamples). A complete
analytic proof for general non-trees remains open.

The $C_4$ exact formula \eqref{eq:c4-scalar} is, to our knowledge,
the first explicit curvature formula for a Lorentzian signed Fisher
metric in the information geometry literature (we note that the
Ruppeiner geometry program computes curvature of positive-definite
Fisher-type metrics in statistical mechanics, but these do not involve
signed metrics with indefinite signature). It reveals the
\emph{factor-of-two law} (numerically verified, $\Scal_{\Met} = \Scal_{\Fish}/2$) and
the strong monotone growth of curvature with coupling.

The vielbein identification (\cref{sec:vielbein}) provides a principled
justification for the signed Fisher metric $\Met = e^T \eta\, e$: the
ansatz receives support from three independent arguments --- Cencov's
uniqueness theorem (determining $\Fish$), spectral gap selection
(determining $\sgn$), and the canonical symmetric square root
(determining $e = \Fish^{1/2}$). While these arguments are mutually
reinforcing and single out the construction uniquely within the class
of vielbein-based metrics, a full derivation from first principles
remains open.

The observer kinematics (\cref{sec:raychaudhuri}) provide a concrete
realization of the gravitational analog: tree observers are inertial
(zero expansion, zero shear, zero vorticity), cycle observers are
focusing (negative expansion, positive Raychaudhuri source), and
vorticity is controlled by the cycle rank. The classification is
verified on 839 graphs, 8348 sign assignments.

The BCC-Chain Conjecture (\cref{subsec:bcc-chain-complete}) characterizes
vorticity: the forward direction ($\betti \leq 1$ implies $\vort = 0$)
is proved in Theorem~\ref{thm:vorticity}(i). The converse ($\vort = 0$
implies $\betti \leq 1$) is stated as Conjecture~\ref{thm:bcc-reverse},
supported by perturbation theory and explicit construction for minimal
obstructions (Diamond, $K_4$, Theta), with census verification on all
$839$ tested graphs ($0$ counterexamples).

The observer consistency formalization
(\cref{subsec:observer-consistency}) gives the BCC-Chain equivalence a direct
multi-observer interpretation.  Three levels of inter-observer consistency
are defined (causal, temporal, full metric), forming a strict hierarchy.
The central result (Theorem~\ref{thm:topology-consistency}) shows that
\emph{automatic temporal consistency} --- every observer pair independently
admits a global time foliation on their shared BCC --- is equivalent to the
BCC-chain condition $\max_i \betti(B_i) \leq 1$.  This is the discrete
analog of the Geroch splitting theorem: global hyperbolicity of the
statistical spacetime is determined by the cycle topology of the graph,
not by the coupling values or observer choices.  Graphs with BCC cycle
rank $\leq 1$ are temporally ordered; those with some BCC having
$\betti \geq 2$ are temporally frustrated, with obstruction to global time
foliation for generic observers.

The $C_n$ Lorentzian selection theorem
(Theorem~\ref{thm:cn-lorentzian}) proves
$W_{\mathrm{best}}(1) > W_{\mathrm{best}}(q)$ for cycle graphs
(sparsest $2$-regular), providing the first analytical proof of
Lorentzian selection for a specific non-tree graph family.
The $K_n$ Lorentzian selection
(Conjecture~\ref{thm:kn-lorentzian}) extends this to complete graphs
(densest possible), verified by exact enumeration for $K_4$--$K_9$
but not yet proved analytically.

The $P_{\mathrm{sel}}$ atlas at $n = 10$ (\cref{sec:psel-atlas})
confirms $n_{\mathrm{bccs}}$ as the dominant predictor for the fourth
consecutive scale, with $0$ BCC-Chain counterexamples at $n = 10$.
The global invariant discriminant identifies topological fragmentation
(BCC count, bridges) as the primary independent predictor of $P_{\mathrm{sel}}$,
with spectral invariants collapsing after size control.

The $O(N \geq 2)$ universality (\cref{sec:on-universality}) shows
that the Lorentzian selection mechanism is robust against the Ising
non-Lorentzian gap: continuous symmetry completely eliminates the gap,
providing a physical mechanism (Mermin--Wagner-type near-diagonality)
for the robustness of spacetime emergence.

The observer trajectory dynamics (\cref{sec:dynamics}) complete the
picture: the observer flow (Law A) exhibits an expansion bounce,
a monotone arrow of time (in Sector A), and most importantly, the
viability self-selection principle $d\Delta W/d\tau \geq 0$. The
observer flow is the unique law (among those tested) that strengthens
its own Lorentzian viability.

The phase transition curvature divergence (\cref{subsec:phase-transition})
provides the first graph-based confirmation of the Ruppeiner program:
$|\Scal(J_c)| \sim L^{5.65}$ on $L \times L$ grids (open BC) at the Ising
critical coupling. The curvature divergence is universal across lattice types
(square, triangular, hexagonal, and 3D cubic), but the \emph{sign} of
$\Scal$ at $J_c$ is lattice-dependent (\cref{rem:lattice-types,rem:3d-lattice}).

For periodic boundary conditions (tori), the curvature scaling exponent is
much smaller and governed by the universality-class formula
$d_R^{\mathrm{PBC}} = (d\nu + 2\eta)/(d\nu + \eta)$ (Conjecture~\ref{conj:pbc-scaling};
\cref{subsec:pbc-scaling}). For 2D Ising tori ($L = 3, \ldots, 8$, exact TM),
$d_R^{\mathrm{PBC}} = 1.1110$ at $L = 7{\to}8$, consistent with the predicted $10/9$.
For 3D Ising tori (hybrid TM+MCMC, $L \leq 7$), the measured
$d_R^{\mathrm{eff}}(6{\to}7) = 1.043 \pm 0.016$ is $1.52\sigma$ from the
predicted $1.0188$ (statistically consistent). The boundary condition thus
qualitatively changes the exponent: open-BC grids show $d_R^{\mathrm{open}} \approx 2.15$
driven by boundary-enhanced correlations, while tori show $d_R^{\mathrm{PBC}} \approx 1.11$
governed purely by bulk anomalous dimensions.
Lattices with high coordination number ($z = 6$: triangular, 3D cubic)
have $\Scal > 0$ at $J_c$; those with lower $z$ ($= 3, 4$: hexagonal,
square) have $\Scal < 0$. The mechanism is the position of the curvature
sign-change coupling $J_{\Scal}$ relative to $J_c$: higher $z$ pushes
$J_c$ below $J_{\Scal}$, placing the critical point in the
positive-curvature regime. This pattern is formalized in
\cref{conj:coord-sign}, which predicts a coordination number
threshold $z^* \in (5, 6]$ separating negative from positive
curvature at criticality (the icosahedron at $z = 5$ yields
$\Scal < 0$, narrowing the bound from below).

A complementary axis is the Potts model parameter $q$
(\cref{rem:potts-sign}): on the $3 \times 3$ grid,
the curvature at $J_c^{(q)}$ transitions from negative ($q = 2$:
$\Scal = -3.06$; $q = 3$: $\Scal \approx -0.08$) to positive
($q = 4$: $\Scal = +0.85$), and remains positive through $q = 8$
($\Scal \approx +0.7$ on $2 \times 3$ grids). This reveals $q$
as a second control parameter that governs the curvature sign at
criticality through the same $J_{\Scal}/J_c$ mechanism. The
positive curvature plateau for $q \geq 4$ suggests a sharp
transition at a critical $q^* \in (3, 4)$.

The Riemann decomposition identity
(Theorem~\ref{thm:riem-decomp} for cycle graphs, plus an empirical
identity beyond cycles) provides the algebraic backbone:
$\Riem|_{\mathrm{lin}} = -2\,\Riem|_{\mathrm{quad}}$ is proven on $C_n$
and verified numerically for all single-block exponential families tested.

\subsection{The Curvature--Correlation Correspondence}

The central conceptual contribution is the precise identification:
\begin{equation}
  \boxed{\text{Curvature of } \Met \text{ is encoded in correlations of } \edgeo_a.}
  \label{eq:curvature-correlations}
\end{equation}
This is not merely analogical. The Fisher matrix $\Fish_{ab} = \Cov(\edgeo_a, \edgeo_b)$
is the covariance of edge observables, and on trees $\Fish$ is diagonal and
$\Met$ is flat (Theorem~\ref{thm:tree-flat}).
Conversely, on non-trees, cycle constraints induce correlations among edge
observables and (conjecturally) source nonzero curvature
(Conjecture~\ref{conj:vacuum-converse}; \cref{rem:vacuum-proof-status}).
Diagonal $\Fish$ means pairwise uncorrelated edge observables; on trees
this coincides with full statistical independence.

The physical meaning is: \emph{flat spacetime is independent measurement}.
When the edge observables can be measured without interfering with each other
(tree graph, no loops), the geometry is flat. When there are loops, any
measurement of one edge induces correlations with other edges, and the
geometry is typically curved (supported by \cref{rem:vacuum-proof-status}).

This gives a new interpretation of the Einstein vacuum equation
$G_{\mu\nu} = 0$ in this setting: on trees, it is the statement that the
degrees of freedom of the observer are statistically independent.
Conversely, cycle-induced correlations act as an effective statistical
source of curvature.

\subsection{Comparison with Jacobson's Derivation}

Jacobson~\cite{Jacobson1995} derived the Einstein equation from the
Clausius relation $\delta Q = T\, dS$ applied to local Rindler horizons,
treating spacetime as a thermodynamic system. Our derivation of the
statistical vacuum equation is complementary: we do not derive the
Einstein \emph{equation} (with sources) but rather characterize the
\emph{vacuum solutions} ($\Ein_{ab} = 0$) and their statistical meaning.

The key structural similarity is: in both cases, the geometric objects
(metric, curvature, Einstein tensor) are related to thermodynamic or
information-theoretic objects (entropy, Fisher information, covariance).
In Jacobson's work, $\nabla_a T^{ab} = 0$ follows from $\delta S \propto
\delta A$ (entropy proportional to area). In our work, $\Ein_{ab} = 0$
holds on trees where $\Cov(\edgeo_a, \edgeo_b) = 0$ for $a \neq b$
(Fisher matrix diagonal), and non-trees exhibit off-diagonal Fisher
correlations and nonzero $\Ein$ in all tested cases.

The Lorentzian signature of $\Met$ plays a crucial role: it distinguishes
the timelike direction (the observer's direction of maximal Fisher
response) from the spacelike directions. The Raychaudhuri equation
in this Lorentzian geometry gives focusing (negative $d\expan/d\tau$)
for cycles, analogous to gravitational attraction.

\subsection{Connection to Holographic Entanglement}

The correspondence between tree flatness and statistical independence
connects to the Ryu--Takayanagi formula~\cite{RyuTakayanagi2006} via
the following identification:
\begin{enumerate}
\item \emph{Boundary theory} = Ising model on a graph.
\item \emph{Boundary region} = a subset $A \subset E$ of edges.
\item \emph{Entanglement entropy} of $A$ = $H(\{\edgeo_a\}_{a \in A})$,
  the Shannon entropy of the edge observable distribution restricted to $A$.
\item \emph{Bulk geometry} = signed Fisher metric $\Met$ on the coupling
  manifold.
\item \emph{Minimal surface} = submanifold of the coupling manifold
  minimizing the $(m-1)$-dimensional area measure of $\Met$.
\end{enumerate}

The Ryu--Takayanagi formula would then read:
$H(A) = \text{Area}(\gamma_A) / (4 G_N)$, where $\gamma_A$ is a minimal
$(m-|A|-1)$-dimensional submanifold in the coupling space and $G_N$ is
an effective ``Newton constant'' (the Fisher metric normalization).

For tree graphs, the minimal surface degenerates (flat geometry, area $= 0$),
consistent with the tree factorization $H(A) = \sum_{a \in A} H(\edgeo_a)$
(edge entropies are additive on trees, no mutual information). For cycles,
the minimal surface is nontrivial, consistent with $\Scal \neq 0$.

This is a sketch; the precise holographic dictionary requires further work.

\subsection{Open Problems}

\begin{enumerate}
\item \textbf{Analytical self-selection proof.} Prove
  $d\Delta W/d\tau \geq 0$ analytically for observer flow in Sector A.
  Currently verified numerically on $C_4$ and $C_6$. The key missing
  ingredient is a formula relating $\partial_a \Delta W$ to the observer
  vector components $\obs_a$.

\item \textbf{$J \to \infty$ regularity (Sector A).} In Sector A,
  $\|\coupl\| \to \infty$ under observer flow. Does the trajectory
  approach a well-defined limit (e.g., a fixed point of the observer
  flow at $J_a \to \infty$)? The asymmetry brake in Sector B prevents
  this for generic initial conditions.

\item \textbf{Linearized Einstein equation.} \emph{(Substantially
  resolved---numerically.)} Perturb around a uniform-coupling reference
  point $J_a = J_0$ for all edges: reduce one edge coupling by
  $\delta$, so $J_{\mathrm{cycle}} = J_0 - \delta$. Numerical computation
  on $C_4$, $C_5$, and $P_6$ with a shortcut confirms
  $\|\Ein(\coupl_0) - \Ein(\coupl_0 - \delta\, e_a)\|_F = A\,\delta^\alpha$
  with $\alpha \approx 0.96$--$0.98$ and $R^2 > 0.999$ across
  $\delta \in [0.005, 0.99]$ (15 data points per graph).
  The linearity holds far beyond the perturbative regime.
  \emph{Caveat}: the naive ``tree plus weak cycle edge'' parameterization
  (all tree edges at $J$, cycle edge at $\epsilon \to 0$) is singular---the
  Fisher metric degenerates at $J_{\mathrm{cycle}} = 0$ (zero covariance),
  so the correct linearization must be performed around a nondegenerate
  reference point. The tree limit is a \emph{boundary} of the Fisher
  manifold. The open part is an analytical derivation of $\Ein_{ab}
  \approx \kappa\, T_{ab}$ with an explicit identification of the
  source tensor $T_{ab}$.

\item \textbf{AdS/CFT bootstrap.} The Lashkari--Van Raamsdonk
  framework~\cite{Lashkari2014} shows that entanglement entropy perturbations
  are dual to canonical energy perturbations in the bulk. In our setting:
  can we derive the linearized Einstein equation from the first law of
  entanglement entropy $\delta S_A = \delta \langle K_A \rangle$?

\item \textbf{Dimensionality selection (``$3{+}1$ from Trees'').}
  \emph{(Partially negative.)} Explicit computation of the spectral
  dimension $d_s$ of finite Cayley trees (branching number $k = 2$--$4$,
  depth $h = 2$--$5$, up to $n = 1365$ vertices) yields $d_s \approx
  1.0$--$1.5$, \emph{not} $d_s = 3$. The effective $d_s$ saturates near
  the Alexander--Orbach value $4/3$ for random trees, far below the
  $d_s = 3$ needed. Moreover, the Tree Fisher Identity
  $\Fish = \operatorname{sech}^2(J)\, I_m$ means all edges carry equal
  Fisher information---the information geometry ``does not see'' the tree
  structure. If there is a ``$3{+}1$ from trees'' mechanism, it cannot
  come from the Fisher matrix eigenstructure alone. A positive result
  would require either the infinite Bethe lattice limit (where $d_s = 3$
  may hold) or a different notion of effective dimension.

\item \textbf{Uniqueness of observer flow.} Is Law A the unique
  law satisfying $d\Delta W/d\tau \geq 0$ for all initial conditions
  and all graphs? A variational characterization would strengthen
  the self-selection principle from an observation to a theorem.

\item \textbf{Sector B physics.} The generic universe lives in Sector B
  (stable couplings, oscillating expansion, no arrow of time). Does this
  correspond to a physically viable cosmological scenario? What is the
  thermodynamic interpretation of the quasi-periodic oscillation of
  $\expan$? Can graph evolution (changing the graph topology, not just
  the couplings) mediate transitions between Sectors A and B?

\item \textbf{$C_n$ exact formulas and universal Riemann decomposition.}
  \emph{(Resolved; extended to all graphs.)} Exact closed-form
  expressions for $\alpha(n, J)$, $\beta(n, J)$, and $\Scal(n, J)$ are
  now derived for $C_3$--$C_6$ as rational functions of $w = e^J$
  (\cref{app:cn-exact}). The Riemann decomposition
  $\Riem_{\mathrm{linear}} = -2\,\Riem_{\mathrm{quad}}$ initially proved
  for $C_n$ is now verified as universal for \emph{all} graphs
  (Theorem~\ref{thm:riem-decomp}): cycles, complete graphs, bipartite,
  Petersen, grids, wheels, and random graphs, with zero violations in
  42 test cases. An analytical proof for the general case remains open.
  For the large-$n$ asymptotics of $C_n$: $\Scal \to 0$ as $n \to \infty$
  (consistent with absence of phase transition in 1D Ising; the Fisher
  manifold is asymptotically flat).

\item \textbf{Universal Lorentzian selection.} Theorem~\ref{thm:cn-lorentzian}
  proves $W_{\mathrm{best}}(1) > W_{\mathrm{best}}(q)$
  for $C_n$, and Conjecture~\ref{thm:kn-lorentzian} extends this to $K_n$
  (verified by enumeration for $K_4$--$K_9$). A natural conjecture is
  that this holds for \emph{all} graphs with $m \geq 4$ edges: for any
  connected graph $G$ with uniform Ising coupling $J > 0$,
  $W_{\mathrm{best}}(1) > W_{\mathrm{best}}(q)$ for all $q \geq 2$.
  The analytical proof for $C_n$ uses only $2f_o > 0$; the $K_n$
  argument uses the Johnson scheme but the analytical inequality from
  the cubic secular equation remains open. A general proof would require
  characterizing which Fisher matrix structures can break the inequality,
  and no counterexample has been found in any tested graph.

\item \textbf{BCC-Chain equivalence: completing the analytical proof.}
  Conjecture~\ref{thm:bcc-reverse} is supported by explicit construction for
  minimal obstructions and census verification. The remaining gap is a
  proof that no graph with a $2$-connected BCC of $\betti \geq 2$ has
  a cancellation mechanism producing $\vort = 0$ for \emph{all} sign
  assignments. This would require showing that the Non-Alignment
  Obstruction (Theorem~\ref{thm:non-alignment}) is not only necessary
  but sufficient for $\vort \neq 0$.

\item \textbf{Observer consistency: causal vs.\ temporal gap.}
  Theorem~\ref{thm:topology-consistency} fully characterizes automatic
  temporal consistency.  An open question is the intermediate regime:
  for a BCC $B$ with $\betti(B) \geq 2$, which pairs $(e^*_A, e^*_B)$
  are still temporally consistent (both happen to give $\vort = 0$ by
  orbit symmetry)?  The Orbit Invariance Theorem
  (Theorem~\ref{thm:orbit-consistency}) shows consistency is constant
  within edge orbits, but the fraction of orbit classes with $\vort = 0$
  in graphs with $\betti \geq 2$ declines rapidly and its exact
  characterization remains open.  A second open question is whether
  causal consistency (Definition~\ref{def:causal-consistency}) can fail,
  or whether it holds universally for all observer pairs in all graphs.

\item \textbf{$P_{\mathrm{sel}}$ nonlinearity.} The global discriminant
  \eqref{eq:psel-discriminant} achieves $R^2 = 0.247$, below the
  BCC-local model $R^2 = 0.351$. The missing explanatory power may
  come from nonlinear interactions between BCC structure and coupling
  strength, or from the discrete nature of the Lorentzian selection
  criterion. A nonlinear model or a coupling-dependent predictor might
  close the gap.

\item \textbf{Phase transition curvature exponent (open BC vs.\ PBC).}
  Exact computation on $L \times L$ grids at $J = J_c$ with \emph{open}
  boundary conditions yields
  $|\Scal(J_c)| \sim L^{\nu_{\Scal}}$ with $\nu_{\Scal} \approx 5.6$
  (equivalently $|\Scal| \sim n^{2.8}$ where $n = L^2$ is the number of
  vertices, or $|\Scal| \sim m^{2.15}$ in terms of edges). Finite-size
  scaling theory predicts $\xi \sim L$ at criticality, so the Ruppeiner
  conjecture $|\Scal| \sim \xi^d$ gives $|\Scal| \sim L^2$ for $d = 2$.
  The observed $\nu_{\Scal} \approx 5.6 \gg 2$ suggests that the Fisher
  manifold curvature contains additional information beyond the correlation
  length, possibly encoding boundary-enhanced correlations.

  A sharply contrasting result holds for \emph{periodic} boundary conditions
  (tori). Exact transfer-matrix data for 2D Ising tori $L = 3, \ldots, 8$
  show $|\Scal(J_c)| \sim m^{d_R^{\mathrm{PBC}}}$ with
  $d_R^{\mathrm{PBC}} \to 10/9 = 1.111\ldots$ (measured $1.1110$ at $L = 7{\to}8$),
  consistent with the conjecture $d_R^{\mathrm{PBC}} = (d\nu + 2\eta)/(d\nu + \eta)$
  (\cref{conj:pbc-scaling}). This PBC exponent is much smaller than
  the open-BC exponent, reflecting the elimination of boundary effects on the torus.

  Three open directions: (i)~extend the open-BC computation to larger grids
  ($5 \times 5$ and beyond) and non-square lattices (triangular, hexagonal)
  to confirm the exponent; (ii)~derive $\nu_{\Scal}$ analytically from
  2D Ising critical exponents; (iii)~confirm PBC predictions for
  $q = 3, 4$ Potts models as exact TM computations complete
  (\cref{subsec:pbc-scaling}).
  The $R(J)$ profile of the $4 \times 4$ grid (open BC) shows curvature sign
  change at $J_{\Scal} \approx 0.32 < J_c$, with $\Scal$ growing to $-1512$ at
  $J = 1.0$: the post-critical curvature grows as approximately
  $|J - J_c|^{-\gamma}$.

\item \textbf{3D lattice curvature and dimensional universality.}
  \emph{(Preliminary data obtained; see \cref{rem:3d-lattice}.)}
  The curvature is positive at $J_c^{3D}$ on small 3D cubic lattices
  ($2{\times}2{\times}2$, $2{\times}2{\times}3$), with $J_{\Scal} \approx
  0.27 > J_c^{3D} = 0.222$. This mirrors the 2D triangular lattice
  behavior and arises from the same mechanism (high coordination number
  pushing $J_c$ below $J_{\Scal}$). Open questions remain:
  does the curvature exponent $\nu_{\Scal}$ depend on the spatial
  dimension $d$? Does the Ruppeiner relation $|\Scal| \sim \xi^d$
  hold with the correct $d = 3$ on finite 3D lattices? Larger 3D
  lattices ($3{\times}3{\times}3$ and beyond) are needed to determine
  the 3D scaling exponent.

\item \textbf{Potts model curvature sign transition.}
  \emph{(See \cref{rem:potts-sign}; data for $q = 2$ through $8$.)}
  The curvature at $J_c^{(q)}$ transitions from negative
  ($q = 2$: $\Scal = -3.06$; $q = 3$: $\Scal = -0.08$) to positive
  ($q = 4$: $\Scal = +0.85$) and remains positive through $q = 8$
  ($\Scal \approx +0.71$). Is there a critical $q^* \in (3, 4)$ at
  which $\Scal(J_c) = 0$ exactly? How does $q^*$ depend on the
  lattice? For the $q = 4$ Potts model on the infinite square lattice
  ($J_c = \ln 3$), the phase transition is first-order: does the
  curvature divergence have a different
  character (discontinuity rather than power-law)?
\end{enumerate}

\subsection{Connections to the Literature}

\textbf{Emergent spacetime.} The program of deriving spacetime geometry
from more fundamental degrees of freedom has a rich history: Sakharov's
induced gravity, Jacobson's thermodynamic derivation, Verlinde's entropic
gravity, and causal set theory. Our contribution is a concrete, computable
realization in which the ``more fundamental'' objects are statistical
observables (edge-edge covariances of an Ising model) and the emergent
geometry is the signed Fisher metric.

\textbf{Tensor networks and holography.} The MERA (multi-scale entanglement
renormalization ansatz) and the Pastawski--Yoshida--Harlow--Preskill
HaPPY code~\cite{Pastawski2015} provide toy models of holography using
tensor networks on trees (or trees plus limited loops). Our result gives
a precise curvature interpretation: the tree structure of these networks
corresponds to flat bulk geometry, and additional loops introduce curvature.

\textbf{Information geometry and gravity.} Caticha~\cite{Caticha2019}
has developed an entropic dynamics approach to gravity in which the metric
is constructed from information-theoretic data. Our construction provides
a specific implementation: the Ising Fisher metric as the vielbein, with
spectral gap selection determining the signature.

\textbf{Note on negative results.} We honestly report the negative results:
(1) The Stab$=$Aut criterion for vorticity is falsified by 139 counterexamples.
(2) The modified viability $V = \Delta W - \lambda\|\vort\|^2$ does not
improve self-selection. (3) The $J \to \infty$ problem in Sector A is
unresolved. (4) Self-selection is not universal: the Diamond graph with
$\vort > 0$ assignment also self-selects, showing $\vort = 0$ is not necessary.
(5) The global discriminant \eqref{eq:psel-discriminant} achieves only
$R^2 = 0.247$: global invariants are insufficient to fully predict $P_{\mathrm{sel}}$,
indicating that local BCC structure carries more information.
(6) The $K_n$ Lorentzian selection (Conjecture~\ref{thm:kn-lorentzian})
is verified by enumeration ($K_4$--$K_9$) but not proved analytically:
the cubic secular equation has been computed but the analytical inequality
proof from it is not yet complete.
These negative results are as important as the positive ones for guiding
future work.


%======================================================================
\appendix
%======================================================================

\section{Exact Formulas for $C_4$ and $C_5$}
\label{app:exact}

\subsection{$C_4$ Partition Function and Fisher Matrix}

For $C_4 = (\{0,1,2,3\}, \{(01),(12),(23),(30)\})$ with uniform coupling
$J$ and zero field, label the four edges $e_1, e_2, e_3, e_4$ in order.
The partition function is computed by summing over all $2^4 = 16$ spin
configurations:
\begin{equation}
  \partf_{C_4}(J) = \sum_{\sigma \in \{-1,+1\}^4}
    e^{J(\sigma_0\sigma_1 + \sigma_1\sigma_2 + \sigma_2\sigma_3 + \sigma_3\sigma_0)}.
  \label{eq:z-c4-def}
\end{equation}
Using the identity $e^{J\sigma_i\sigma_j} = \cosh(J) + \sigma_i\sigma_j\sinh(J)$
and summing, one obtains:
\begin{equation}
  \partf_{C_4}(J) = (2\cosh J)^4 + (2\sinh J)^4
  = 16\bigl(\cosh^4 J + \sinh^4 J\bigr)
  = 4\bigl(\cosh(4J) + 3\bigr),
  \label{eq:z-c4}
\end{equation}
where the last equality follows from $\cosh^4 x + \sinh^4 x = (\cosh(4x)+3)/4$.

The log-partition function is $A_{C_4}(J) = \log(4(\cosh(4J)+3))$.

\medskip
The Fisher matrix entries are covariances of the edge observables
$T_a(\sigma)=\sigma_{i(a)}\sigma_{j(a)}$:
$F_{ab}=\Cov(T_a,T_b)$. At uniform coupling, by $\mathbb{Z}_4$ symmetry,
$\Fish_{C_4}$ is equi-off-diagonal:
$F_{aa}=f_d$ and $F_{ab}=f_o$ for $a\neq b$.

Let $t=\tanh(J)$. Direct computation (e.g., via transfer matrix on the ring,
or by exact enumeration of the $2^4$ states) gives:
\begin{align}
  f_d &= \frac{(1-t^2)(1-t^6)}{(1+t^4)^2}, \label{eq:fd-exact}\\
  f_o &= \frac{t^2(1-t^2)^2}{(1+t^4)^2}. \label{eq:fo-exact}
\end{align}

\subsection{$C_4$ Eigenvalues and Square Root}

By the equi-off-diagonal structure (\cref{lem:equi-off-diag}), the
eigenvalues of $\Fish_{C_4}$ are:
\begin{align}
  \lambda_+ &= f_d + 3 f_o \quad (\text{multiplicity } 1), \label{eq:c4-lam-plus}\\
  \lambda_- &= f_d - f_o \quad (\text{multiplicity } 3). \label{eq:c4-lam-minus}
\end{align}
The eigenvector for $\lambda_+$ is $\mathbf{v}_+ = \tfrac{1}{2}(1,1,1,1)^T$
and the three eigenvectors for $\lambda_-$ are orthogonal to $\mathbf{v}_+$,
e.g., $\tfrac{1}{\sqrt{2}}(1,-1,0,0)^T$, etc.

The symmetric square root $\Fish^{1/2}_{C_4}$ has eigenvalues
$\sqrt{\lambda_+}$ and $\sqrt{\lambda_-}$, with the same eigenvectors.
Explicitly:
\begin{equation}
  \Fish^{1/2} = \sqrt{f_d - f_o}\, I
    + \frac{\sqrt{f_d + 3 f_o} - \sqrt{f_d - f_o}}{4}
      \mathbf{1}\mathbf{1}^T,
  \label{eq:c4-sqrt}
\end{equation}
where $\mathbf{1} = (1,1,1,1)^T$.

\subsection{$C_4$ Scalar Curvature: 2D Reduction vs.\ 4D Sign Change}

Even for $C_4$, the scalar curvature depends on whether we restrict to a
symmetry-reduced coupling submanifold or work in the full $m=4$ coupling space.

\textbf{2D reduced model.}
Restrict to the $\mathbb{Z}_2$-symmetric submanifold
$(J_1,J_2,J_3,J_4)=(J_a,J_b,J_a,J_b)$. In this two-parameter model, the
scalar curvature at the uniform point $J_a=J_b=J$ has the closed form:
\begin{equation}
  \Scal_{C_4}^{(2)}(J) = \frac{4\sinh^2(4J)}{(3\cosh(4J)+1)^2}.
  \label{eq:c4-scalar-app}
\end{equation}
In particular, $\Scal_{C_4}^{(2)}(J)>0$ for all $J>0$ and
$\Scal_{C_4}^{(2)}(J)\to 4/9$ as $J\to\infty$.

\textbf{Full 4D coupling space.}
In the full $m=4$ parameter space, the scalar curvature of the Fisher metric
changes sign as $J$ increases (a spherical-to-hyperbolic transition), and the
signed metric scalar curvature is observed numerically to satisfy
$\Scal_{\text{signed}}=\Scal_{\text{Fisher}}/2$ at uniform coupling.

\subsection{$C_5$ and $C_6$ Numerical Eigenvalues}

\begin{table}[ht]
\centering
\caption{Fisher matrix eigenvalues $\lambda_\pm$ for $C_5$ and $C_6$
at selected coupling values. By the equi-off-diagonal structure,
$\lambda_+ = f_d + (n-1)f_o$ (multiplicity 1) and
$\lambda_- = f_d - f_o$ (multiplicity $n-1$). Values from
exact partition function enumeration.}
\label{tab:c5-c6-eigs}
\renewcommand{\arraystretch}{1.3}
\begin{tabular}{llcc}
\toprule
Graph & $J$ & $\lambda_+$ & $\lambda_-$ \\
\midrule
$C_5$ & 0.3 & $8.98 \times 10^{-1}$ & $8.87 \times 10^{-1}$ \\
$C_5$ & 0.5 & $8.13 \times 10^{-1}$ & $7.88 \times 10^{-1}$ \\
$C_5$ & 1.0 & $4.23 \times 10^{-1}$ & $3.83 \times 10^{-1}$ \\
$C_6$ & 0.3 & $8.95 \times 10^{-1}$ & $8.89 \times 10^{-1}$ \\
$C_6$ & 0.5 & $8.05 \times 10^{-1}$ & $7.96 \times 10^{-1}$ \\
$C_6$ & 1.0 & $4.09 \times 10^{-1}$ & $4.02 \times 10^{-1}$ \\
\bottomrule
\end{tabular}
\end{table}

The ratio $\lambda_+/\lambda_- \to 1$ as $n \to \infty$ (girth grows),
reflecting the tree limit: for large girth, $f_o \to 0$ and
$\lambda_+ \approx \lambda_- \approx \sech^2(J)$.

\textbf{Thermodynamic limit ($n \to \infty$).}
Direct numerical computation of $\Scal(n, J)$ for $n = 3$ to $200$
confirms that $\Scal \to 0$ as $n \to \infty$ for all $J > 0$.
The decay is exponential in $n$ at fixed $J$: at $J = 0.5$,
$|\Scal| < 10^{-10}$ already at $n = 30$; at $J = 1.0$,
$|\Scal| < 10^{-5}$ at $n = 75$; at $J = 2.0$, $|\Scal|$ decays
as $n^{-1.56}$, reaching $|\Scal| \approx 26$ at $n = 200$.
The underlying mechanism is $f_o \to 0$ exponentially (as $\tanh^n(J)$),
which reduces the Fisher metric to $f_d \cdot I_m$ (diagonal, flat).
Physically, the 1D Ising model has no phase transition, and accordingly
the Fisher manifold curvature is a \emph{finite-size effect} that vanishes
in the thermodynamic limit. This raises the critical question: does $\Scal$
\emph{survive} in the thermodynamic limit for 2D grids, where a genuine
phase transition exists at $J_c \approx 0.4407$?

%======================================================================
\subsection{General $C_n$ Exact Christoffel Symbols and Einstein Tensor}
\label{app:cn-exact}
%======================================================================

We record here the exact analytical formulas for the Christoffel symbols
and Einstein tensor of the (unsigned Riemannian) Fisher metric on $C_n$,
derived symbolically in SageMath using the transfer matrix and equi-off-diagonal
structure. All formulas use $w = e^J$ and are verified numerically to $10^{-6}$
for $n = 3, 4, 5, 6$ and $J \in [0.3, 2.0]$.

\textbf{Transfer matrix setup.} The partition function of $C_n$ with uniform
coupling $J$ is
\begin{equation}
  Z_n = \lambda_+^n + \lambda_-^n, \quad
  \lambda_\pm = w \pm w^{-1} = 2\cosh(J) \text{ or } 2\sinh(J),
\end{equation}
and the edge mean observable is $m = \frac{1}{n}\frac{d\log Z_n}{dJ}$.
The Fisher metric components are $f_d = 1 - m^2$ and
\begin{equation}
  f_o = \frac{\Var[S] - n f_d}{n(n-1)}, \quad
  S = \sum_{e} X_e \text{ (total edge spin)}.
\end{equation}

\textbf{Proved theorems for $C_n$.}
\begin{enumerate}
  \item $\Gamma^m_{mm} = -m$ (the Christoffel diagonal equals the negative edge mean).
  \item $\Gamma^m_{ab} = 0$ whenever $m \neq a = b$ (off-diagonal diagonal vanishes).
  \item For every Riemann tensor component: $\text{Riem}_\text{linear} = -2\,\text{Riem}_\text{quad}$,
        hence $\text{Riem}_\text{total} = -\text{Riem}_\text{quad}$.
        \emph{Note:} This identity holds not only for $C_n$ but for every graph
        tested (Theorem~\ref{thm:riem-decomp}).
\end{enumerate}

These give exact Riemann quadratic formulas (using the three Christoffel types
$G_{00} = -m$, $G_{01}$, $G_{12}$):
\begin{align}
  \text{Riem}[1,0,1,0]_\text{quad} &= G_{01} G_{00} - G_{01}^2 - (n-2)G_{12}^2, \\
  \text{Riem}[0,0,0,1]_\text{quad} &= G_{01}^2 + (n-2)G_{01}G_{12}, \\
  \text{Riem}[2,0,2,1]_\text{quad} &= G_{01}^2 + G_{00}G_{12} + (n-4)G_{01}G_{12} - (n-3)G_{12}^2.
\end{align}

\textbf{Ricci tensor (exact for all $C_n$):}
\begin{align}
  \text{Ric}_d &= -(n-1)[G_{01}G_{00} - G_{01}^2 - (n-2)G_{12}^2], \\
  \text{Ric}_o &= -(n-1)G_{01}^2 - (n-2)G_{12}[G_{00} + (n-3)G_{01}] + (n-2)(n-3)G_{12}^2.
\end{align}

\textbf{Exact formulas for $C_3$} ($G_{12} = 1/2 = \mathrm{const}$, $R = 3/2 = \mathrm{const}$):
\begin{align}
  G_{00} &= -\frac{(w^2+1)(w+1)(w-1)}{w^4+3}, \\
  G_{01} &= \frac{G_{00}}{2} = -\frac{(w^2+1)(w+1)(w-1)}{2(w^4+3)}, \quad G_{12} = \frac{1}{2}, \\
  R &= \frac{3}{2} \quad (\text{constant, Einstein manifold}), \\
  \alpha(J) &= -\frac{2(w^4+1)}{(w^4+3)^2}, \quad
  \beta(J) = -\frac{(w^2+1)(w+1)(w-1)}{(w^4+3)^2}.
\end{align}

\textbf{Exact formulas for $C_4$:}
\begin{align}
  G_{00} &= -\frac{(w^4+1)(w^2+1)(w+1)(w-1)}{w^8+6w^4+1}, \\
  G_{01} &= -\frac{(w^4+1)(w^2+1)^3(w+1)^3(w-1)^3}{(3w^8+2w^4+3)(w^8+6w^4+1)}, \\
  G_{12} &= \frac{(w^4+1)(w^2+1)(w+1)(w-1)}{2(3w^8+2w^4+3)}, \\
  R(J) &= -\frac{3(3w^8-14w^4+3)(w^4+1)^2(w^2+1)^2(w+1)^2(w-1)^2}{4(3w^8+2w^4+3)^2 w^4}, \\
  \alpha(J) &= \frac{3(3w^{16}-38w^8+3)(w^4+1)^2(w^2+1)^2(w+1)^2(w-1)^2}{(3w^8+2w^4+3)^2(w^8+6w^4+1)^2}, \\
  \beta(J) &= \frac{6(w^{16}+2w^{12}+10w^8+2w^4+1)(w^4+1)^2(w^2+1)^2(w+1)^2(w-1)^2}{(3w^8+2w^4+3)^2(w^8+6w^4+1)^2}.
\end{align}

Note: $R(J)$ changes sign at $3w^8-14w^4+3=0$, i.e., $w^4=(7\pm\sqrt{40})/3$,
giving a critical coupling $J_c = \frac{1}{4}\log\frac{7+\sqrt{40}}{3} \approx 0.371$.

\textbf{Exact formulas for $C_5$:} (abbreviated; full SageMath output in \texttt{scripts/legacy/sage\_cn\_final\_alpha\_beta.py})
\begin{align}
  G_{00} &= -\frac{(w^4+3)(w^2+1)(w+1)(w-1)}{w^8+10w^4+5}, \\
  R(J) &= -\frac{5(3w^{16}+3w^{12}-27w^8-55w^4-20)(w^2+1)^4(w+1)^4(w-1)^4}{2(w^8+2w^4+5)^2(3w^4+1)^3}.
\end{align}

\textbf{Exact formulas for $C_7$:} (scalar curvature)
\begin{align}
  G_{00} &= -\frac{w^{12} + 9w^8 - 5w^4 - 5}{w^{12} + 21w^8 + 35w^4 + 7}, \\
  R(J) &= -\frac{21}{2}\frac{(30w^{32}+325w^{28}+\cdots+525)(w^2+1)^7(w+1)^7(w-1)^7}{(3w^{16}+20w^{12}+114w^8+84w^4+35)^2(5w^8+10w^4+1)^3}.
\end{align}

\textbf{Numerical values of $R$, $\alpha$, $\beta$ for $n = 3, \ldots, 9$:}

\begin{table}[ht]
\centering
\caption{Exact Einstein tensor components $\alpha(J)$ and $\beta(J)$ for
cycle graphs $C_n$, $n=3,\ldots,9$. All values from closed-form symbolic
formulas (SageMath). For $n=7,8,9$: new results (Session 33, 2026-02-18).
The strong-coupling limit $\alpha(n,\infty)$ is given in the last row for each $n$.}
\label{tab:cn-exact-alpha-beta}
\renewcommand{\arraystretch}{1.2}
\begin{tabular}{llrrr}
\toprule
$n$ & $J$ & $R(J)$ & $\alpha(J)$ & $\beta(J)$ \\
\midrule
3 & 0.3 & $+3/2$ & $-0.21631$ & $-0.05808$ \\
3 & 0.5 & $+3/2$ & $-0.15545$ & $-0.05920$ \\
3 & 1.0 & $+3/2$ & $-0.03352$ & $-0.01616$ \\
3 & 2.0 & $+3/2$ & $-0.00067$ & $-0.00034$ \\
& $\infty$ & $+3/2$ & $0$ & $0$ \\
\midrule
4 & 0.3 & $+0.1295$ & $-0.00832$ & $+0.10123$ \\
4 & 0.5 & $-0.5603$ & $+0.17986$ & $+0.22771$ \\
4 & 1.0 & $-12.172$ & $+0.78742$ & $+0.54827$ \\
4 & 2.0 & $-743.7$ & $+0.99554$ & $+0.66414$ \\
& $\infty$ & $-\infty$ & $1$ & $2/3$ \\
\midrule
5 & 0.3 & $+0.0028$ & $+0.00256$ & $+0.03911$ \\
5 & 0.5 & $-0.4998$ & $+0.14536$ & $+0.13761$ \\
5 & 1.0 & $-13.023$ & $+1.08933$ & $+0.51821$ \\
5 & 2.0 & $-825.8$ & $+1.54502$ & $+0.71744$ \\
& $\infty$ & $-\infty$ & $14/9$ & $\cdots$ \\
\midrule
6 & 0.3 & $-0.0101$ & $+0.00352$ & $+0.01344$ \\
6 & 0.5 & $-0.3261$ & $+0.09484$ & $+0.07333$ \\
6 & 1.0 & $-12.375$ & $+1.25613$ & $+0.43216$ \\
6 & 2.0 & $-835.2$ & $+2.04266$ & $+0.70514$ \\
& $\infty$ & $-\infty$ & $33/16$ & $\cdots$ \\
\midrule
7 & 0.3 & $-0.00621$ & $+0.00208$ & $+0.00459$ \\
7 & 0.5 & $-0.19988$ & $+0.05906$ & $+0.03818$ \\
7 & 1.0 & $-11.385$ & $+1.33442$ & $+0.34689$ \\
7 & 2.0 & $-830.3$ & $+2.52696$ & $+0.68303$ \\
& $\infty$ & $-\infty$ & $64/25$ & $52/75$ \\
\midrule
8 & 0.3 & $-0.00273$ & $+0.00094$ & $+0.00156$ \\
8 & 0.5 & $-0.11954$ & $+0.03611$ & $+0.01980$ \\
8 & 1.0 & $-10.303$ & $+1.34524$ & $+0.27316$ \\
8 & 2.0 & $-822.4$ & $+3.00475$ & $+0.66097$ \\
& $\infty$ & $-\infty$ & $55/18$ & $85/126$ \\
\midrule
9 & 0.3 & $-0.00107$ & $+0.00038$ & $+0.00053$ \\
9 & 0.5 & $-0.07003$ & $+0.02166$ & $+0.01025$ \\
9 & 1.0 & $-9.2141$ & $+1.30582$ & $+0.21266$ \\
9 & 2.0 & $-814.1$ & $+3.47729$ & $+0.64081$ \\
& $\infty$ & $-\infty$ & $174/49$ & $129/196$ \\
\bottomrule
\end{tabular}
\end{table}

\textbf{Remark (strong-coupling asymptotics).}
The exact limits $\alpha(n,\infty)$ are rational numbers: $1, 14/9, 33/16, 64/25, 55/18, 174/49$
for $n=4,\ldots,9$. These are \emph{not} equal to $(n-2)/2$ (which would give $1, 3/2, 2, 5/2, 3, 7/2$),
a discrepancy that was missed in earlier numerical work restricted to $J\leq 2$.
The Christoffel symbol limits are exact for all $n$:
\begin{equation}
  G_{00} \to -1, \quad G_{01} \to -\frac{1}{n-1}, \quad G_{12} \to \frac{1}{(n-1)(n-2)}.
\end{equation}

\section{Computational Methods}
\label{app:computation}

All computations were performed in Python 3.11 with NumPy/SciPy on an
Apple M3 Max processor. Symbolic computations used SageMath 10.8. The
code implements a pipeline: graph specification $\to$ partition function
$\to$ Fisher matrix $\to$ signed metric $\to$ curvature $\to$ Raychaudhuri.

\subsection{Partition Function and Fisher Matrix}

\textbf{Small systems ($n \leq 20$):} The partition function
$\partf(\coupl) = \sum_{\boldsymbol{\sigma} \in \{-1,+1\}^n}
\exp(\sum_{e} J_e \sigma_i \sigma_j)$ is computed by brute-force
enumeration (exact, $2^n$ terms). The log-partition function
$A(\coupl) = \log \partf$ is computed with 64-bit double precision.

\textbf{Large systems ($n > 20$):} Wolff cluster Monte Carlo is used
with $10^6$ warmup sweeps and $10^7$ measurement sweeps. The Fisher
matrix is estimated as the sample covariance of edge observables:
$\Fish_{ab} \approx \langle \edgeo_a \edgeo_b \rangle_{\mathrm{MC}}
- \langle \edgeo_a \rangle_{\mathrm{MC}} \langle \edgeo_b \rangle_{\mathrm{MC}}$.

\textbf{Fisher matrix via finite differences:}
\begin{equation}
  \Fish_{ab} \approx \frac{A(\coupl + h\mathbf{e}_a + h\mathbf{e}_b)
    - A(\coupl + h\mathbf{e}_a - h\mathbf{e}_b)
    - A(\coupl - h\mathbf{e}_a + h\mathbf{e}_b)
    + A(\coupl - h\mathbf{e}_a - h\mathbf{e}_b)}{4h^2},
  \label{eq:fisher-num}
\end{equation}
with $h = 10^{-4}$. Verified against exact covariance on $C_4$:
relative error $< 10^{-6}$.

\textbf{Symmetric square root:} $\Fish^{1/2}$ is computed via
eigendecomposition: $\Fish = V \Lambda V^T$ (with $V$ orthogonal,
$\Lambda = \diag(\lambda_1, \ldots, \lambda_m)$), then
$\Fish^{1/2} = V \Lambda^{1/2} V^T$. Verified: $(\Fish^{1/2})^2 = \Fish$
to precision $< 10^{-10}$ for all tested matrices.

\subsection{Christoffel Symbols}

The signed metric $\Met = \Fish^{1/2}\sgn\Fish^{1/2}$ is computed at
each coupling point. The Christoffel symbols are computed by centered
finite differences:
\begin{equation}
  \partial_b \Met_{cd} \approx \frac{\Met_{cd}(\coupl + h\mathbf{e}_b)
    - \Met_{cd}(\coupl - h\mathbf{e}_b)}{2h},
  \label{eq:metric-diff}
\end{equation}
with $h = 10^{-3}$, then
\begin{equation}
  \Chr{a}{bc} = \frac{1}{2}(\Met^{-1})^{ad}
    \bigl(\partial_b \Met_{cd} + \partial_c \Met_{bd}
      - \partial_d \Met_{bc}\bigr).
\end{equation}
The metric inverse $(\Met^{-1})$ is computed by \texttt{numpy.linalg.inv};
nonsingularity is checked ($|\det\Met| > 10^{-10}$). On singular points
(where the metric degenerates), the coupling is perturbed by $10^{-6}$.

\textbf{Consistency check:} On trees, the analytic result
$\Chr{a}{aa} = -\tanh(J_a)$ (all other zero) is verified numerically
to residual $< 10^{-6}$ for $P_3$, $P_4$, $P_6$, $K_{1,4}$, $K_{1,5}$
at coupling $J \in \{0.1, 0.3, 0.5, 1.0\}$ (20 configurations, 100\%
agreement).

\subsection{Riemann Tensor and Scalar Curvature}

The Riemann tensor components are:
\begin{equation}
  \Ric^a{}_{bcd} = \partial_c \Chr{a}{bd} - \partial_d \Chr{a}{bc}
    + \Chr{a}{ce}\Chr{e}{bd} - \Chr{a}{de}\Chr{e}{bc},
\end{equation}
computed numerically by finite-differencing the Christoffel symbols
(step $h = 10^{-3}$). The Ricci tensor $\Ric_{bd} = \Ric^a{}_{bad}$
and scalar curvature $\Scal = (\Met^{-1})^{bd}\Ric_{bd}$ follow by
contraction.

\textbf{Precision:} On $C_4$ at $J = 0.5$, the numerical scalar
curvature is $\Scal_{\mathrm{num}} = 0.06914$ vs.\ the exact formula
\eqref{eq:c4-scalar} giving $\Scal_{\mathrm{exact}} = 0.06913$
(relative error $< 0.01\%$).

\textbf{Tree flatness:} On $P_4$, $P_6$, $K_{1,4}$, star graphs with
$n \leq 10$: $\Scal_{\mathrm{num}} < 10^{-9}$ (within numerical noise),
confirming $\Scal = 0$ for all tested trees. The nonzero residual on
trees is due to finite-difference truncation error, not physical curvature.

\subsection{Raychaudhuri Decomposition}

The observer vector $\obs^a$ is computed as the eigenvector of $\Met$
with the most negative eigenvalue (\texttt{numpy.linalg.eigh}, returns
sorted eigenvalues). The observer is normalized: $\Met_{ab}\obs^a\obs^b
= -1$ (verified; adjustment if needed by rescaling).

The spatial projector $h_{ab} = \Met_{ab} + \obs_a \obs_b$ satisfies
$h_{ab}\obs^b = 0$ (projection property, verified numerically).

The kinematic tensor $B_{ab} = h_a{}^c h_b{}^d \nabla_c \obs_d$ is
computed by:
\begin{enumerate}
\item Compute $\nabla_c \obs_d = \partial_c \obs_d - \Chr{e}{cd}\obs_e$
  numerically.
\item Project: $B_{ab} = h_a{}^c h_b{}^d (\nabla_c \obs_d)$.
\item Decompose: $\expan = h^{ab} B_{ab}$,
  $\shear_{ab} = B_{(ab)} - (\expan/(m-1)) h_{ab}$,
  $\vort_{ab} = B_{[ab]}$.
\end{enumerate}

\textbf{Raychaudhuri identity verification (code consistency check):}
The Raychaudhuri equation is a mathematical identity that holds by
construction for any pseudo-Riemannian manifold; its numerical
verification therefore serves as a code consistency check rather than
physical evidence.  We verify \eqref{eq:raychaudhuri} by computing
both sides independently and checking agreement. Mean residual
$< 10^{-5}$ on all 285 tested configurations ($P_3, P_4, P_6, C_4,
C_5, C_6, C_8$ and random trees with $n \leq 12$, at
$J \in \{0.1, 0.3, 0.5, 1.0\}$, with both $q = 0$ and $q = 1$
sign assignments).

\subsection{$O(N)$ Computations}

For $O(N)$ models, the partition function is computed via Monte Carlo
(Wolff algorithm adapted for continuous spins). The Fisher matrix is
estimated as the sample covariance of edge inner products
$\edgeo_e = \mathbf{s}_i \cdot \mathbf{s}_j$.

For the $N^*$ census (500 graphs, $n = 4$--$12$): for each graph,
$N$ is swept from 1 to 6, and $N^*$ is the minimum $N$ achieving
$> 99\%$ Lorentzian fraction across 10 random coupling vectors in
$[0.1, 2.0]^m$. Each coupling/sign-assignment pair uses $10^5$ Monte
Carlo sweeps (after $10^3$ burn-in); the resulting Fisher matrix
estimates have typical entry-level standard errors $< 2\%$ as estimated
by block-bootstrap resampling.

\subsection{Atlas Census (Vorticity Classification)}

The atlas census of 839 graphs uses the McKay--Piperno graph database
(all connected graphs with $n \leq 9$ vertices and specific $\betti$
values). For each graph:
\begin{enumerate}
\item All $q = 1$ sign assignments ($m$ choices) are enumerated.
\item For each assignment and coupling $J \in \{0.3, 0.5, 1.0\}$,
  compute $\|\vort\|_F = \sqrt{\sum_{ab} \vort_{ab}^2}$.
\item Classify: $\vort = 0$ iff $\|\vort\|_F < 10^{-8}$.
\end{enumerate}
Total: $839 \times \langle m \rangle \times 3 \approx 8348 \times 3
\approx 25000$ computations. Wall time: $\approx 4$ hours on M3 Max.

\section{Vielbein Verification}
\label{app:vielbein}

\subsection{Gauge Invariance Test}

The gauge group of the vielbein construction is the subgroup of $O(m)$
that preserves $\sgn = \diag(+1,\ldots,+1,-1)$:
\begin{equation}
  \mathcal{G} = \{O \in O(m) : O^T \sgn O = \sgn\}
    = O(m-1) \times O(1) \cong O(m-1,1) \cap O(m).
\end{equation}
For each graph and coupling, we generate 100 random elements of
$\mathcal{G}$ (block diagonal with a random $O \in O(m-1)$ on the
spacelike block and $\pm 1$ on the timelike block) and verify:
\begin{equation}
  \|(O\Fish^{1/2})^T \sgn (O\Fish^{1/2}) - \Met\|_F < 10^{-10}.
  \label{eq:gauge-test}
\end{equation}
Result: 100\% agreement on all tested graphs
($P_4, C_4, C_6, K_4, K_{3,3}$, random graphs with $n \leq 12$).

\subsection{Uniqueness of Symmetric Square Root}

\begin{lemma}[Uniqueness of PD Square Root]
\label{lem:sqrt-unique}
For any symmetric positive-definite matrix $A \in \mathbb{R}^{m \times m}$,
there exists a unique symmetric positive-definite $B$ satisfying $B^2 = A$.
\end{lemma}

\begin{proof}
Existence: $A = V\Lambda V^T$ (spectral decomposition, $\Lambda > 0$);
set $B = V\Lambda^{1/2}V^T$ (symmetric PD). Then $B^2 = V\Lambda V^T = A$.

Uniqueness: Suppose $B_1^2 = B_2^2 = A$ with $B_1, B_2$ symmetric PD.
Then $B_1 B_2 B_1 = B_1 (B_2^2) = B_1 A$ and also $B_1 B_2 B_1
= (B_1 B_2) B_1$. Setting $C = B_1 B_2^{-1}$ (well-defined since
$B_2$ is PD), we get $C B_1 = B_1 C^{-1}$, and taking transposes
(using symmetry) $C^T B_1 = B_1 (C^{-1})^T$. Combined with
$CB_1 = B_1 C^{-1}$, this gives $C = C^{-T}$, so $C$ is orthogonal.
A symmetric PD matrix has a unique positive square root, so $C$ must
be the identity, giving $B_1 = B_2$.
\end{proof}

\noindent Numerical verification: for 50 random symmetric PD matrices $A$
with $m \in \{3, 4, 5, 6\}$, both the eigendecomposition method and
the Newton--Schulz iteration give the same $A^{1/2}$ to precision
$\|B_1 - B_2\|_F < 10^{-12}$.

\subsection{Comparison: $\Met = F^{1/2}SF^{1/2}$ vs.\ $M = FSF$}

The signed mass tensor $M = FSF$ (from \cite{Paper1}) and the signed
Fisher metric $\Met = F^{1/2}SF^{1/2}$ are different matrices. We
verify their distinct properties:

\begin{table}[ht]
\centering
\caption{Comparison of $M = FSF$ (signed mass tensor) and
$\Met = \Fish^{1/2}\sgn\Fish^{1/2}$ (signed Fisher metric) on $C_4$
at $J = 0.5$, $q = 1$. Both are indefinite (Lorentzian) but have
different eigenvalue spectra and curvatures.}
\label{tab:m-vs-metric}
\renewcommand{\arraystretch}{1.3}
\begin{tabular}{lcc}
\toprule
Property & $M = FSF$ & $\Met = F^{1/2}SF^{1/2}$ \\
\midrule
Eigenvalues (negative) & $\lambda_1(FSF)$ & $\lambda_1(F^{1/2}SF^{1/2})$ \\
Scalar curvature & $\Scal_M$ & $\Scal_{\Met} = \Scal_M/2$ \\
Vielbein form & $F(SF) = e^T (F\eta F) e'$ (not standard) & $e^T \eta e$ (standard) \\
Tree flatness & Yes ($FSF = \sech^4(J) S \diag$) & Yes ($F^{1/2}SF^{1/2} = \sech^2(J) S \diag$) \\
Sign independence & Yes & Yes \\
\bottomrule
\end{tabular}
\end{table}

The factor-of-two relationship $\Scal_{FSF} = 2\Scal_{F^{1/2}SF^{1/2}}$
holds numerically on all tested cycles (exact on $C_4$ from
\eqref{eq:c4-scalar} and the analogous $FSF$ formula). The vielbein
interpretation favors $\Met = F^{1/2}SF^{1/2}$ as the canonical
construction.


%======================================================================
\begin{thebibliography}{99}
%======================================================================

\bibitem{Amari1998}
S.-i.\ Amari and H.\ Nagaoka,
\textit{Methods of Information Geometry},
AMS Translations of Mathematical Monographs, vol.\ 191 (2000).
Originally published as \textit{Joho Kika no Hoho}, Iwanami Shoten (1993).

\bibitem{Cencov1982}
N.\ N.\ Cencov,
\textit{Statistical Decision Rules and Optimal Inference},
AMS Translations of Mathematical Monographs, vol.\ 53 (1982).
[Uniqueness of the Fisher metric as the only monotone Riemannian
metric on statistical manifolds.]

\bibitem{Matsueda2013}
H.\ Matsueda,
``Geodesic distance in Fisher information space for massive Dirac fermions,''
\textit{arXiv}:1305.6955 (2013).
[First computation of Christoffel symbols for Ising-type Fisher metrics.]

\bibitem{Machta2013}
B.\ B.\ Machta, R.\ Chachra, M.\ K.\ Transtrum, and J.\ P.\ Sethna,
``Parameter space compression underlies emergent theories and
predictive models,''
\textit{Science} \textbf{342}, 604--607 (2013).

\bibitem{Greensite1993}
J.\ Greensite,
``Dynamical origin of the Lorentzian signature of spacetime,''
\textit{Physics Letters B} \textbf{300}, 34--37 (1993).
[Dynamical signature selection via path integral; closest antecedent
to vielbein construction.]

\bibitem{Lashkari2014}
N.\ Lashkari and M.\ Van Raamsdonk,
``Canonical energy is quantum Fisher information,''
\textit{Journal of High Energy Physics} \textbf{2016}, 153 (2016).
[arXiv:1508.00897. Connects bulk metric perturbations to boundary
Fisher information.]

\bibitem{RyuTakayanagi2006}
S.\ Ryu and T.\ Takayanagi,
``Holographic derivation of entanglement entropy from AdS/CFT,''
\textit{Physical Review Letters} \textbf{96}, 181602 (2006).

\bibitem{Jacobson1995}
T.\ Jacobson,
``Thermodynamics of spacetime: The Einstein equation of state,''
\textit{Physical Review Letters} \textbf{75}, 1260 (1995).

\bibitem{Verlinde2011}
E.\ Verlinde,
``On the origin of gravity and the laws of Newton,''
\textit{Journal of High Energy Physics} \textbf{2011}, 29 (2011).

\bibitem{Pastawski2015}
F.\ Pastawski, B.\ Yoshida, D.\ Harlow, and J.\ Preskill,
``Holographic quantum error-correcting codes: Toy models for the
bulk/boundary correspondence,''
\textit{JHEP} \textbf{2015}, 149 (2015).

\bibitem{Uhlmann1976}
A.\ Uhlmann,
``The transition probability in the state space of a *-algebra,''
\textit{Reports on Mathematical Physics} \textbf{9}, 273--279 (1976).
[Bures metric as $\mathrm{Tr}(\rho^{1/2} A \rho^{1/2})$.]

\bibitem{Paper1}
M.~Zhuravlev,
``Where the Lovelock Bridge Breaks: Negative Results and New Directions
for Connecting Discrete and Continuous Spacetime Emergence,''
Preprint (2026).
[Tree Fisher Identity; PSD obstruction.]

\bibitem{Paper5}
M.~Zhuravlev,
``Spectral Gap Selection of Lorentzian Signature from Observer Fisher
Information Geometry,''
Preprint (2026).
[Theorems A--C; $W(q=1) > W(q\geq 2)$; O(N) hierarchy.]

\bibitem{Caticha2019}
A.\ Caticha,
``The Entropic Dynamics approach to Quantum Mechanics,''
\textit{Entropy} \textbf{21}, 943 (2019).
[Information-theoretic approach to deriving spacetime geometry.]

\bibitem{Rao1945}
C.\ R.\ Rao,
``Information and the accuracy attainable in the estimation of
statistical parameters,''
\textit{Bulletin of the Calcutta Mathematical Society} \textbf{37},
81--91 (1945).
[Original derivation of the Fisher metric.]

\bibitem{Ruppeiner1979}
G.\ Ruppeiner,
``Thermodynamics: A Riemannian geometric model,''
\textit{Physical Review A} \textbf{20}, 1608--1613 (1979).
[Original Ruppeiner metric: curvature of thermodynamic fluctuation geometry.]

\bibitem{Ruppeiner1995}
G.\ Ruppeiner,
``Riemannian geometry in thermodynamic fluctuation theory,''
\textit{Reviews of Modern Physics} \textbf{67}, 605--659 (1995).
[Review of thermodynamic geometry; $|R| \sim \xi^d$ near critical points.]

\bibitem{BrodyRitz2003}
D.\ C.\ Brody and A.\ Ritz,
``Information geometry of finite Ising models,''
\textit{Journal of Geometry and Physics} \textbf{47}, 207--220 (2003).
[Finite-size information-geometric analysis for small Ising systems.]

\bibitem{Janke2004}
W.\ Janke, D.\ A.\ Johnston, and R.\ Kenna,
``Information geometry and phase transitions,''
\textit{Physica A} \textbf{336}, 181--186 (2004).
[Information-geometric curvature for Ising, Potts, and spherical models
on 2D thermodynamic manifold; confirms Ruppeiner scaling.]

\bibitem{Erdmenger2020}
J.\ Erdmenger, K.\ Franze, C.\ M. Melby-Thompson, and R.\ Meyer,
``Information geometry of the Ising model on planar random graphs,''
\textit{Journal of High Energy Physics} \textbf{05}, 104 (2020).
[Coupling-space information geometry in low-dimensional parameter families.]

\bibitem{Vanchurin2022}
V.\ Vanchurin,
``Towards a theory of quantum gravity from neural networks,''
\textit{Entropy} \textbf{24}, 7 (2022).
[Type II-style emergent-gravity motivation for observer parameter geometry.]

\bibitem{PredLetter}
M.~Zhuravlev,
``Fisher Curvature Scaling at Critical Points:
An Exact Information-Geometric Exponent,''
Preprint (2026).
[Conjecture $d_R^{\mathrm{PBC}} = (d\nu+2\eta)/(d\nu+\eta)$;
2D Ising exact TM verification; 3D Ising MCMC; Potts $q=3,4$.]

\end{thebibliography}

\end{document}
